\section{Theorem S III: the spectral argument and assembly}\label{sec:pf:thmS}

\providecommand{\thmSClosed}{\sfname{Closed}}
\providecommand{\thmSEv}{\sfname{Ev}}
\providecommand{\thmSdC}{d_{C}}
\providecommand{\thmSepsL}{\varepsilon_{L}}
\providecommand{\thmSId}{\mathrm{Id}}
\providecommand{\thmSDeg}{\mathrm{Deg}}
\providecommand{\thmSNG}{\mathrm{NG}}
\providecommand{\thmSGen}{\mathrm{Gen}}
\providecommand{\thmSNZ}{\mathrm{NZ}}
\providecommand{\thmSpart}{\operatorname{part}}
\providecommand{\thmSvp}{\operatorname{vp}}
\providecommand{\thmSps}{\operatorname{ps}}
\providecommand{\thmSpa}{\operatorname{pa}}
\providecommand{\thmSrow}{\operatorname{row}}
\providecommand{\thmSlev}{\operatorname{lev}}
\providecommand{\thmSjoin}{\operatorname{join}}
\providecommand{\thmSblk}{\operatorname{blk}}
\providecommand{\thmSmaj}{\operatorname{maj}}
\providecommand{\thmSSf}{\mathrm{Sf}}
\providecommand{\thmSNst}{\mathrm{N}}
\providecommand{\thmSNz}{\mathrm{N}_0}
\providecommand{\thmSNh}{\widehat{\mathrm{N}}}
\providecommand{\thmSDg}{\mathrm{Dg}}
\providecommand{\thmSgram}{\mathrm{gr}}
\providecommand{\thmStt}{\mathrm{tt}}

This section proves the three ``local rows'' $\mathsf{sLp}$, $\mathsf{sO}$, $\mathsf{hTF}$ of the hypothesis
package from the Siegel--Weil formula, and, on the way, Theorem~S (\cref{thm:theorem-S}). The two largest inputs of
Theorem~S, the nineteen hypotheses of its pointwise form and the lower bound for Poincar\'e amplitudes, are proved
elsewhere in Part~II (\lref{GlobalTheory.CarrierGap.theoremS_leaves_of_carrierG} and
\lref{GlobalTheory.CarrierGap.carrierP}); here we give the assembly and the spectral argument which turns these
inputs into the bound of Theorem~S.

\paragraph{The idea.}
Let $u\ne u'$ be two words of a datum, and let $C_u$, $C_{u'}$ be the finite-dimensional spaces of automorphic forms
attached to them, of dimensions $d_u$, $d_{u'}$. Theorem~S asserts that the operator relating $C_{u'}$ to $C_u$ has
Hilbert--Schmidt mass at most $(1+b^{-6})\,d_ud_{u'}/D_*$, where $D_*$ is the main term of the relevant space; that is,
$C_u$ and $C_{u'}$ are, up to a factor $1+b^{-6}$, as orthogonal as two independent random subspaces
(\cref{rem:families:random}). The proof follows the relative trace formula strategy of \cite{Announcement2026}.
The operator has an orbital expansion $K=\sum_t\mathrm{op}(t)$ over \emph{terms} $t$: a term is a pair of row
lattices $\Lambda,\Lambda'$ with a class of rational isometries $\gamma$ satisfying $\gamma\Lambda'=\Lambda$, and
$\mathrm{op}(t)$ is an integral of rank-one operators over a fundamental domain. Terms are identity, degenerate,
generic or non-generic; identity and degenerate terms contribute nothing, and $K_0$, the leading coefficient, is
the sum of the generic and non-generic parts. The Hilbert--Schmidt norm of each part is expanded by a Gram
identity over the finitely many ``Levi types'' of the terms (a finite Fourier analysis on products of groups
$\GL_k(\F_q)$ at two auxiliary primes); after a positive semi-definite comparison of the Gram matrix with a scalar,
it becomes a double sum over pairs of terms in the same row. Pairs are sorted by a corank $s_0\ge1$: pairs of corank
at least two and non-scalar pairs of corank one (the \emph{loop rows}) and non-generic pairs of corank one are
negligible, and the scalar pairs of corank one give the \emph{main row}, which is bounded by $(1+\delta)e_1\,
d_ud_{u'}/D_*$ through a Bessel-type inequality for a tower of positive semi-definite kernels (``Theorem~A'') and a
chain of identities between volumes and masses. Explicit numerics then give the factor $1+b^{-6}$.

\paragraph{Organisation.}
\Cref{sec:pf:thmS:conv} fixes the notions of families, closed properties and selections used throughout.
\Cref{sec:pf:thmS:rows} proves the three local rows (the root result, \cref{pf:thmS:lg1}) from Theorem~S, the loop
rows and the rank and trace-frame rows. \Cref{sec:pf:thmS:fam} deduces Theorem~S (\cref{pf:thmS:main}) from its
pointwise form, and \cref{sec:pf:thmS:pw} reduces the pointwise form to an abstract statement about term systems.
\Cref{sec:pf:thmS:abs} contains the spectral argument for term systems, \cref{sec:pf:thmS:main} the main row and
Theorem~A, and \cref{sec:pf:thmS:loops} the loop rows. \Cref{sec:pf:thmS:deferred} proves the remaining inputs of these
arguments and several results used in other sections of Part~II (density of $K$-finite vectors, a Capelli identity, the
joint tensor identity of the carrier, lattice facts and size bounds), and \cref{sec:pf:thmS:remarks} records
peculiarities of the formal statements.

% ---------------------------------------------------------------------------
\subsection{Conventions: families, closed properties, selections}\label{sec:pf:thmS:conv}

We use the notions of \cref{def:families:closed,def:families:rows}. A family $F\in\sfname{FamP}(\eps)$ has a threshold
$x_1$, index sets $I_x$, members $(D_{x,i},L_{x,i})$ for $x>x_1$, and lists of constants: a list $c_{\mathrm{out}}$,
a list $c_{\mathrm{mid}}(C_H)$ depending on an auxiliary constant $C_H$, a last list, and named constants such as
$C_{\mathrm{YL}}$, $C_{\mathrm{Tu}}$, $C_\Gamma$. Two derived constants are used constantly:
$s_{C'}(F)$ (the constant $C_{k,\mathrm{low}}$ of $F$) and $s_{C_l}(F)$ (the minimum of six constants of $F$).
For a datum $D$ we write $k,m,f,r,r_S$ for its ranks, $E$ for its CM host field, $T_u$ for the level prime of a word
$u$, $d^S_u\in\N$ for the number attached to $u$ (the ``size'' of $u$), $D_*$ for its main term, and for size data
$L$ we write $b=b(L)\ge2$ for the scale quantity $b=e^{x^\eps}$.

\begin{definition}[Eventually; closed]\label{pf:thmS:closed}
For a family $F$ and a predicate $P(C',C_l,x,i)$, $\thmSEv_F\,P$ means: there is $x_0$ such that
$P(s_{C'}(F),s_{C_l}(F),x,i)$ holds for all $x\ge x_0$ and $i\in I_x$. For a class $\mathcal C$ and a property
$P(F,C_H)$, $\thmSClosed(\mathcal C,P)$ means: there is $\eps_1>0$ such that for every $\eps\in(0,\eps_1]$ there is a
staged threshold $K$ (a constant $K_{\mathrm{pre}}$ and functions $K_{\mathrm{mid}},K_{\mathrm{last}}$ of lists of
reals) with $P(F,C_H)$ for all $F\in\mathcal C\cap\sfname{FamP}(\eps)$ and all $C_H$ such that the constants of
$F$ are \emph{large} for $(K,C_H)$: every entry of the first list is $\ge K_{\mathrm{pre}}$, every entry of
$c_{\mathrm{mid}}(C_H)$ is $\ge K_{\mathrm{mid}}(c_{\mathrm{out}})$, and every entry of the last list is
$\ge K_{\mathrm{last}}(c_{\mathrm{out}}\mathbin{+\!\!+}c_{\mathrm{mid}}(C_H))$.
\end{definition}

The following shared lemma is used to derive closed statements; it is proved in
\lref{Formalism.L2.SG.ThmSOwn.closed_with}: \emph{if $\thmSClosed(\mathcal C,Q)$ holds and $\theta_m,\theta$ are
functions of lists of reals such that $Q(F,C_H)$ implies $P(F,C_H)$ whenever $F\in\mathcal C$,
$\theta_m(c_{\mathrm{out}})\le C_\Gamma(F)$, $\theta(c_{\mathrm{out}}\mathbin{+\!\!+}c_{\mathrm{mid}}(C_H))\le
s_{C'}(F)$ and $\theta(c_{\mathrm{out}}\mathbin{+\!\!+}c_{\mathrm{mid}}(C_H))\le s_{C_l}(F)$, then
$\thmSClosed(\mathcal C,P)$.} In words: in a closed statement, the constants $s_{C'}(F),s_{C_l}(F)$ may be assumed
larger than any prescribed function of the earlier constants. Closed statements are also stable under
conjunction (take the largest of the staged thresholds).

\paragraph{Spectral selections and term systems.}
A \emph{spectral selection} $\sigma$ for a class $\mathcal C$ assigns to every $F\in\mathcal C$ a threshold $x_0(F)$
and, for $x>x_1$, $x\ge x_0(F)$ and $i\in I_x$, \emph{spectral ties} $\sigma(F)_{x,i}$ for $D_{x,i}$. Spectral ties
$T$ for a datum $D$ consist of a Hida setting $X=T.X$ (an automorphic instance of a unitary group over $E$, with a
level subgroup character, discrete series data at the infinite places and a Haar measure on the finite adelic
points), an adelic bridge, and a \emph{full system} $Y=T.Y$ (flags, row representatives, fundamental sets,
column systems and a typed reduction at the two type primes). For words $u,u'$ the full system defines a
\emph{term system} $Y.t_O(u,u')$ in the sense of \cref{pf:thmS:terms} below, whose terms are the triples
$(\Lambda,\Lambda',[\gamma])$ described in the introduction, with operators in the Hilbert--Schmidt space of maps
between the column spaces of $u'$ and $u$; we write $K_0(u,u')$ for its leading coefficient. The pair $(X,Y)$
is \emph{concretely tied} if, among other things, the dimension of the space $C_u$ of $X$ equals $d^S_u$ and the
main term $D_*(X)$ of $X$ equals $D_*$. We put
\[
\thmSdC(u,u')=\frac{\dim C_u\,\dim C_{u'}}{D_*(X)},\qquad \thmSepsL(D)=\exp\Bigl(-\tfrac13\,k f\log k\Bigr).
\]

\paragraph{Standing and count hypotheses.}
Data and size data carry \emph{standing hypotheses} and \emph{standing numerical data}, depending on constants; we use
them only through the consequences cited below. The \emph{admissible bounds} are the standing numerical data with
constant $C_0$, the standing hypotheses with constants $(C',C_l,C_{\mathrm{YL}}+1,\frac14,3,\frac12,C_{\mathrm{Tu}})$, and
$1\le C_0$. The count hypotheses $\mathrm{CH}_0(D,L;C',C_l,C_g,c_b,C_u,c_T,C_T,C_0)$ are the standing hypotheses and data
with these constants and $1\le C_0$; the count hypotheses $\mathrm{CH}(D,L,u,u';\dots,C_0,A)$ for two words add the
\emph{standard conditions} $\mathrm{St}(D,u,u')$ (the level primes $T_u,T_{u'}\ge3$ and the type primes $g_0\ne g_1$ are
primes split completely in $E$, $g_S(D)$ is prime to $T_uT_{u'}$, $T_u\equiv T_{u'}\bmod g_S(D)$, and $r_S\ge1$), the
comparisons $T_u/T_{u'},T_{u'}/T_u\le A$ and $\mathrm{vand}(u_\tau)\le A^k\mathrm{vand}(u'_\tau)$ (and symmetrically) for
the Vandermonde products of the Harish-Chandra parameters at every embedding type $\tau$, and the inequalities
$C_lm^6\Lambda\le k\log k$, $(\log f)^2\le Am^2\Lambda$, $\log f\le Ar_S^2$, where $\Lambda$ is the scale quantity of $L$
entering the size conditions.

\begin{definition}[Loop rows, Theorem S along a selection, trace frame]\label{pf:thmS:defrows}
Let $\mathcal C$ be a class, $\sigma$ a spectral selection for $\mathcal C$ and $C_0\in\R$. Write $X,Y$ for the
components of $\sigma(F)_{x,i}$ and $D=D_{x,i}$, $L=L_{x,i}$.
\begin{enumerate}
\item (Loop rows.) $\mathcal{L}_{C_0}(\mathcal C,\sigma)$: for every function $a_F$ of lists of reals,
$\thmSClosed(\mathcal C,P)$, where $P(F,C_H)$ says that $\thmSEv_F$ of the following holds: for all admissible
$x$ and all words $u,u'$ of $D$, if the count hypotheses $\mathrm{CH}(D,L,u,u';C',C_l,C_{\mathrm{YL}}+1,\tfrac14,3,
\tfrac12,C_{\mathrm{Tu}},C_0,a_F(c_{\mathrm{out}}))$ hold, $(X,Y)$ is concretely tied and the term system
$Y.t_O(u,u')$ is absolutely convergent, then the row bounds $\mathrm{R}_{\mathrm{big}}$ and $\mathrm{R}_{\mathrm{nsc}}$ of
\cref{pf:thmS:terms} hold for $Y.t_O(u,u')$ with $\varepsilon=\thmSepsL(D)$ and $dd=\thmSdC(u,u')$.
\item (Theorem S along the selection.) $\mathcal{S}(F,C_0,\sigma(F))$: $\thmSEv_F$ of the following: if the admissible
bounds hold, then for all admissible $x$ and all words $u\ne u'$ of $D$,
\begin{equation}\label{eq:pf:thmS:S}
\|K_0(u,u')\|^2\le\bigl(1+b^{-6}\bigr)\,\frac{d^S_u\,d^S_{u'}}{D_*}.
\end{equation}
\item (Trace frame.) $\mathrm{TF}(F;C_g,c_b,C_u,c_T,C_T,C_0)$: $\thmSEv_F$ of: if the count hypotheses $\mathrm{CH}_0$ hold for $(D,L)$ with the constants
$(C',C_l,C_g,c_b,C_u,c_T,C_T,C_0)$, then $1\le T_w\,\kappa_S(D)^{2r_S}$ for every word $w$,
where $\kappa_S(D)$ is the quotient of the constant $\kappa$ of the host by a constant $g(D)$ of the datum.
\item (Mass formulas.) $\mathrm{MF}(F;C_0,\dots)$: $\thmSEv_F$ of: under the corresponding standing hypotheses, the \emph{mass
formulas} hold for $D$: for every hermitian $\OO_E$-lattice $L$ with a symmetric form which is orthogonal of rank $k$,
$d^S_u$ equals the mass of $L$ times explicit local and archimedean factors (a product of Vandermonde values of the
Harish-Chandra parameters of $u$, factorials, and $(T_u^k-1)^f$), and for $L$ orthogonal of rank $r$, $D_*$ equals the
mass of $L$ times explicit factors. The \emph{rank row} for $\mathcal C$ is
$\mathcal{R}_{C_0}(\mathcal C)=\thmSClosed(\mathcal C,\mathrm{MF}(\cdot;C_0,C_{\mathrm{YL}}+1,\tfrac14,3,\tfrac12,
C_{\mathrm{Tu}}))$.
\end{enumerate}
\end{definition}

The three rows of the package treated here are, for a slot $Q$ and a linkage $\mathrm{Lk}$ for bundles (in the
sense of \cref{def:families:rows}),
\begin{align*}
\mathsf{sLp}:&\ \ \sfname{Row}\bigl(Q,\mathrm{Lk},(\mathcal A,\sigma)\mapsto
\mathcal{L}_{3}(\mathcal D_3,\sigma^{\mathrm{sp}})\bigr),\\
\mathsf{sO}:&\ \ \sfname{Row}\bigl(Q,\mathrm{Lk},(\mathcal A,\sigma)\mapsto
\thmSClosed\bigl(\mathcal D_3,F\mapsto\mathcal{S}(F,3,\sigma^{\mathrm{sp}}(F))\bigr)\bigr),\\
\mathsf{hTF}:&\ \ \sfname{Row}\bigl(Q,\mathrm{Lk},(\mathcal A,\sigma)\mapsto
\thmSClosed\bigl(\mathcal D_3,\mathrm{TF}(\cdot;C_{\mathrm{YL}}+1,\tfrac14,3,\tfrac12,C_{\mathrm{Tu}},3)\bigr)\bigr),
\end{align*}
where $\mathcal D_3=D3(3,\mathcal A)$ is the class of designed families of \cref{def:families:closed}, and for a raw
selection $\sigma$ of bundles over $\mathcal A$, $\sigma^{\mathrm{sp}}$ is the spectral selection with threshold $x_1$
whose spectral ties at $(F,x,i)$ are those of the spectral record of the bundle $\sigma(F,x,i)$.
For a predicate $\mathrm{KP}$ on bases and frames, slots $Q,Q'$ and a property $r$ of slots, $\mathrm{Lg}(\mathrm{KP},
Q,Q',r)$ means: there is a slot $R$ such that, eventually in the scale and for large twisting data, every core of $Q$
satisfying $\mathrm{KP}$ lies in $R$, and $r(Q'\wedge R)$ holds.

% ---------------------------------------------------------------------------
\subsection{The three local rows}\label{sec:pf:thmS:rows}

\begin{theorem}[Local rows from the Siegel--Weil formula]\label{pf:thmS:lg1}
Assume the Siegel--Weil formula in the convergent range (\cref{lit:families:sw}). Then for every predicate
$\mathrm{KP}$ on bases and frames, all slots $Q,Q'$ and every linkage $\mathrm{Lk}$ for bundles,
$\mathrm{Lg}(\mathrm{KP},Q,Q',S\mapsto\mathsf{sLp}(S,\mathrm{Lk}))$, $\mathrm{Lg}(\mathrm{KP},Q,Q',
S\mapsto\mathsf{sO}(S,\mathrm{Lk}))$ and $\mathrm{Lg}(\mathrm{KP},Q,Q',S\mapsto\mathsf{hTF}(S,\mathrm{Lk}))$ hold.
\formallabel{GlobalTheory.LocalRows.lg1_sLp_sO_hTF}
\end{theorem}

\begin{proposition}[The rows at every slot]\label{pf:thmS:everyslot}
Under the Siegel--Weil formula, $\mathsf{sLp}(Q,\mathrm{Lk})$, $\mathsf{sO}(Q,\mathrm{Lk})$ and $\mathsf{hTF}(Q,\mathrm{Lk})$
hold for every slot $Q$ and every linkage $\mathrm{Lk}$.
\formallabel{GlobalTheory.LocalRows.sLp_sO_hTF_of_siegelWeil}
\end{proposition}

\begin{proposition}[The rows at a selection]\label{pf:thmS:atchoice}
Assume the Siegel--Weil formula. For every class $\mathcal A$ and every raw selection $\sigma$ of bundles over
$\mathcal A$, with $\mathcal D_3=D3(3,\mathcal A)$:
\begin{enumerate}
\item $\mathcal{L}_{3}(\mathcal D_3,\sigma^{\mathrm{sp}})$;
\item $\thmSClosed(\mathcal D_3,F\mapsto\mathcal{S}(F,3,\sigma^{\mathrm{sp}}(F)))$;
\item $\thmSClosed(\mathcal D_3,\mathrm{TF}(\cdot;C_{\mathrm{YL}}+1,\tfrac14,3,\tfrac12,C_{\mathrm{Tu}},3))$.
\end{enumerate}
\formallabel{GlobalTheory.LocalRows.at_choice_alt}
\end{proposition}

We first record the rank and trace-frame rows. A \emph{twist selection} for a class $\mathcal C$ chooses, for every
member of every $F\in\mathcal C$ past a threshold, a tied record (theta data with real structure, column data and
levels); a raw selection $\sigma$ of bundles over $\mathcal A$ induces one on $\mathcal D_3$, by taking the records
of the bundles.

\begin{lemma}[Rank and trace-frame rows]\label{pf:thmS:ranks}
Assume the mass-ratio identity for every datum $D$ (if the volumes of $D$ are tied, they satisfy the volume ratio of \cref{prop:families:volratio}(i)). Then for every class $\mathcal C$ admitting a twist selection,
$\mathcal{R}_{3}(\mathcal C)$ holds, and so does
$\thmSClosed(\mathcal C,\mathrm{TF}(\cdot;C_{\mathrm{YL}}+1,\tfrac14,3,\tfrac12,C_{\mathrm{Tu}},3))$.
\formallabel{GlobalTheory.ranks_and_frame_of_tied_class_alt}
\end{lemma}

\begin{proof}
For the rank row it suffices to prove $\mathrm{MF}(F;C_0,\dots)$ for every family $F$ with a twist selection and
all constants (a property holding for every member of the class is closed). This is
\lref{GlobalTheory.exact_ranks_along_tied_family_alt}: past the threshold of the selection, the tied record
gives the real structure of the theta data and the equality of the column levels with the symmetric level; the
class number formula of the development gives the rank formula for the levels; the type primes split
completely (by the standing hypotheses); the volume of the definite space equals the mass by the mass-ratio
identity; and \lref{GlobalTheory.exact_ranks_of_record} rewrites $d^S_u$ and $D_*$ through the mass of orthogonal
lattices, the product over the places above the type primes, and the exact local indices of columns and rows. The
trace-frame row follows from the rank row by \cref{pf:thmS:tf} below with $C_0=3$, applied with the formal-degree identity and the representation property of CM fields (\cref{pf:thmS:rem1}(2)), which are theorems of the development
(\lref{Formalism.L2.SG.OwnFacts.FD.factDsigHC} and \lref{GlobalTheory.hermRepUnram}); this is
\lref{GlobalTheory.tFrame_of_ranksRow_three}.
\end{proof}

\begin{corollary}\label{pf:thmS:bothrows}
The Siegel--Weil formula implies, for every class with a twist selection, the rank row $\mathcal{R}_{3}$ and the
trace-frame row of \cref{pf:thmS:ranks}.
\formallabel{GlobalTheory.bothRowsFromSiegelWeil_alt}
\end{corollary}

\begin{proof}
The mass-ratio identity for every datum follows from the Siegel--Weil formula by
\lref{GlobalTheory.VolRatioFromSiegelWeil.volumeRatio} (\cref{prop:families:volratio}(i)); apply
\cref{pf:thmS:ranks}.
\end{proof}

\begin{proof}[Proof of \cref{pf:thmS:atchoice}]
Let $T$ be the twist selection on $\mathcal D_3$ induced by $\sigma$. By \cref{pf:thmS:ranks} and the mass-ratio
identity from the Siegel--Weil formula, $\mathcal{R}_{3}(\mathcal D_3)$ and the trace-frame row hold; the latter is
(3). Members of $\mathcal D_3$ satisfy the size condition $\sfname{TuSize}$ (it is part of the definition of designed
families). Hence (1) is the loop-row theorem \cref{pf:thmS:loops}, applied to $\mathcal D_3$, $T$ and
$\sigma^{\mathrm{sp}}$, with finite-dimensionality of the row spaces $\mathrm{FD}$
(\lref{GlobalTheory.rowsFinDim_link}) and the mass-ratio identity. For (2) apply Theorem~S
(\cref{pf:thmS:main}) to $\mathcal D_3$ and $\sigma^{\mathrm{sp}}$: its hypotheses are $\sfname{SWXS}$ and
$\sfname{RowNGP}$, which follow from the Siegel--Weil formula
(\lref{GlobalTheory.St.swxs_of_siegelWeil} and \lref{GlobalTheory.rowNGP_of_siegelWeil},
\cref{prop:families:volratio}(ii),(iii)), $\mathrm{FD}$, the local ties of the selected data (a field of the
spectral record of every bundle), the rank row and (1).
\end{proof}

\begin{proof}[Proof of \cref{pf:thmS:everyslot} and \cref{pf:thmS:lg1}]
A row $\sfname{Row}(Q,\mathrm{Lk},H)$ holds as soon as $H(\mathcal A,\sigma)$ holds for every class $\mathcal A$ and
every raw selection $\sigma$ (take for the bundle property of \cref{def:families:rows} the true property; this is
the shared lemma \lref{GlobalTheory.Fc.Row.of_sel}). By \cref{pf:thmS:atchoice} this is the case for the three
rows, for every $Q$ and $\mathrm{Lk}$; this proves \cref{pf:thmS:everyslot}. For \cref{pf:thmS:lg1} take for $R$ the
slot containing every core: the large-base condition holds trivially for it
(\lref{GlobalTheory.Fc.largeBaseProperty_true}), and the three rows hold at the slot $Q'\wedge R$ by
\cref{pf:thmS:everyslot}.
\end{proof}

% ---------------------------------------------------------------------------
\subsection{Theorem S along families}\label{sec:pf:thmS:fam}

\begin{theorem}[\cref{thm:theorem-S}]\label{pf:thmS:main}
Assume $\sfname{SWXS}$, $\sfname{RowNGP}$ and $\mathrm{FD}$. Let $\mathcal C$ be a class and $\sigma$ a
spectral selection for $\mathcal C$ such that every datum $D_{x,i}$ with $F\in\mathcal C$, $x>x_1$, $x\ge x_0(F)$ is
locally tied, the rank row $\mathcal{R}_{3}(\mathcal C)$ holds and $\mathcal{L}_{3}(\mathcal C,\sigma)$
holds. Then $\thmSClosed(\mathcal C,F\mapsto\mathcal{S}(F,3,\sigma(F)))$: closedly on $\mathcal C$,
eventually in $x$, under the admissible bounds, \eqref{eq:pf:thmS:S} holds for all $i\in I_x$ and all distinct
words $u\ne u'$ of $D_{x,i}$.
\formallabel{GlobalTheory.CarrierGap.offDiag}
\end{theorem}

Here $\sfname{SWXS}$ is the Siegel--Weil volume identity for the CM and hermitian steps of every Hida setting together
with a first-Chern compatibility of stabilisers, $\sfname{RowNGP}$ is the bound for the non-generic terms
(\cref{prop:families:volratio}(iii)), and $\mathrm{FD}$ is the finite-dimensionality of the spaces of
automorphic forms cut out by a locally constant idempotent on a fundamental domain. ``Locally tied'' means that at
every place of the auxiliary set, the local data of $D$ are isomorphic to the canonical local data.

The proof combines the pointwise form of Theorem~S (\cref{sec:pf:thmS:pw}) with the following statements.

\begin{definition}[Pointwise Theorem S]\label{pf:thmS:distinct}
$\mathcal{S}^{\mathrm{pt}}$ is the statement: for every $C_{\mathrm{thr}}$ there are $C_4\ge1$,
$C_A\ge\max(1,C_{\mathrm{thr}})$ and $0\le C\le C_A$ such that for all $C_0\ge1$, $C_G\ge0$, $C_\Delta>4$, $c_T>0$,
$C_b>0$ with $c_TC_b>1$ there is $k_0\in\N$ such that for all $C_g,c_b,C_w,c_{Tf},C_T,C_{0f}$ with $c_b>0$, $c_{Tf}>0$
there are $C_n\ge0$, $C'_0$, $C_{l,0}$ such that for all $C'\ge C'_0$ and $C_l\ge C_{l,0}$ the following holds for
every datum $D$, size data $L$, spectral ties $T=(X,Y,\dots)$ for $D$ and distinct words $u\ne u'$: if $D$ is tied,
$b\ge0$, the size hypotheses $\mathrm{Sz}(C_A,C_0,C_G,C_\Delta,c_T,C_b,C_n,k_0;D,L,T,u,u')$, the count
hypotheses $\mathrm{CH}_0(D,L;C',C_l,C_g,c_b,C_w,c_{Tf},C_T,C_{0f})$, the standard conditions
$\mathrm{St}(D,u,u')$ on the level and type primes, the hypotheses $\mathrm{H}_{K}$, the gap hypothesis
$\mathrm{Gap}(D)$ ($r\ge1801$ and consecutive Harish-Chandra parameters differ by at least $16r^{19}$), local ties,
the trace frame $1\le T_w\kappa_S(D)^{2r_S}$ for all words $w$, and the two loop-row bounds $\mathrm{R}_{\mathrm{big}}$,
$\mathrm{R}_{\mathrm{nsc}}$ for $Y.t_O(u,u')$ with $(\thmSepsL(D),\thmSdC(u,u'))$ hold, then \eqref{eq:pf:thmS:S} holds.
\end{definition}

\begin{lemma}[Pointwise to distinct words]\label{pf:thmS:distinctlemma}
The conclusion of \cref{pf:thmS:pointwise} below implies $\mathcal{S}^{\mathrm{pt}}$.
\formallabel{Formalism.L2.SG.ThmSOwn.theoremS_distinct_of_closed}
\end{lemma}

\begin{proof}
\Cref{pf:thmS:pointwise} has the same quantifier structure and hypotheses, and its conclusion is the conjunction of
\eqref{eq:pf:thmS:S} with a second inequality; keep the first conjunct.
\end{proof}

\begin{lemma}[Selections are admissible]\label{pf:thmS:selok}
Assume $\mathrm{Fin}$ and Wu decay $\mathrm{Wu}$ (for every datum $D$ with $\mathrm{Gap}(D)$ and every
Hida setting $X$ for $D$, the Wu decay estimate holds and the archimedean matrix-coefficient function of $X$ is
integrable). Let $\sigma$ be a spectral selection for $\mathcal C$ whose selected data are locally tied, and assume the
trace-frame row $\thmSClosed(\mathcal C,\mathrm{TF}(\cdot;C_{\mathrm{YL}}+1,\tfrac14,3,\tfrac12,C_{\mathrm{Tu}},C_0))$.
Then closedly on $\mathcal C$, eventually in $x$ and under the admissible bounds, every selected $\sigma(F)_{x,i}$ is
\emph{admissible} (its pair $(X,Y)$ is concretely tied, and $D_{x,i}$ is locally tied and satisfies
the trace frame).
\formallabel{Formalism.L2.SG.ThmSOwn.selOK_of}
\end{lemma}

\begin{proof}
Apply the closing lemma \lref{Formalism.L2.SG.ThmSOwn.closed_with} to the trace-frame row, with thresholds $0$ for
$C_\Gamma$ and $29$ for $s_{C'}$, $s_{C_l}$. Under the admissible bounds with $s_{C'}(F),s_{C_l}(F)\ge29$, the standing hypotheses
give $\mathrm{Gap}(D_{x,i})$ (the shared lemma \lref{Formalism.L2.SG.ThmSOwn.gap_hypothesis_of_hypotheses}:
the standing hypotheses with $C_l\ge16$ and $C'\ge29$ imply the gap hypothesis); the
trace frame holds at $D_{x,i}$ by the row; and the shared lemma \lref{Formalism.L2.SG.ThmSOwn.sTiesOK_of} states
that $\mathrm{Fin}$, Wu decay, the gap hypothesis, local ties and the trace frame make spectral ties admissible.
\end{proof}

\begin{lemma}\label{pf:thmS:finxc}
Wu decay and $\mathrm{FD}$ imply $\mathrm{Fin}$ (finiteness of the spaces of the full systems of all spectral
ties for data with the gap hypothesis).
\formallabel{Formalism.L2.SG.Pf.finXC_of}
\end{lemma}

\begin{proof}
This is the shared statement \lref{Formalism.L2.SG.Pf.FinXP.finX_holds}, applied with Wu decay,
$\mathrm{FD}$, the finiteness of Levi quotients \lref{Formalism.L2.SG.Pf.FinLeviPf.finLeviC_holds} and
the kernel statement \lref{Formalism.L2.SG.Pf.KernelOfEmbedding.IotaKer_holds}.
\end{proof}

\begin{lemma}[Size rows from the rank row]\label{pf:thmS:tf}
Let $C_0\ge1$ and assume $\mathcal{R}_{C_0}(\mathcal C)$. Then:
\begin{enumerate}
\item $\thmSClosed(\mathcal C,\mathrm{TF}(\cdot;C_{\mathrm{YL}}+1,\tfrac14,3,\tfrac12,C_{\mathrm{Tu}},C_0))$;
\item closedly, eventually and under the admissible bounds, the size rows hold at $D_{x,i}$: the
Schur formal-degree identity, the mass formulas of $\mathrm{MF}$, the split size condition, and
the lower bound $\ge\frac12$ for the relative formal degrees at all infinite places;
\item with $a(C_0)=e^{3/4}C_0e$: closedly, $1\le a(C_0)$ and, eventually and under the admissible bounds,
$\mathrm{LR}(D_{x,i},a(C_0))$ holds, i.e.\ $T_u/T_{u'}\le a(C_0)$ and
$\mathrm{vand}(u'_\tau)\le a(C_0)^k\,\mathrm{vand}(u_\tau)$ for all words $u,u'$ and all embeddings $\tau$, where
$\mathrm{vand}$ is the Vandermonde product of the Harish-Chandra parameters.
\end{enumerate}
\formallabel{Formalism.L2.SG.Pf.SzHypFam.hTF_of_ranks}
\formallabel{GlobalTheory.tFrame_of_ranksRow_three}
\formallabel{Formalism.L2.SG.Pf.SzHypFam.hRw_of_ranks}
\formallabel{Formalism.L2.SG.Pf.SzHypFam.hA_of_ranks}
\end{lemma}

\begin{proof}
In each case apply \lref{Formalism.L2.SG.ThmSOwn.closed_with} to the rank row, with thresholds $40$, $1$ and $40$
for $s_{C_l}$. (For $C_0=3$ the auxiliary propositions taken as instance arguments are theorems of the development,
\cref{pf:thmS:rem1}(2); this is the form of (1) used for the trace-frame row.) (1) Under $s_{C_l}(F)\ge40$, the trace frame is deduced from the rank formulas by
\lref{Formalism.L2.SG.Pf.SzHypFam.tFrame_of_ranks}: the Siegel mass ratio for CM fields unramified over their
real subfield and the rank-two mass bound give the first part of the family Bessel bound, the critical-level lower
bound gives the local split exponents, and the level frame lemma combines them. (2) Under $s_{C_l}(F)\ge1$ one has
$1\le s_{C_l}(F)\,L$, and \lref{Formalism.L2.SG.Pf.SzHypFam.szRowsAt_of_ties} gives the four size rows from the
ties of $D$, the standing data and hypotheses, and the mass formulas. (3) $a(C_0)\ge1$ since $C_0\ge1$; the bound
$\mathrm{LR}$ is \lref{Formalism.L2.SG.Pf.SzHypFam.sizes3b_of_ranks}, from the critical-level ratio and the
level-ratio bound for the lowest $K$-types.
\end{proof}

We also use two closed statements about the thresholds and sizes of selected data, both deferred:
\lref{Formalism.L2.SG.Pf.ThresholdLevels.thrRow_holds} states that for every $C_A\ge1$, closedly, eventually and
under the admissible bounds with $C_0$, for all distinct words $u\ne u'$ the threshold set of the data at places has
a largest element of the form $\Theta^{-f}$ with $\Theta>0$ (the threshold condition at places); and
\lref{Formalism.L2.SG.Pf.SzHypOf.sizesRowS_of} states that the size rows of \cref{pf:thmS:tf}(2) and these
threshold rows imply the closed size row $\mathrm{SR}_{C_0}(\mathcal C,\sigma)$: for all $C_A\ge1$ and all
functions $k_0,c_n$ of lists of reals, closedly, eventually and under the admissible bounds, the size hypotheses
$\mathrm{Sz}(C_A,C_0,C_\Gamma,5,\tfrac12,3,c_n(\cdot),k_0(\cdot);D,L,\sigma(F)_{x,i},u,u')$ hold for all $u\ne u'$.

\begin{proposition}[Passage to families]\label{pf:thmS:transport}
Assume $\mathcal{S}^{\mathrm{pt}}$ and $C_0\ge1$. Let $\sigma$ be a spectral selection for $\mathcal C$ such that,
closedly, the selected spectral ties are admissible (\cref{pf:thmS:selok}); assume the size row
$\mathrm{SR}_{C_0}(\mathcal C,\sigma)$, a function $a_F$ with $\thmSClosed(\mathcal C,F\mapsto 1\le
a_F(c_{\mathrm{out}})\wedge\thmSEv_F(\text{admissible bounds}\Rightarrow\mathrm{LR}(D,a_F(c_{\mathrm{out}}))))$, the
loop rows $\mathcal{L}_{C_0}(\mathcal C,\sigma)$, and absolute convergence of all tied term systems
($\mathrm{AC}$). Then $\thmSClosed(\mathcal C,F\mapsto\mathcal{S}(F,C_0,\sigma(F)))$.
\formallabel{Formalism.L2.SG.ThmSOwn.offDiagonalK0_of_rows}
\end{proposition}

\begin{proof}
Apply $\mathcal{S}^{\mathrm{pt}}$ with $C_{\mathrm{thr}}=1$; it gives $C_A\ge1$. For a list $l$ of reals let
$C_G(l)$ be its entry of index $17$ (for $l=c_{\mathrm{out}}\mathbin{+\!\!+}c_{\mathrm{mid}}(C_H)$ this entry is
$C_\Gamma(F)\ge0$). Applying the statement with $C_0$, $C_G=C_G(l)$, $C_\Delta=5$, $c_T=\frac12$, $C_b=3$ (so
$c_TC_b>1$) and with the count constants $(C_g,c_b,C_w,c_{Tf},C_T,C_{0f})=(l_6+1,\frac14,3,\frac12,l_{13},C_0)$ (the
entries $6$ and $13$ of the list are $C_{\mathrm{YL}}$ and $C_{\mathrm{Tu}}$) defines functions $k_0(l)$, $c_n(l)$,
$c'(l)$, $c_l(l)$ such that the conclusion of the pointwise theorem holds for all $C'\ge c'(l)$, $C_l\ge c_l(l)$
(for lists with negative entry $17$ the functions are set to $0$). Take the conjunction of the closed statements:
admissibility, the size row $\mathrm{SR}_{C_0}$ with $(C_A,k_0,c_n)$, the loop rows with $a_F$, and the bound
$\mathrm{LR}$; and apply \lref{Formalism.L2.SG.ThmSOwn.closed_with} with the threshold $0$ for $C_\Gamma$ and
$\max(c'(l),c_l(l),29)$ for $s_{C'}$, $s_{C_l}$. Fix $F$, $C_H$ in the resulting situation, choose $x_0$ larger than
the four thresholds ``eventually'' of the statements, and let $x\ge x_0$, $i\in I_x$, $u\ne u'$, under the admissible
bounds. Then: the selected $(X,Y)$ is concretely tied (admissibility); the count hypotheses
$\mathrm{CH}$ hold with $a_F(c_{\mathrm{out}})$, by \lref{Formalism.MainNotions.count_hypotheses_of_sizes}
(the admissible bounds, $s_{C_l}(F)\ge29$ and $\mathrm{LR}$); the term system is absolutely convergent
($\mathrm{AC}$); so the loop rows give $\mathrm{R}_{\mathrm{big}}$ and $\mathrm{R}_{\mathrm{nsc}}$ with
$(\thmSepsL(D),\thmSdC(u,u'))$. Further $b\ge2\ge0$; the size row gives $\mathrm{Sz}$; the conditions
$\mathrm{St}(D,u,u')$ and $\mathrm{H}_{K}$ follow from the standing data and hypotheses
(\lref{Formalism.MainNotions.Stand_of} and \lref{Formalism.MainNotions.NonGenericAndKTwoTheorems.Hyp_of}); the gap
hypothesis follows from the standing hypotheses since $s_{C'}(F)\ge29$; local ties and the trace frame are part of
admissibility; and $D_{x,i}$ is tied (a field of the family). The pointwise statement with $C'=s_{C'}(F)$,
$C_l=s_{C_l}(F)$ gives \eqref{eq:pf:thmS:S}.
\end{proof}

\begin{lemma}\label{pf:thmS:absconv}
$\mathrm{AC}$ holds: for every datum $D$, every Hida setting $X$ and full system $Y$ which are concretely tied
and all words $u,u'$, the term system $Y.t_O(u,u')$ is absolutely convergent.
\formallabel{Formalism.L2.SG.ThmSOwn.absConvTied_holds}
\end{lemma}

\begin{proof}
This is the shared implication \lref{Formalism.L2.SG.RowsConc.absConvTied_of_full} (folding, existence of
fundamental sets, integrability of the orbital integrals and absolute convergence for full systems imply
$\mathrm{AC}$), applied to the four statements
\lref{Formalism.L2.SG.Closed.FoldOKC_outright}, \lref{Formalism.L2.SG.Closed.FundExistsC_holds}, \lref{Formalism.L2.SG.Closed.OOIntAllC_outright}, and \lref{Formalism.L2.SG.Closed.AbsConvC_outright}
proved in the development.
\end{proof}

\begin{proof}[Proof of \cref{pf:thmS:main}]
Wu decay holds (\lref{CompanionTheory.Lime.BM.wuDecayOfGap_holds'}), hence $\mathrm{Fin}$ by
\cref{pf:thmS:finxc}. The rank row gives the trace-frame row (\cref{pf:thmS:tf}(1) with $C_0=3$), hence
admissibility of the selection (\cref{pf:thmS:selok}). The nineteen hypotheses (H1)--(H19) of the pointwise theorem are supplied by \lref{GlobalTheory.CarrierGap.theoremS_leaves_of_carrierG}, which derives
them from the carrier bound (\lref{GlobalTheory.CarrierGap.carrierP}: a lower bound for Poincar\'e amplitudes), $\sfname{RowNGP}$,
$\mathrm{FD}$, uniqueness of Bessel models with large gap (\lref{CompanionTheory.Lime.BM.uniqGap_holds})
and $\sfname{SWXS}$; so \cref{pf:thmS:pointwise,pf:thmS:distinctlemma} give $\mathcal{S}^{\mathrm{pt}}$. The size
row $\mathrm{SR}_{3}$ follows from \cref{pf:thmS:tf}(2) and the threshold row
(\lref{Formalism.L2.SG.Pf.ThresholdLevels.thrRow_holds} and \lref{Formalism.L2.SG.Pf.SzHypOf.sizesRowS_of}); the bound
$\mathrm{LR}$ with $a_F=a(3)$ is \cref{pf:thmS:tf}(3); the loop rows are assumed; and $\mathrm{AC}$ is
\cref{pf:thmS:absconv}. \Cref{pf:thmS:transport} with $C_0=3$ concludes. (The auxiliary statements used as instance arguments by the lemmas of this subsection are theorems of the development, see \cref{pf:thmS:rem1}(2).)
\end{proof}

% ---------------------------------------------------------------------------
\subsection{The pointwise theorem and its reduction to term systems}\label{sec:pf:thmS:pw}

The pointwise theorem has nineteen hypotheses (H1)--(H19), supplied by
\lref{GlobalTheory.CarrierGap.theoremS_leaves_of_carrierG}. We list them by their role; the terms ``datum'',
``spectral ties'', ``term system'' are those of \cref{sec:pf:thmS:conv,sec:pf:thmS:abs}.
\begin{enumerate}[label=(H\arabic*)]
\item\label{pf:thmS:H1} (Local models.) For every datum $D$, the abstract kernels attached to $D$ have the local
model of \cref{pf:thmS:localmodel} below.
\item (Density inputs.) For every locally tied datum, in particular the density input
$\mathrm{DI}(k,r_S,f,s_d)$ of \cref{pf:thmS:mainrowthm} for the density $s_d$ of the data, and positivity and
boundedness of the local factors $\alpha_p$.
\item (Row inputs.) For every datum, positive semi-definiteness of the top kernel minus the row coefficients
(see $\mathrm{RI}$ in \cref{pf:thmS:mainrowthm}).
\item (Non-generic terms.) $\sfname{RowNGP}$, the bound for the non-generic pairs (\cref{prop:families:volratio}(iii)).
\item (Degenerate terms.) For distinct words, under $\mathrm{H}_{K}$ and a concrete tie, every degenerate term
of $Y.t_O(u,u')$ has operator $0$.
\item (Carrier bound.) The carrier lower bound for Poincar\'e amplitudes (\lref{GlobalTheory.CarrierGap.carrierP}).
\item (Kernel expansion.) The kernel relating the words is the sum of the term operators.
\item[(H8--10)] Folding, integrability of the orbital integrals and absolute convergence for full systems.
\item[(H11)] (Finiteness.) $\mathrm{Fin}$: the spaces of the full systems of all spectral ties of data with the gap hypothesis are finite.
\item[(H12)] (Wu decay.) $\mathrm{Wu}$ (\cref{pf:thmS:selok}).
\item[(H13)] (Choices.) Existence of point representatives, a pinned archimedean measure, base vectors,
local splitting data, uniformisers at the auxiliary places, a row labelling and a pair system.
\item[(H14--16)] (Formula, row sums, row forms.) The counting formula for the terms,
the resulting sum of rows, and the row forms with their normalising constants.
\item[(H17)] (Mass data.) The thirteen identities $\mathrm{MI}$ of \cref{pf:thmS:massdata} for the
mass data of the terms.
\item[(H18)] (Links.) The link $\mathrm{LC}$ of \cref{pf:thmS:pbfromlocal} between the term system and the
mass data.
\item[(H19)] (Split inputs.) The inputs $\mathrm{SI}$ of \cref{pf:thmS:pbfromlocal} (Gram expansion, Gram
bound, row orthogonality, positive corank, vanishing of the identity part, $\Gamma(i,i)=\Gamma_1$).
\end{enumerate}

\begin{theorem}[Pointwise Theorem S]\label{pf:thmS:pointwise}
Assume (H1)--(H19). Then the statement $\mathcal{S}^{\mathrm{pt}}$ of
\cref{pf:thmS:distinct} holds with the conclusion strengthened to
\[
\|K_0(u,u')\|^2\le\bigl(1+b^{-6}\bigr)\frac{d^S_ud^S_{u'}}{D_*}\quad\text{and}\quad
\mathrm{hs}(u,u')\le\bigl(1+b^{-6}\bigr)\frac{d^S_ud^S_{u'}}{D_*},
\]
where $\mathrm{hs}(u,u')$ is the Hilbert--Schmidt square of the block $(u,u')$ recorded in the size data $L$.
\formallabel{Formalism.L2.SG.TheoremSConclusion.theoremS_of_conj}
\formallabel{Formalism.L2.SG.PoincareAmplitudeLowerBounds.theoremS_closedP}
\formallabel{Formalism.L2.SG.PoincareAmplitudeLowerBounds.theoremS_rec3P}
\formallabel{Formalism.L2.SG.PoincareAmplitudeLowerBounds.theoremS_rec2P}
\formallabel{Formalism.L2.SG.SpectralTiesReal.theoremS_rec}
\end{theorem}

The proof passes through three intermediate statements in which more and more of the hypotheses are replaced by the
explicit data they provide; we describe them as steps.

\begin{lemma}[Hilbert--Schmidt mass and leading coefficient]\label{pf:thmS:packagetie}
Let $D$ be a datum with size data $L$, $b\ge0$, and assume the isometry of norms of the development. Let $X$, $Y$ be a
Hida setting and a full system with finite spaces, isomorphic to the spectral data of $D$, and $u\ne u'$ words with
the dimension ties, the $K$-expansion, absolute convergence, identity terms of weight $0$, vanishing degenerate part
and the norm tie. If for some global size datum $G$ and parameters with the size tie one has $\|K_0(u,u')\|^2\le(1+b^{-6})\,d_ud_{u'}/D_*(G)$, then
$\mathrm{hs}(u,u')\le(1+b^{-6})\,d^S_ud^S_{u'}/D_*$.
\formallabel{Formalism.L2.SG.PackageTie2.packageTieConc}
\end{lemma}

\begin{proof}
By the norm tie (applicable under the dimension ties, the expansion, absolute convergence and the vanishing of the
identity-degenerate and degenerate parts) and the isometry of norms, $\mathrm{hs}(u,u')=\|K_0(u,u')\|^2$. The size tie
gives $d_u=d^S_u$, $d_{u'}=d^S_{u'}$ and $D_*(G)=D_*$, and $b^{-6}=(b^6)^{-1}$ for $b\ge0$.
\end{proof}

\begin{proof}[Proof of \cref{pf:thmS:pointwise}]
\emph{Step 1: unpacking the sizes.}
The size hypotheses $\mathrm{Sz}$ provide a bijection $e$ between the infinite places of $E$ and the embedding
types $\tau$, a type $\tau$, constants $C_p,\Theta,C_f$ with $\Theta>0$, and for the global size datum $G=g(D,e)$
attached to $D$ and the word data $d_e=\mathrm{de}(D,\tau,u)$, $d'_e=\mathrm{de}(D,\tau,u')$, $l=l(d_e,d'_e)$: $k_0\le
G.k$, the size hypotheses $\mathrm{SH}(G,C_A,C_0,C_G,C_\Delta,c_T,C_b,b,C_p,\Theta,d_e,d'_e,T_u,T_{u'},d^S_u,
d^S_{u'},l)$, the inequalities $16\le c_{\exp}(C_\Delta,c_T,C_b)\,C_f$, $C_f\log b\le G.f\log G.k$, $2r_S+2+f+5\le k$,
the numerical conditions $\mathrm{SN}(G.k,G.r,G.f,b,G.g_1,G.g_2,\frac13,C_n)$ (\cref{pf:thmS:sizenumbers}),
$D_*(G.\mathrm{size}(d_e,T_u,d^S_u))=D_*$ and $T_u\equiv T_{u'}\bmod g_S(D)^2$. By $\mathrm{Fin}$ and Wu decay the
full system of $T$ has finite spaces and the archimedean coefficient is integrable; with the level-subgroup properties,
neatness of the level subgroups (shared statements of the development), local ties and the trace frame, the pair
$(T.X,T.Y)$ is concretely tied. (H13) provides the choices, and (H14) the counting formula for them. Since $(kb)^8\le g_1,g_2$ with $k\ge32$ and $b\ge16$ (from $\mathrm{SN}$), both type
primes satisfy $g_i\ge3$. Folding, integrability and absolute convergence give absolute convergence of the term
system; (H5) gives the vanishing of its degenerate part (the shared lemma
\lref{Formalism.L2.SG.RowsOther.kppOff_of_termZero}); identity terms have weight $0$ (the shared lemma
\lref{Formalism.L2.SG.Pf.SatGram.idDegenerate_tO}). Together with (H7) and (H15)--(H19) at these
choices, all hypotheses of Step~2 hold.

\emph{Step 2: the carrier bound.}
The carrier bound (H6) provides $C_4\ge1$, $C_{A,0}$ and, for every $C_A\ge C_{A,0}$, the carrier estimate for the data of Step~1 beyond a further threshold for $k$. One applies Step~3 with
$C_{\mathrm{thr}}$ replaced by $\max(C_{\mathrm{thr}},C_{A,0})$ and with $k_0$ the larger of the two thresholds.

\emph{Step 3: the non-generic row.}
The carrier record identifies the global size datum: $G.k=k$, $G.r=r_S$, $G.f=f$. The concrete tie gives
$\dim C_u=d^S_u$, $\dim C_{u'}=d^S_{u'}$, $D_*(X)=D_*$; hence $\thmSdC(u,u')=d^S_ud^S_{u'}/D_*(G.\mathrm{size})$.
$\sfname{RowNGP}$ provides $C_4'\ge0$ and the non-generic row bound with
$\varepsilon_{\mathrm{NG}}=2^{-(k-2r_S-5)}\exp(C_4'(2r_S+2)^5f\log b)$; local ties follow from the concrete tie; and
$\thmSepsL(D)=\exp(-\frac13 kf\log k)$. So the three row bounds of Step~4 hold with $c_5=\frac13$ and $C_n=C'_4$.

\emph{Step 4: instantiation.} The carrier record contains a
structure $S$ from which the abstract data of \cref{sec:pf:thmS:abs,sec:pf:thmS:main} are built: a TB-datum $D^\flat$ over the field $E$, built from $S$ and the data of $Y$, (\cref{pf:thmS:tbdata}) which is an instance
$\mathrm{Inst}(D^\flat,S.\mathrm{pl},S.G,C_A,S.d_e,S.d'_e,T_u,T_{u'},S.l)$ of the sizes of $G$; (H1) gives its local model; (H2) the density input; (H3), together with the $K_V$-identity and its consequences $K_V\ge0$ and the bound for $K_V$, gives the row input $\mathrm{RI}$ for the row coefficients $c$; and the energy link
$\mathrm{En}(B,\mathrm{jn},s_d,C,\Theta,r,f)$ for the mass data $B$ of the terms is
\lref{Formalism.L2.SG.MassDataConc.energyLink_bCF}. $E$ is unramified over its maximal real subfield
(a property of the host). The term system $Y.t_O(u,u')$ satisfies $\mathrm{LC}$ and
$\mathrm{SI}$ with the mass data, and the three row bounds of Step~3. Now \cref{pf:thmS:pbapplied} gives
$\|K_0(u,u')\|^2\le(1+b^{-6})\,d^S_ud^S_{u'}/D_*(G.\mathrm{size})$, which is the first inequality since
$D_*(G.\mathrm{size})=D_*$; \cref{pf:thmS:packagetie} (with the dimension ties, the norm tie and the isometry of
norms, which are shared statements of the development) gives the second.
\end{proof}

% ---------------------------------------------------------------------------
\subsection{The spectral argument for term systems}\label{sec:pf:thmS:abs}

\begin{definition}[Term systems]\label{pf:thmS:terms}
A \emph{term system} $\mathcal X$ consists of a set $\mathcal T$ of terms, complex inner product spaces $H$ and $H_1$,
maps $\mathrm{op}:\mathcal T\to H$ and $\mathrm{vec}:\mathcal T\to H_1$, a set of rows with a map
$\thmSrow:\mathcal T\to\mathrm{Rw}$, a finite set $I$ (of \emph{types}) with a map $\iota:\mathrm{Rw}\to I$,
predicates $\mathrm{isId}$ and $\mathrm{gen}$ on $\mathcal T$, a weight $G:\mathcal T\to\R$, a matrix $\Gamma:I\times
I\to\C$, a real number $\Gamma_1$, a corank function $s_0:\mathcal T\times\mathcal T\to\N$ and a relation
$\mathrm{scal}$ on $\mathcal T$. Put
\begin{gather*}
\thmSId=\{\mathrm{isId}\},\quad \thmSDeg=\{G=0,\ \neg\mathrm{isId}\},\quad \thmSNG=\{G\ne0,\ \neg\mathrm{gen}\},\\
\thmSGen=\{G\ne0,\ \mathrm{gen}\},\quad \thmSNZ=\{G\ne0\},\quad \Phi(S)=\Gamma_1\sum_{i\in I}\|\thmSvp(S,i)\|^2,\\
\thmSpart(S)=\sum_{t\in S}\mathrm{op}(t),\qquad \thmSvp(S,i)=\sum_{t\in S,\ \iota(\thmSrow t)=i}\mathrm{vec}(t),\\
K=\sum_t\mathrm{op}(t),\qquad K_0=\thmSpart(\thmSId)+\thmSpart(\thmSGen)+\thmSpart(\thmSNG)
\end{gather*}
(sums in the sense of unconditional summation). For $S\subseteq\mathcal T$ let $P(S)$ be the set of pairs
$(t,t')\in S\times S$ with $\thmSrow t=\thmSrow t'$, and $\mathrm{big}(S)$, $\mathrm{one}(S)$, $\mathrm{nsc}(S)$,
$\mathrm{sc}(S)$ the subsets of $P(S)$ where $s_0(t,t')\ge2$, resp.\ $s_0=1$, resp.\ $s_0=1$ and not
$\mathrm{scal}(t,t')$, resp.\ $s_0=1$ and $\mathrm{scal}(t,t')$. For a set $Q$ of pairs put
\[
\thmSps(Q)=\Gamma_1\sum_{(t,t')\in Q}\Re\langle\mathrm{vec}\,t,\mathrm{vec}\,t'\rangle,\qquad
\thmSpa(Q)=\Gamma_1\sum_{(t,t')\in Q}|\langle\mathrm{vec}\,t,\mathrm{vec}\,t'\rangle|.
\]
The system is \emph{absolutely convergent} if $\sum_t\|\mathrm{op}(t)\|<\infty$, $\sum_t\|\mathrm{vec}(t)\|<\infty$ and
$\Gamma_1>0$. We use the following conditions:
\begin{itemize}
\item $\mathrm{ID}$: $G(t)=0$ for identity terms; $\mathrm{IZ}$: $\thmSpart(\thmSId)=0$;
\item $\mathrm{GE}(S)$: $\|\thmSpart(S)\|^2=\Re\sum_{i,j\in I}\Gamma(i,j)\langle\thmSvp(S,i),\thmSvp(S,j)\rangle$;
\item $\mathrm{GB}(P_g)$: the matrix $P_g\Gamma_1\cdot\one-\Gamma$ is positive semi-definite;
\item $\mathrm{RO}$: $\langle\mathrm{vec}\,t,\mathrm{vec}\,t'\rangle=0$ if $\iota(\thmSrow t)=\iota(\thmSrow t')$ and
$\thmSrow t\ne\thmSrow t'$;
\item $\mathrm{CP}$: $s_0(t,t')\ge1$ if $\thmSrow t=\thmSrow t'$, $G(t)\ne0$ and $G(t')\ne0$;
\item the \emph{row bounds}, for $\varepsilon,A,dd\in\R$: $\mathrm{R}_{\mathrm{big}}(\varepsilon,dd)$:
$\thmSpa(\mathrm{big}(\thmSNZ))\le\varepsilon\,dd$; $\mathrm{R}_{\mathrm{nsc}}(\varepsilon,dd)$:
$\thmSpa(\mathrm{nsc}(\thmSGen))\le\varepsilon\,dd$; $\mathrm{R}_{\mathrm{main}}(A,dd)$: $\thmSps(\mathrm{sc}(\thmSGen))\le A\,dd$;
$\mathrm{R}_{\mathrm{ng}}(\varepsilon,dd)$: $\thmSps(\mathrm{one}(\thmSNG))\le\varepsilon\,dd$.
\end{itemize}
\end{definition}

For the term systems $Y.t_O(u,u')$ of \cref{sec:pf:thmS:conv}, $H$ is the Hilbert--Schmidt space of maps between the
column spaces of $u'$ and $u$; $\mathrm{op}(t)$ is the block operator whose $(\rho,\rho')$ entry, for columns $\rho$,
$\rho'$ over $\Lambda$, $\Lambda'$, is a normalised value of the level character at the finite part of $t$ times the
orbital integral $\int_{\Omega_t}\phi_u(x)\otimes\overline{\phi_{u'}(\gamma^{-1}x)}\,dx$ of rank-one operators over a
fundamental set $\Omega_t$; $H_1$ is a sum of Hilbert--Schmidt spaces indexed by the pairs $(\Lambda,\Lambda')$; the row
of $t$ is $(\Lambda,\Lambda',\iota(t))$ with $\iota(t)\in I=\prod_v\GL_k(\F_{q_v})$ the reduction of the finite part of
$t$ at the places above the type primes; $\Gamma(i,j)$ is the complex conjugate of a product of character sums
evaluated at $j^{-1}i$, and $\Gamma_1$ the real part of this product at $1$; $G(t)=1$ or $0$ according as an invariant
of $t$ is non-zero or zero; and $s_0$, $\mathrm{scal}$ compare the frames of two terms with the same lattices. We do not
use this description: the argument below only uses the axioms listed.

\begin{lemma}[Gram sums]\label{pf:thmS:gram}
Let $I$ be finite, $\Gamma:I\times I\to\C$, $c\in\R$ and $a:I\to H_1$. If $c\cdot\one-\Gamma$ is positive
semi-definite, then $\Re\sum_{i,j}\Gamma(i,j)\langle a_i,a_j\rangle\le c\sum_i\|a_i\|^2$.
\formallabel{Formalism.FourierDecayAlt.P.GrowthPolynomialOperator.t_EqualIota}
\end{lemma}

\begin{proof}
Write $M=c\cdot\one-\Gamma=\sum_lv_lv_l^*$, i.e.\ $M(i,j)=\sum_lv_l(i)\overline{v_l(j)}$. With inner products linear in
the second variable, $\sum_{i,j}M(i,j)\langle a_i,a_j\rangle=\sum_l\|\sum_i\overline{v_l(i)}a_i\|^2\ge0$, and
$\sum_{i,j}M(i,j)\langle a_i,a_j\rangle=c\sum_i\|a_i\|^2-\sum_{i,j}\Gamma(i,j)\langle a_i,a_j\rangle$. Take real parts.
\end{proof}

\begin{lemma}[Price of the types]\label{pf:thmS:price}
If $\mathrm{GE}(S)$ and $\mathrm{GB}(P_g)$ hold, then $\|\thmSpart(S)\|^2\le P_g\,\Phi(S)$.
\formallabel{Formalism.FourierDecayAlt.P.GrowthPolynomialOperator.t_PriceOfIota}
\end{lemma}

\begin{proof}
Apply \cref{pf:thmS:gram} with $c=P_g\Gamma_1$ and $a_i=\thmSvp(S,i)$.
\end{proof}

\begin{lemma}[Pair expansion]\label{pf:thmS:pairsplit}
Let $\mathcal X$ be absolutely convergent with $H_1$ complete, $\mathrm{RO}$ and $\mathrm{CP}$. Then
$\Phi(S)=\thmSps(P(S))$ for every $S$, and
\begin{equation}\label{eq:pf:thmS:split}
\Phi(\thmSGen)=\thmSps(\mathrm{big}(\thmSGen))+\thmSps(\mathrm{nsc}(\thmSGen))+\thmSps(\mathrm{sc}(\thmSGen)),\qquad
\Phi(\thmSNG)=\thmSps(\mathrm{big}(\thmSNG))+\thmSps(\mathrm{one}(\thmSNG)).
\end{equation}
\formallabel{Formalism.FourierDecayAlt.P.GrowthPolynomialOperator.t_PairSplit}
\end{lemma}

\begin{proof}
For an absolutely summable family $(f_t)$ in a Hilbert space, $\|\sum_tf_t\|^2=\sum_{(t,t')}\Re\langle f_t,f_{t'}\rangle$,
the double sum being absolutely convergent since $|\langle f_t,f_{t'}\rangle|\le\|f_t\|\|f_{t'}\|$. Hence
$\|\thmSvp(S,i)\|^2$ is the sum of $\Re\langle\mathrm{vec}\,t,\mathrm{vec}\,t'\rangle$ over the pairs in $S\times S$
with $\iota(\thmSrow t)=\iota(\thmSrow t')=i$; summing over $i$ gives the sum over the pairs with
$\iota(\thmSrow t)=\iota(\thmSrow t')$. By $\mathrm{RO}$ the pairs with distinct rows contribute $0$, so
$\Phi(S)=\thmSps(P(S))$. For $(t,t')\in P(\thmSGen)$ or $P(\thmSNG)$ both weights are non-zero, so $s_0(t,t')\ge1$ by
$\mathrm{CP}$; the sets $\mathrm{big}$, $\mathrm{nsc}$, $\mathrm{sc}$ (resp.\ $\mathrm{big}$, $\mathrm{one}$)
therefore partition $P(\thmSGen)$ (resp.\ $P(\thmSNG)$), and \eqref{eq:pf:thmS:split} follows by additivity of
absolutely convergent sums.
\end{proof}

\begin{lemma}\label{pf:thmS:subfamily}
If $\mathcal X$ is absolutely convergent and $Q\subseteq Q'$, then $\thmSpa(Q)\le\thmSpa(Q')$ and
$|\thmSps(Q)|\le\thmSpa(Q)$.
\formallabel{Formalism.FourierDecayAlt.P.GrowthPolynomialOperator.t_SubFamily}
\end{lemma}

\begin{proof}
The family $|\langle\mathrm{vec}\,t,\mathrm{vec}\,t'\rangle|$ is summable over pairs (as above), $\Gamma_1>0$, and
$|\Re z|\le|z|$.
\end{proof}

\begin{lemma}\label{pf:thmS:textclosing}
In a normed space, if $dd\ge0$, $\|x\|^2\le(1+\eta+\varepsilon)\,dd$ and $\|y\|^2\le\varepsilon\,dd$, then
$\|x+y\|^2\le(\sqrt{1+\eta+\varepsilon}+\sqrt\varepsilon)^2\,dd$.
\formallabel{Formalism.FourierDecayAlt.P.GrowthPolynomialOperator.t_TextClosing}
\end{lemma}

\begin{proof}
$\|x+y\|\le\|x\|+\|y\|\le(\sqrt{1+\eta+\varepsilon}+\sqrt\varepsilon)\sqrt{dd}$.
\end{proof}

\begin{proposition}[Abstract pair bound]\label{pf:thmS:pba}
Let $\mathcal X$ be an absolutely convergent term system with $H_1$ complete, and let $P_g,A,\varepsilon,\varepsilon',dd$
be reals with $P_g\ge0$ and $dd\ge0$. Assume $\mathrm{GE}(\thmSGen)$, $\mathrm{GE}(\thmSNG)$,
$\mathrm{GB}(P_g)$, $\mathrm{RO}$, $\mathrm{CP}$, $\mathrm{IZ}$ and the row bounds
$\mathrm{R}_{\mathrm{big}}(\varepsilon,dd)$, $\mathrm{R}_{\mathrm{nsc}}(\varepsilon,dd)$, $\mathrm{R}_{\mathrm{main}}(A,dd)$,
$\mathrm{R}_{\mathrm{ng}}(\varepsilon',dd)$. Then
\begin{equation}\label{eq:pf:thmS:pba}
\|K_0\|^2\le\Bigl(\sqrt{P_g(A+2\varepsilon)}+\sqrt{P_g(\varepsilon+\varepsilon')}\Bigr)^2\,dd.
\end{equation}
\formallabel{Formalism.FourierDecayAlt.P.GrowthPolynomialOperator.t_PairBoundAbstract}
\end{proposition}

\begin{proof}
Since $\mathrm{big}(\thmSNG)\subseteq\mathrm{big}(\thmSNZ)$, \cref{pf:thmS:pairsplit,pf:thmS:subfamily} give
\[
\Phi(\thmSNG)\le|\thmSps(\mathrm{big}(\thmSNG))|+\thmSps(\mathrm{one}(\thmSNG))
\le\thmSpa(\mathrm{big}(\thmSNZ))+\varepsilon'dd\le(\varepsilon+\varepsilon')\,dd.
\]
Similarly, since $\mathrm{big}(\thmSGen)\subseteq\mathrm{big}(\thmSNZ)$,
\[
\Phi(\thmSGen)\le\thmSpa(\mathrm{big}(\thmSNZ))+\thmSpa(\mathrm{nsc}(\thmSGen))+\thmSps(\mathrm{sc}(\thmSGen))
\le(A+2\varepsilon)\,dd.
\]
By \cref{pf:thmS:price}, $\|\thmSpart(\thmSGen)\|^2\le P_g(A+2\varepsilon)\,dd$ and
$\|\thmSpart(\thmSNG)\|^2\le P_g(\varepsilon+\varepsilon')\,dd$. By $\mathrm{IZ}$,
$K_0=\thmSpart(\thmSGen)+\thmSpart(\thmSNG)$, and \cref{pf:thmS:textclosing} (with
$\eta=P_g(A+2\varepsilon)-1-P_g(\varepsilon+\varepsilon')$ and $P_g(\varepsilon+\varepsilon')$ in place of
$\varepsilon$) gives \eqref{eq:pf:thmS:pba}.
\end{proof}

\begin{lemma}\label{pf:thmS:fourparts}
If $\mathrm{ID}$ holds, the four sets $\thmSId,\thmSDeg,\thmSNG,\thmSGen$ partition $\mathcal T$, and for an
absolutely convergent system with $H$ complete, $K=K_0+\thmSpart(\thmSDeg)$.
\end{lemma}

\begin{proof}
A term is an identity term, or else has weight $0$ (degenerate) or non-zero weight, and in the last case is generic or
not; identity terms have weight $0$ by $\mathrm{ID}$, so the sets are disjoint. The second claim follows by additivity of
absolutely convergent sums. Thus, when the degenerate part vanishes (as in Step~1 of the proof of
\cref{pf:thmS:pointwise}), $K_0$ is the whole operator.
\end{proof}

\paragraph{Numerics.}
\begin{definition}\label{pf:thmS:sizenumbers}
$\mathrm{SN}(k,r,f,b,g_1,g_2,c,C_n)$ means: $b\ge16$, $k\ge32$, $f\le k$, $4r+12\le k$, $(kb)^8\le g_1$,
$(kb)^8\le g_2$, $70\log b\le k$, $1\le5cf\log k$ and $8C_n(2r+2)^5f\log b\le k$.
\end{definition}

\begin{lemma}[Numerical inequalities]\label{pf:thmS:num}
Let $k,r,f\in\N$ and $b,g_1,g_2,c,C,C_n\in\R$.
\begin{enumerate}
\item If $b\ge1$, $k\ge1$, $f\le k$ and $(kb)^8\le g_1,g_2$, then $f/g_j\le b^{-7}$ and
$0\le(1+3/g_1)^f(1+3/g_2)^f\le e^{6b^{-7}}$.
\item If $b>0$ and $70\log b\le k$, then $e^{-k/5}\le b^{-14}$.
\item If $f\le k$, $k\ge32$ and $4r+12\le k$, then $12f\,2^{-(k-2r-2)}\le f\,2^{-k/2}\le e^{-k/5}$.
\item If $1\le5cf\log k$, then $\exp(-ckf\log k)\le e^{-k/5}$.
\item If $8Cm^5f\log b\le k$, then $2^{-k/2}\exp(Cm^5f\log b)\le e^{-k/5}$.
\item If $4r+10\le k$, then $2^{-(k-2r-5)}\le2^{-k/2}$.
\item If $b\ge16$, $0\le P_g\le e^{6b^{-7}}$, $A\le(1+b^{-7})e^{b^{-14}}$ and $\varepsilon,\varepsilon'\le b^{-14}$, then
$\bigl(\sqrt{P_g(A+2\varepsilon)}+\sqrt{P_g(\varepsilon+\varepsilon')}\bigr)^2\le1+b^{-6}$.
\end{enumerate}
\formallabel{Formalism.FourierDecayAlt.P.GrowthPolynomialOperator.t_NumPrice}
\formallabel{Formalism.FourierDecayAlt.P.GrowthPolynomialOperator.t_NumExp}
\formallabel{Formalism.FourierDecayAlt.P.GrowthPolynomialOperator.t_NumDens}
\formallabel{Formalism.FourierDecayAlt.P.GrowthPolynomialOperator.t_NumLoops}
\formallabel{Formalism.FourierDecayAlt.P.GrowthPolynomialOperator.t_NumNG}
\formallabel{Formalism.FourierDecayAlt.P.GrowthPolynomialOperator.t_NGExponent}
\formallabel{Formalism.FourierDecayAlt.P.GrowthPolynomialOperator.t_EtaZero}
\end{lemma}

\begin{proof}
(1) $fb^7\le kb^7\le(kb)^8\le g_j$, and $(1+3/g)^f\le e^{3f/g}\le e^{3b^{-7}}$. (2) $-k/5\le-14\log b$. (3) Write
$k=2r+2+j$; then $j\ge2r+10$, so $2^{-j}\cdot2^{k/2}=2^{-(j-2r-2)/2}\le2^{-4}<\frac1{12}$, which is the first
inequality. For the second, $f2^{-k/2}\le k2^{-k/2}=\exp(\log k-\frac k2\log2)$; from $\log(k/32)\le k/32-1$ and
$0.69<\log2<0.7$ one gets $\log k-\frac k2\log2\le5\log2-1+k(\frac1{32}-0.345)\le-\frac k5$ for $k\ge32$. (4)
$ckf\log k\ge k/5$. (5) $\exp(-\frac k2\log2+Cm^5f\log b)\le\exp(-0.345k+\frac k8)\le e^{-k/5}$. (6) $k-2r-5\ge k/2$.
(7) Put $u=b^{-7}\le16^{-7}<10^{-3}$. Then $e^{6u}\le1+7u$, $e^{u^2}\le1+u$, so $A+2\varepsilon\le(1+u)^2+2u^2\le1+3u$,
$P_g(A+2\varepsilon)\le(1+7u)(1+3u)\le1+11u$ and $P_g(\varepsilon+\varepsilon')\le2u^2(1+7u)\le3u^2$. Hence
$\sqrt{P_g(A+2\varepsilon)}\le1+\frac{11}2u$, $\sqrt{P_g(\varepsilon+\varepsilon')}\le\frac74u$, and
$(1+\frac{29}4u)^2\le1+16u\le1+b^{-6}$, the last step because $16u=16b^{-7}\le b^{-6}$.
\end{proof}

\begin{proposition}[Pair bound]\label{pf:thmS:pairbound}
Let $\mathcal X$ be an absolutely convergent term system with $H_1$ complete, and assume
$\mathrm{SN}(k,r,f,b,g_1,g_2,c,C_n)$, $dd\ge0$, $0\le\delta\le b^{-7}$, $0\le e_1\le
\exp(12f\,2^{-(k-2r-2)})$, the Gram expansions on $\thmSGen$ and $\thmSNG$, $\mathrm{GB}((1+3/g_1)^f(1+3/g_2)^f)$,
$\mathrm{RO}$, $\mathrm{CP}$, $\mathrm{IZ}$, and the row bounds
\begin{gather*}
\mathrm{R}_{\mathrm{big}}(\varepsilon,dd),\quad \mathrm{R}_{\mathrm{nsc}}(\varepsilon,dd)\quad\text{with }\varepsilon=e^{-ckf\log k},\qquad
\mathrm{R}_{\mathrm{main}}((1+\delta)e_1,dd),\\
\mathrm{R}_{\mathrm{ng}}\bigl(2^{-(k-2r-5)}e^{C_n(2r+2)^5f\log b},dd\bigr).
\end{gather*}
Then $\|K_0\|^2\le(1+b^{-6})\,dd$.
\formallabel{Formalism.FourierDecayAlt.P.GrowthPolynomialOperator.t_PairBound}
\end{proposition}

\begin{proof}
Put $P_g=(1+3/g_1)^f(1+3/g_2)^f$, $A=(1+\delta)e_1$, $\varepsilon'=2^{-(k-2r-5)}e^{C_n(2r+2)^5f\log b}$. By
\cref{pf:thmS:num}: $0\le P_g\le e^{6b^{-7}}$ by (1); $e_1\le\exp(12f2^{-(k-2r-2)})\le\exp(e^{-k/5})\le e^{b^{-14}}$ by (3)
and (2), so $A\le(1+b^{-7})e^{b^{-14}}$; $\varepsilon\le e^{-k/5}\le b^{-14}$ by (4), (2); and $\varepsilon'\le
2^{-k/2}e^{C_n(2r+2)^5f\log b}\le e^{-k/5}\le b^{-14}$ by (6) (as $4r+10\le k$), (5) with $m=2r+2$, and (2).
\Cref{pf:thmS:pba} and \cref{pf:thmS:num}(7) give the claim.
\end{proof}

\begin{theorem}[Pair bound from the main row]\label{pf:thmS:pbfromlocal}
Assume the main-row statement (\cref{pf:thmS:mainrowthm}, as an implication from its inputs: the local inputs, the global inputs, the lattice link, the local-model axioms and the theta step), and
the pair bound (\cref{pf:thmS:pairbound}) with its inputs (\cref{pf:thmS:pba,pf:thmS:num} and the lemmas used there).
Then for every $C_4\ge1$ and $C_{\mathrm{thr}}$ there are $C_A\ge\max(C_{\mathrm{thr}},1)$ and $0\le C\le C_A$ such that
for all $C_0\ge1$, $C_G\ge0$, $C_\Delta>4$, $c_T,C_b>0$ with $c_TC_b>1$ there is $k_0$ with the following property. Let
the data of the main row be given, satisfying the hypotheses of \cref{pf:thmS:mainrowthm} with these constants: a CM
field $E$ unramified over its real subfield, a TB-datum $D$ over $E$, a global size datum $G$ with $k_0\le G.k$ and the
size hypotheses, an instance $\mathrm{Inst}$, the local model, the real inputs, the density input, mass data $B$ with
$\mathrm{RI}$, $\mathrm{MI}$ and $\mathrm{En}$. Let $\mathcal X$ be a term system with $H_1$
complete, $c_5,C_n\in\R$, and assume:
\begin{itemize}
\item $\mathrm{LC}$: $\Gamma_1(\mathcal X)=\Gamma_1(B)$, $\thmSps(\mathrm{sc}(\thmSGen))=\Gamma_1(B)\Sigma(B)$, there is a
bijection between the index set of $B$ and $\{0,1\}\times\{\text{infinite places of }E\}$ under which $g_p=g_1$ or $g_2$
according to the first coordinate, $d(B)=d_u$, $d'(B)=d_{u'}$ and $D_*(B)=D_*(G.\mathrm{size}(d_e,T,d_u))$;
\item $\mathrm{SI}(\mathcal X,\thmStt(B))$: absolute convergence, the Gram expansions on $\thmSGen$ and
$\thmSNG$, $\mathrm{GB}(\thmStt(B))$, $\mathrm{RO}$, $\mathrm{CP}$, $\mathrm{IZ}$, and $\Gamma(i,i)=\Gamma_1$;
\item $\mathrm{SN}(G.k,G.r,G.f,b,G.g_1,G.g_2,c_5,C_n)$, $\mathrm{R}_{\mathrm{big}}$ and $\mathrm{R}_{\mathrm{nsc}}$ with
$(e^{-c_5G.kG.f\log G.k},dd)$, and $\mathrm{R}_{\mathrm{ng}}(2^{-(G.k-2G.r-5)}e^{C_n(2G.r+2)^5G.f\log b},dd)$, where
$dd=d_ud_{u'}/D_*(G.\mathrm{size}(d_e,T,d_u))$.
\end{itemize}
Then $\|K_0\|^2\le(1+b^{-6})\,dd$.
\formallabel{Formalism.FourierDecayAlt.P.GrowthPolynomialOperator.t_PairBoundFromLocal}
\end{theorem}

\begin{proof}
Take the constants $C_A,C$ and $k_0$ of the main row theorem \cref{pf:thmS:mainrowthm}. It gives
$\Gamma_1(B)\Sigma(B)\le(1+\delta)e_1\,d(B)d'(B)/D_*(B)$ with $0\le\delta\le b^{-7}$ and $1\le e_1\le
\exp(12G.f\,2^{-(G.k-2G.r-2)})$. By $\mathrm{LC}$ the left side is $\thmSps(\mathrm{sc}(\thmSGen))$ and the right
side is $(1+\delta)e_1\,dd$: this is $\mathrm{R}_{\mathrm{main}}((1+\delta)e_1,dd)$. The bijection of $\mathrm{LC}$ gives
$\thmStt(B)=\prod_p(1+3/g_p)=(1+3/g_1)^{G.f}(1+3/g_2)^{G.f}$, $G.f$ being the number of infinite places. Positivity of
the mass data gives $D_*(B)>0$, hence $dd\ge0$. Now apply \cref{pf:thmS:pairbound}.
\end{proof}

\begin{theorem}[The abstract Theorem S]\label{pf:thmS:pbapplied}
The conclusion of \cref{pf:thmS:pbfromlocal} holds unconditionally.
\formallabel{Formalism.L2.SG.TheoremS.pbApplied}
\end{theorem}

\begin{proof}
All hypotheses of \cref{pf:thmS:pbfromlocal} are theorems: the main row is \cref{pf:thmS:mainrowthm}; the pair bound
is \cref{pf:thmS:pairbound}, with \cref{pf:thmS:pba,pf:thmS:price,pf:thmS:gram,pf:thmS:pairsplit,pf:thmS:subfamily,pf:thmS:num};
and the inputs of the main row are the local and global inputs
\cref{pf:thmS:localcores,pf:thmS:globalcores},
the lattice link \lref{Formalism.FourierDecayAlt.P.GrowthPolynomialOperator.t_LatticeLink}, the local-model axioms
\lref{Formalism.FourierDecayAlt.P.GrowthPolynomialOperator.t_AxiomsOfLocalModel} and the theta step
\cref{pf:thmS:thetastep} (see \cref{sec:pf:thmS:main}).
\end{proof}

% ---------------------------------------------------------------------------
\subsection{The main row: Theorem A and the mass identities}\label{sec:pf:thmS:main}

The main row bounds the contribution of the scalar generic pairs. In the concrete situation this contribution is,
after the counting formula and the row forms, a multiple of a sum $\sum_\rho\|\sum_pc(\rho,p)V(p)\|^2$ of norms of
combinations of vectors $V(p)$ attached to the points $p$ of a tower of flags, with coefficients $c(\rho,p)$ dominated
by a positive semi-definite kernel. \Cref{sec:pf:thmS:tbdef,sec:pf:thmS:steps,sec:pf:thmS:thmA} prove a Bessel-type
inequality for such sums (Theorem~A); \cref{sec:pf:thmS:mass} converts it into the main row through identities
between volumes and masses.

\subsubsection{Tower data}\label{sec:pf:thmS:tbdef}

\begin{definition}[Points, towers, TB-data]\label{pf:thmS:tbdata}
Let $E$ be a field and $W$ a set (of \emph{words}). A \emph{point} is a finite sequence
$p=((s_1,t_1),\dots,(s_n,t_n))$ of pairs in $\N\times E$; its level is $\thmSlev(p)=\sum_js_j$, and
$\thmSjoin(h,s,t)$ is $h$ followed by $(s,t)$. A point $p$ is a point of level $i$ if $\thmSlev(p)=i$ and $s_j\ge1$,
$t_j\ne0$ for all $j$. A \emph{tower} $T$ consists of finite sets $\mathrm{Sg}(b,s)\subseteq E\smallsetminus\{0\}$
($b,s\in\N$) and finite sets $X_i$ of points, where $X_i$ is the set of $p$ with $\thmSlev(p)=i$, $s_j\ge1$ and
$t_j\in\mathrm{Sg}(s_1+\dots+s_{j-1}+1,s_j)$ for all $j$; thus $X_0=\{\emptyset\}$. A kernel $\kappa$ on points is
\emph{positive semi-definite (PSD) on} a finite set $Y$ if the matrix $(\kappa(x,y))_{x,y\in Y}$ is.

A \emph{TB-datum} $D$ over $(E,W)$ consists of integers $r,f,K$; functions $a_v:E\to\R$ ($1\le v\le f$); a weight
$\mathrm{wt}:W\to\R$ and a count $\mathrm{cnt}:W\to\N$; a finite set $W_{\mathrm{all}}\subseteq W$ containing an
element $\varnothing$ with $\mathrm{cnt}(\varnothing)=0$, $\mathrm{wt}(\varnothing)=1$; kernels
$\thmSgram_{\eps,i,w}(p,q)$ ($\eps\in\R$, $i\in\N$, $w\in W$), $K_i(p,q)$ with $K_0(\emptyset,\emptyset)=1$, and
$\Theta_{i,s,t}(h,h')$ ($t\in E$); reals $e_i$; additive subgroups $\mathrm{La}(b,s)$, $\mathrm{Lap}(b,s)$,
$\mathrm{Lad}(b)$ of $E$; complex numbers $\thmSNh(\eps,i,w,Y,\kappa,l)$ (for finite sets $Y$ of points, kernels
$\kappa$ and $l\in E$); and reals $\mathrm{szD}(i,v)$, $\mathrm{szL}(b,s,v)$, $\mathrm{vol}(b,s)$. Put
$N_{\mathrm{top}}=4K(r+1)$, $\mathrm{bud}(i)=4K(r+1-i)$, $W(n)=\{w\in W_{\mathrm{all}}:\mathrm{cnt}(w)\le n\}$, and
\begin{gather*}
\thmSSf(\eps,i,w,Y)=\Re\sum_{p,q\in Y}K_i(p,q)\,\thmSgram_{\eps,i,w}(p,q),\qquad Q(T,\eps,i,w)=\thmSSf(\eps,i,w,X_i)/\mathrm{wt}(w),\\
\thmSNst(\eps,i,s,w,Y,\kappa;t,t')=\mathrm{wt}(w)^{-1}\sum_{h,h'\in Y}\kappa(h,h')\,
\thmSgram_{\eps,i,w}(\thmSjoin(h,s,t),\thmSjoin(h',s,t')),\\
\thmSNz(\eps,j,w,Y,\kappa)=\mathrm{wt}(w)^{-1}\Re\sum_{h,h'\in Y}\kappa(h,h')\,\thmSgram_{\eps,j,w}(h,h'),\\
\thmSblk(\kappa,s,t,t')(h,h')=\kappa(\thmSjoin(h,s,t),\thmSjoin(h',s,t')),\qquad
\thmSmaj(\rho,l)=\prod_{v=1}^f\min\bigl(1,(\rho_v/a_v(l))^K\bigr),\\
\thmSDg(T,\eps,i,s,w)=\sum_{t\in\mathrm{Sg}(i-s+1,s)}\Re\thmSNst(\eps,i,s,w,X_{i-s},\thmSblk(K_i,s,t,t);t,t),
\end{gather*}
and let $X_{i,s}\subseteq X_i$ be the set of points whose last block has size $s$.
\end{definition}

\begin{definition}[Hypotheses on TB-data]\label{pf:thmS:tbhyp}
Let $D$ be a TB-datum, $C_4,\mathrm{ep},\mathrm{jn},\Theta\in\R$, $V$ a map from points to a Hilbert space $H$ and
$c:\N\times\{\text{points}\}\to\C$. In the following, $1\le s\le i\le r+1$, ``$Y$ of level $j$'' means a finite set of
points of level $j$, and $\eps>0$.
\begin{itemize}
\item (Positivity.) $\mathrm{GP}$: every $\thmSgram_{\eps,i,w}$ is PSD on every finite set; $\mathrm{KP}$: every $K_i$ is; the
conjunction of $\mathrm{KP}$ with the same property of all $\Theta_{i,s,t}$ is $\mathrm{KP}^{+}$;
$\mathrm{WP}$: $\mathrm{wt}>0$; $\mathrm{E}$: $0\le e_i\le1$ for $1\le i\le r+1$.
\item (Support.) $\mathrm{KS}$: let $h,h'$ be points of level $i-s$ and let $t,t'\ne0$ be distinct. If the value
$K_i(\thmSjoin(h,s,t),\thmSjoin(h',s,t'))$ is non-zero, then $(i,s)\ne(r+1,1)$ and $t^{-1}-t'^{-1}$ lies in
$\mathrm{Lap}(i-s+1,s)$.
\item (Theta domination.) $\mathrm{TD}$: $\mathrm{E}$; $\Theta_{i,s,t}-\thmSblk(K_i,s,t,t)$ is PSD on every $Y$ of level
$i-s$ for $t\ne0$; $\Theta_{i,s,t}$ depends only on $t$ modulo $\mathrm{La}(i-s+1,s)$; and for every finite $S\subseteq E$
of elements pairwise incongruent modulo $\mathrm{La}(i-s+1,s)$, $(1+e_i)K_{i-s}-\sum_{t\in S}\Theta_{i,s,t}$ is PSD on every
$Y$ of level $i-s$.
\item (Domination on towers.) $\mathrm{DT}(T,\Theta')$, for a tower $T$ and kernels $\Theta'$: the three domination properties of
$\mathrm{TD}$ for $\Theta'$, restricted to $t,t'\in\mathrm{Sg}(i-s+1,s)$, $S\subseteq\mathrm{Sg}(i-s+1,s)$ and
$Y=X_{i-s}$. $\mathrm{DT}_{\forall}$: for every tower $T$ there is $\Theta'$ with $\mathrm{DT}(T,\Theta')$.
\item (Wave estimate.) $\mathrm{EW}$: for $w\in W_{\mathrm{all}}$ with $\mathrm{cnt}(w)+K\le N_{\mathrm{top}}$ there are probability
vectors $\theta_A$ on $W(\mathrm{cnt}(w)+K)$, $A\subseteq\{1,\dots,f\}$, such that for $Y$ of level $i-1$ and $\kappa$
PSD on $Y$: $\thmSNh(\eps,i,w,Y,\kappa,0)=\thmSNz(\eps,i-1,w,Y,\kappa)$, and for $0\ne l\in\mathrm{Lad}(i)$ there is $A$
with $|\thmSNh(\eps,i,w,Y,\kappa,l)|\le\thmSmaj(\mathrm{szD}(i),l)\sum_{w'}\theta_A(w')\thmSNz(\eps,i-1,w',Y,\kappa)$.
\item (Poisson step.) $\mathrm{PS}$: for $w\in W$, $Y$ of level $i-1$, $\kappa$ PSD on $Y$, $t_0\in E$ and $B\in\R$: if
$\sum_{l\in L}|\thmSNh(\eps,i,w,Y,\kappa,l)|\le B$ for every finite $L\subseteq\mathrm{Lad}(i)\smallsetminus\{0\}$, then
$\sum_{t\in S}\Re\thmSNst(\eps,i,1,w,Y,\kappa;t,t)\le|\thmSNh(\eps,i,w,Y,\kappa,0)|+B$ for every finite
$S\subseteq(t_0+\mathrm{La}(i,1))\smallsetminus\{0\}$.
\item (Lift estimate.) $\mathrm{EL}$: for $(i,s)\ne(r+1,1)$ and $w$ with $\mathrm{cnt}(w)+K\le N_{\mathrm{top}}$ there are
probability vectors $\theta_\pi$ on $W(\mathrm{cnt}(w)+K)^2$, $\pi\in\{1,2,3\}^f$, such that for $Y$ of level $i-s$,
finite $S\subseteq E$, $\kappa$ PSD on $\{\thmSjoin(h,s,t):h\in Y,t\in S\}$ and $t,t'\in S\smallsetminus\{0\}$ with
$t^{-1}\ne t'^{-1}$ there is $\pi$ with
\[
|\thmSNst(\eps,i,s,w,Y,\thmSblk(\kappa,s,t,t');t,t')|\le\thmSmaj(\mathrm{szL}(i-s+1,s),t^{-1}-t'^{-1})
\sum_{(x,x')}\theta_\pi(x,x')\,\frac{a_x(t)+a_{x'}(t')}2,
\]
where $a_x(t)=\Re\thmSNst(\eps,i,s,x,Y,\thmSblk(\kappa,s,t,t);t,t)$.
\item (Block estimate.) $\mathrm{EB}(C_4)$: for $s\ge2$, $w$ with $\mathrm{cnt}(w)+(s+1)K\le N_{\mathrm{top}}$ and
$\mathrm{vol}(i-s+1,s)\le1$ there is a probability vector $\theta$ on $W(\mathrm{cnt}(w)+(s+1)K)$ such that for $Y$ of
level $i-s$, $\kappa$ PSD on $Y$ and finite $S\subseteq(t_0+\mathrm{La}(i-s+1,s))\smallsetminus\{0\}$,
$\sum_{t\in S}\Re\thmSNst(\eps,i,s,w,Y,\kappa;t,t)\le C_4^{fs}\prod_{1\le s'<s}\mathrm{vol}(i-s+1,s')
\sum_{w'}\theta(w')\thmSNz(\eps,i-s,w',Y,\kappa)$.
\item (Lattice sums.) $\mathrm{LS}(\mathrm{ep})$: $\sum_{l\in L}\thmSmaj(\mathrm{szD}(i),l)\le\mathrm{ep}$ for finite
$L\subseteq\mathrm{Lad}(i)\smallsetminus\{0\}$, and $\sum_{l\in L}\thmSmaj(\mathrm{szL}(i-s+1,s),l)\le\mathrm{ep}$ for finite
$L\subseteq\mathrm{Lap}(i-s+1,s)\smallsetminus\{0\}$ when $(i,s)\ne(r+1,1)$.
\item (Energy.) $\mathrm{EN}(\mathrm{jn})$: $\Re\thmSgram_{\eps,0,w}(\emptyset,\emptyset)/\mathrm{wt}(w)\to\mathrm{jn}$ as
$\eps\to0^+$, for $w\in W(N_{\mathrm{top}})$.
\item (Cutoff.) $\mathrm{CO}(V)$: for every finite set $Y$ of points of level $r+1$, for all sufficiently small $\eps>0$,
$\thmSgram_{\eps,r+1,\varnothing}(p,q)=\langle V(q),V(p)\rangle$ whenever $p,q\in Y$ and $K_{r+1}(p,q)\ne0$.
\item (Leading domination.) $\mathrm{LD}(V,c)$: $c(\rho,p)\ne0$ only for points $p$ of level $r+1$; for every finite set $Y$
of such points and finite $R\subseteq\N$, $K_{r+1}-\sum_{\rho\in R}c(\rho,\cdot)\overline{c(\rho,\cdot)}$ is PSD on $Y$;
and $\sum_p|c(\rho,p)|\,\|V(p)\|<\infty$ for every $\rho$.
\item (Sizes.) $\mathrm{SB}(\Theta)$: $\mathrm{vol}\ge0$, and $\mathrm{vol}(b,s)\le\Theta^{-2f}$ for $b,s\ge1$, $b+s\le r+2$.
\end{itemize}
\end{definition}

We use two elementary facts throughout: by the Schur product theorem \cite[Thm.~7.5.3]{HornJohnson2013}, the entrywise
product of two PSD matrices is PSD, so the sum of its entries is $\ge0$; consequently, under $\mathrm{GP}$ and
$\mathrm{WP}$, $\kappa\mapsto\thmSNz(\eps,j,w,Y,\kappa)$ is linear, $\ge0$ on kernels PSD on $Y$, and monotone for
the PSD order on $Y$, and the same holds for $\kappa\mapsto\Re\thmSNst(\eps,i,s,w,Y,\kappa;t,t)$ (the kernel
$(h,h')\mapsto\thmSgram_{\eps,i,w}(\thmSjoin(h,s,t),\thmSjoin(h',s,t))$ being PSD).

\subsubsection{The one-level inequality}\label{sec:pf:thmS:steps}

Fix a TB-datum $D$ with $\mathrm{GP}$, $\mathrm{KP}^{+}$ and $\mathrm{WP}$, a tower $T$ and $\eps>0$; let
$1\le s\le i\le r+1$ and $w\in W$.

\begin{lemma}[Five steps]\label{pf:thmS:fivesteps}
\begin{enumerate}
\item (Split by size.) $\sqrt{\thmSSf(\eps,i,w,X_i)}\le\sum_{s=1}^i\sqrt{\thmSSf(\eps,i,w,X_{i,s})}$.
\item (Schur step.) Assume $\mathrm{KS}$, $\mathrm{EL}$ and $\mathrm{LS}(\mathrm{ep})$. If
$w\in W_{\mathrm{all}}$, $\mathrm{cnt}(w)+K\le N_{\mathrm{top}}$ and $\thmSDg(T,\eps,i,s,w')\le M$ for all
$w'\in W(\mathrm{cnt}(w)+K)$, then $\thmSSf(\eps,i,w,X_{i,s})/\mathrm{wt}(w)\le\thmSDg(T,\eps,i,s,w)+3^f\mathrm{ep}\,M$.
\item (Diagonal step.) Assume $\mathrm{TD}$. If $R\subseteq E$ is finite and every
$t\in\mathrm{Sg}(i-s+1,s)$ is congruent modulo $\mathrm{La}(i-s+1,s)$ to some $t_0\in R$, then
$\thmSDg(T,\eps,i,s,w)\le\sum_{t_0\in R}\sum_{t}\Re\thmSNst(\eps,i,s,w,X_{i-s},\Theta_{i,s,t_0};t,t)$, the inner sum over the
$t\in\mathrm{Sg}(i-s+1,s)$ congruent to $t_0$.
\item (One step.) Assume $\mathrm{TD}$, $\mathrm{EW}$, $\mathrm{PS}$ and
$\mathrm{LS}(\mathrm{ep})$. If $w\in W_{\mathrm{all}}$, $\mathrm{cnt}(w)+K\le N_{\mathrm{top}}$, $R$ consists of
elements pairwise incongruent modulo $\mathrm{La}(i,1)$, and $\thmSNz(\eps,i-1,w',X_{i-1},K_{i-1})\le M$ for all
$w'\in W(\mathrm{cnt}(w)+K)$, then
\[
\sum_{t_0\in R}\sum_{t\equiv t_0}\Re\thmSNst(\eps,i,1,w,X_{i-1},\Theta_{i,1,t_0};t,t)\le(1+e_i)\bigl(\thmSNz(\eps,i-1,w,X_{i-1},K_{i-1})
+2^f\mathrm{ep}\,M\bigr).
\]
\item (Multi step.) Assume $\mathrm{TD}$, $\mathrm{EB}(C_4)$, $C_4\ge1$ and $\mathrm{vol}\ge0$. If
$s\ge2$, $w\in W_{\mathrm{all}}$, $\mathrm{cnt}(w)+(s+1)K\le N_{\mathrm{top}}$, $\mathrm{vol}(i-s+1,s)\le1$, $R$ consists of
elements pairwise incongruent modulo $\mathrm{La}(i-s+1,s)$, and $\thmSNz(\eps,i-s,w',X_{i-s},K_{i-s})\le M$ for all
$w'\in W(\mathrm{cnt}(w)+(s+1)K)$, then
\[
\sum_{t_0\in R}\sum_{t\equiv t_0}\Re\thmSNst(\eps,i,s,w,X_{i-s},\Theta_{i,s,t_0};t,t)\le C_4^{fs}\prod_{1\le s'<s}
\mathrm{vol}(i-s+1,s')\,(1+e_i)M.
\]
\end{enumerate}
\formallabel{Formalism.FourierDecayAlt.P.splitByLastSize}
\formallabel{Formalism.FourierDecayAlt.P.schurRecursion_true}
\formallabel{Formalism.FourierDecayAlt.P.diagStep}
\formallabel{Formalism.FourierDecayAlt.P.singleBlockStep}
\formallabel{Formalism.FourierDecayAlt.P.longBlockStep}
\end{lemma}

\begin{proof}
(1) The kernel $K_i\cdot\thmSgram_{\eps,i,w}$ (entrywise product) is PSD on $X_i$ by the Schur product theorem, so
$q(x)=\Re\sum_{p,q}x_p\overline{x_q}K_i(p,q)\thmSgram_{\eps,i,w}(p,q)$ is a positive semi-definite quadratic form and
$\sqrt q$ is a seminorm. Since $i\ge1$, $X_i$ is the disjoint union of the $X_{i,s}$, $1\le s\le i$; apply the triangle
inequality for $\sqrt q$ to the decomposition of the indicator vector of $X_i$ into those of the $X_{i,s}$.

(2) Put $Y=X_{i-s}$ and $S=\mathrm{Sg}(i-s+1,s)$. The map $(h,t)\mapsto\thmSjoin(h,s,t)$ is a bijection from $Y\times S$
onto $X_{i,s}$, so $\thmSSf(\eps,i,w,X_{i,s})/\mathrm{wt}(w)=\sum_{t,t'\in S}\Re\thmSNst(\eps,i,s,w,Y,\thmSblk(K_i,s,t,t');t,t')$.
The diagonal terms sum to $\thmSDg(T,\eps,i,s,w)$. For $t\ne t'$, $\mathrm{KS}$ shows that the term vanishes
unless $(i,s)\ne(r+1,1)$ and $0\ne t^{-1}-t'^{-1}\in\mathrm{Lap}(i-s+1,s)$; in that case $\mathrm{EL}$ (with
$\kappa=K_i$, PSD on the points) bounds it by $m(t,t')\sum_\pi\sum_{(x,x')}\theta_\pi(x,x')(a_x(t)+a_{x'}(t'))/2$, where
$m(t,t')=\thmSmaj(\mathrm{szL}(i-s+1,s),t^{-1}-t'^{-1})$ and the sum over all $3^f$ maps $\pi$ dominates the term for
the $\pi$ given by $\mathrm{EL}$, all terms being $\ge0$. Since $t'\mapsto t^{-1}-t'^{-1}$ is injective,
$\mathrm{LS}(\mathrm{ep})$ gives $\sum_{t'}m(t,t')\le\mathrm{ep}$ and $\sum_tm(t,t')\le\mathrm{ep}$ (sums over the
admissible pairs). Hence, for each $\pi$ and $(x,x')$,
$\sum_{t\ne t'}m(t,t')\frac{a_x(t)+a_{x'}(t')}2\le\frac{\mathrm{ep}}2\bigl(\sum_ta_x(t)+\sum_{t'}a_{x'}(t')\bigr)\le\mathrm{ep}\,M$,
as $\sum_ta_x(t)=\thmSDg(T,\eps,i,s,x)\le M$ (Schur's test). Summing against the probability vectors $\theta_\pi$ and over
the $3^f$ maps $\pi$ gives the bound $3^f\mathrm{ep}\,M$ for the off-diagonal part.

(3) For $t\in\mathrm{Sg}(i-s+1,s)$ choose $t_0\in R$ with $t\equiv t_0$. By $\mathrm{TD}$,
$\Theta_{i,s,t}-\thmSblk(K_i,s,t,t)$ is PSD on $X_{i-s}$, so $\Re\thmSNst(\dots,\thmSblk(K_i,s,t,t);t,t)\le
\Re\thmSNst(\dots,\Theta_{i,s,t};t,t)=\Re\thmSNst(\dots,\Theta_{i,s,t_0};t,t)$ by monotonicity and periodicity; the right
side of (3) contains this term, and its other terms are $\ge0$.

(4) Fix $t_0\in R$ and put $\kappa=\Theta_{i,1,t_0}$, PSD on $X_{i-1}$. $\mathrm{EW}$ gives probability vectors
$\theta_A$ with $\thmSNh(\dots,\kappa,0)=\thmSNz(\eps,i-1,w,X_{i-1},\kappa)\ge0$ and
$|\thmSNh(\dots,\kappa,l)|\le\thmSmaj(\mathrm{szD}(i),l)\sum_Ag_A(t_0)$ for $0\ne l\in\mathrm{Lad}(i)$, where
$g_A(t_0)=\sum_{w'}\theta_A(w')\thmSNz(\eps,i-1,w',X_{i-1},\Theta_{i,1,t_0})\ge0$. By $\mathrm{LS}$,
$\sum_{l\in L}|\thmSNh(\dots,l)|\le\mathrm{ep}\sum_Ag_A(t_0)$ for all finite $L$, and $\mathrm{PS}$ gives
$\sum_{t\equiv t_0}\Re\thmSNst(\eps,i,1,w,X_{i-1},\kappa;t,t)\le\thmSNz(\eps,i-1,w,X_{i-1},\Theta_{i,1,t_0})+\mathrm{ep}\sum_Ag_A(t_0)$.
Now sum over $t_0\in R$: since $(1+e_i)K_{i-1}-\sum_{t_0\in R}\Theta_{i,1,t_0}$ is PSD on $X_{i-1}$, linearity and
positivity of $\thmSNz$ give $\sum_{t_0}\thmSNz(\eps,i-1,w',X_{i-1},\Theta_{i,1,t_0})\le(1+e_i)\thmSNz(\eps,i-1,w',X_{i-1},
K_{i-1})$ for every $w'$; for $w'=w$ this bounds the first part, and averaging against $\theta_A$ gives
$\sum_{t_0}g_A(t_0)\le(1+e_i)M$ for each of the $2^f$ subsets $A$.

(5) For each $t_0\in R$, $\mathrm{EB}(C_4)$ with $\kappa=\Theta_{i,s,t_0}$ bounds the inner sum by
$C_4^{fs}\prod\mathrm{vol}\cdot\sum_{w'}\theta(w')\thmSNz(\eps,i-s,w',X_{i-s},\Theta_{i,s,t_0})$; sum over $t_0$ and use, as in
(4), that $\sum_{t_0}\thmSNz(\dots,\Theta_{i,s,t_0})\le(1+e_i)\thmSNz(\dots,K_{i-s})\le(1+e_i)M$; the prefactor is $\ge0$.
\end{proof}

\begin{proposition}[One-level inequality]\label{pf:thmS:onelevel}
Let $D$ satisfy $\mathrm{GP}$, $\mathrm{KP}^{+}$, $\mathrm{WP}$, $\mathrm{KS}$,
$\mathrm{TD}$, $\mathrm{EW}$, $\mathrm{PS}$, $\mathrm{EL}$, $\mathrm{EB}(C_4)$ and
$\mathrm{LS}(\mathrm{ep})$, with $\mathrm{ep}\ge0$, $C_4\ge1$, $\mathrm{vol}\ge0$ and $\mathrm{vol}(b,s)\le1$ for
$b,s\ge1$, $b+s\le r+2$. Let $T$ be a tower, $\eps>0$, $1\le i\le r+1$, and $\sigma_k\ge0$ ($k<i$) with
$Q(T,\eps,k,w)\le\sigma_k^2$ for all $w\in W(\mathrm{bud}(k))$. Then for every $w\in W(\mathrm{bud}(i))$,
\begin{multline}\label{eq:pf:thmS:onelevel}
\sqrt{Q(T,\eps,i,w)}\le\sqrt{(1+3^f\mathrm{ep})(1+2^f\mathrm{ep})}\sqrt{1+e_i}\,\sigma_{i-1}\\
+\sum_{s=2}^i\sqrt{(1+3^f\mathrm{ep})(1+e_i)C_4^{fs}\textstyle\prod_{1\le s'<s}\mathrm{vol}(i-s+1,s')}\;\sigma_{i-s}.
\end{multline}
\formallabel{Formalism.FourierDecayAlt.P.OneLevelCoefficients.oneLevel_of_steps}
\formallabel{Formalism.FourierDecayAlt.P.one_level_inequality_split_by_size}
\end{proposition}

\begin{proof}
Let $w\in W(\mathrm{bud}(i))$, so $\mathrm{cnt}(w)\le4K(r+1-i)$. For $1\le s\le i$ one has
$\mathrm{cnt}(w)+(s+2)K\le\min(N_{\mathrm{top}},\mathrm{bud}(i-s))$, since $s+2\le4s$ and $s\le i$; hence every $w''$
reached below lies in $W(\mathrm{bud}(i-s))$, where $Q(T,\eps,i-s,w'')\le\sigma_{i-s}^2$. Note
$\thmSNz(\eps,j,w'',X_j,K_j)=Q(T,\eps,j,w'')$. Choose for each $s$ a set $R\subseteq\mathrm{Sg}(i-s+1,s)$ of
representatives of the classes modulo $\mathrm{La}(i-s+1,s)$ met by $\mathrm{Sg}(i-s+1,s)$.

For $s=1$ and $w'\in W(\mathrm{cnt}(w)+K)$, \cref{pf:thmS:fivesteps}(3),(4) with $M=\sigma_{i-1}^2$ give
$\thmSDg(T,\eps,i,1,w')\le(1+e_i)(1+2^f\mathrm{ep})\sigma_{i-1}^2$; then (2) gives
$\thmSSf(\eps,i,w,X_{i,1})/\mathrm{wt}(w)\le(1+3^f\mathrm{ep})(1+e_i)(1+2^f\mathrm{ep})\sigma_{i-1}^2$. For $s\ge2$,
\cref{pf:thmS:fivesteps}(3),(5) with $M=\sigma_{i-s}^2$ (and $\mathrm{vol}(i-s+1,s)\le1$) and then (2) give
$\thmSSf(\eps,i,w,X_{i,s})/\mathrm{wt}(w)\le(1+3^f\mathrm{ep})C_4^{fs}\prod_{s'<s}\mathrm{vol}(i-s+1,s')(1+e_i)\sigma_{i-s}^2$.
By (1), divided by $\sqrt{\mathrm{wt}(w)}$, $\sqrt{Q(T,\eps,i,w)}\le\sum_{s=1}^i\sqrt{\thmSSf(\eps,i,w,X_{i,s})/\mathrm{wt}(w)}$,
and \eqref{eq:pf:thmS:onelevel} follows.
\end{proof}

\begin{corollary}[One-level inequality on towers]\label{pf:thmS:onelevelafter}
\Cref{pf:thmS:onelevel} holds with $\mathrm{KP}^{+}$ and $\mathrm{TD}$ replaced by $\mathrm{KP}$,
$\mathrm{E}$ and the existence of kernels $\Theta'$ with $\mathrm{DT}(T,\Theta')$.
\formallabel{Formalism.FourierDecayAlt.P.RecursionWithDensities.B1.oneLevelAfterOfOld}
\formallabel{Formalism.FourierDecayAlt.P.RecursionWithDensities.B1.oneLevelAfter}
\end{corollary}

\begin{proof}
Let $D^T$ be $D$ with $K_j$ replaced by its restriction to $X_j\times X_j$ (zero elsewhere) and $\Theta$ replaced as
follows: for $1\le s\le i\le r+1$ and $t$ whose class modulo $\mathrm{La}(i-s+1,s)$ meets $\mathrm{Sg}(i-s+1,s)$,
$\Theta^T_{i,s,t}$ is the restriction to $X_{i-s}\times X_{i-s}$ of $\Theta'_{i,s,t^\sharp}$, $t^\sharp$ a chosen element of
$\mathrm{Sg}(i-s+1,s)$ in that class; otherwise $\Theta^T_{i,s,t}=0$. Restrictions of PSD kernels are PSD, and
$\Theta'_{i,s,t}$ is PSD on $X_{i-s}$ since $\Theta'_{i,s,t}-\thmSblk(K_i,s,t,t)$ and $\thmSblk(K_i,s,t,t)$ are; so $D^T$
satisfies $\mathrm{KP}^{+}$. Zeros of $K_i$ remain zeros, so $\mathrm{KS}$ persists. $\mathrm{DT}$,
$\mathrm{KP}$ and $\mathrm{E}$ give $\mathrm{TD}$ for $D^T$ (on points outside the tower the kernels
vanish, blocks $\thmSblk(K^T_i,s,t,t)$ with $t\notin\mathrm{Sg}(i-s+1,s)$ vanish, and a finite set $S$ of incongruent
elements contributes only through its classes meeting the tower, each represented once). The other hypotheses do not involve $K$ or $\Theta$ except through PSD kernels supplied as
arguments, and $Q^{D^T}(T,\eps,k,w)=Q^D(T,\eps,k,w)$. Apply \cref{pf:thmS:onelevel} to $D^T$.
\end{proof}

\subsubsection{Theorem A}\label{sec:pf:thmS:thmA}

\begin{lemma}[Recursive inequalities]\label{pf:thmS:rect}
Let $r\in\N$, $0\le a\le\frac12$, $0\le\eps\le a$ with $ra\le1$, and let $e_i,\sigma_i\ge0$ satisfy, for $1\le i\le r+1$,
\begin{equation}\label{eq:pf:thmS:rec}
\sigma_i\le(1+\eps)\sqrt{1+e_i}\,\sigma_{i-1}+\sum_{t=2}^ia^{t-1}\sigma_{i-t}.
\end{equation}
Then $\sigma_{r+1}^2\le\bigl(1+e^9(r+1)a\bigr)\prod_{i=1}^{r+1}(1+e_i)\,\sigma_0^2$.
\formallabel{Formalism.FourierDecayAlt.P.RecText.recT}
\formallabel{Formalism.FourierDecayAlt.P.coarseBound}
\end{lemma}

\begin{proof}
Put $\lambda=1+\eps+2a$ and $Q_i=\prod_{j=1}^i\sqrt{1+e_j}$, which is $\ge1$ and non-decreasing. We show
$\sigma_i\le\lambda^iQ_i\sigma_0$ for $i\le r+1$ by strong induction: with $M=\lambda^{i-1}Q_i\sigma_0$, the induction
hypothesis gives $\sqrt{1+e_i}\,\sigma_{i-1}\le M$ and $\sigma_{i-t}\le M$ for $t\ge2$, and $\sum_{t=2}^ia^{t-1}\le2a$
(geometric series, $a\le\frac12$), so $\sigma_i\le(1+\eps+2a)M=\lambda^iQ_i\sigma_0$. Hence $\sigma_{r+1}^2\le
\lambda^{2(r+1)}\prod_{i\le r+1}(1+e_i)\sigma_0^2$. With $x=(r+1)a\le\frac32$ we have
$\lambda^{2(r+1)}\le e^{2(r+1)(\eps+2a)}\le e^{6x}$, and by convexity of $\exp$ on $[0,9]$,
$e^{6x}=e^{(2x/3)\cdot9}\le1-\frac{2x}3+\frac{2x}3e^9\le1+e^9x$.
\end{proof}

\begin{corollary}[Recursion with parameters]\label{pf:thmS:jointA}
For every $C_0\ge1$ there are $C'$ and $C$ (namely $C'=2C_0$, $C=e^9$) such that for all $K_x\ge4$, $f\ge1$,
$\Theta\ge C'$, $r$ with $r(C_0/\Theta)^f\le1$, and $e_i,\sigma_i\ge0$ satisfying \eqref{eq:pf:thmS:rec} with
$\eps=(C_0\Theta^{-(K_x-5/2)})^f$ and $a=(C_0/\Theta)^f$, one has
$\sigma_{r+1}^2\le(1+C(r+1)(C_0/\Theta)^f)\prod_{i=1}^{r+1}(1+e_i)\sigma_0^2$.
\formallabel{Formalism.FourierDecayAlt.P.K.jointTheoremA}
\end{corollary}

\begin{proof}
$\Theta\ge2C_0\ge2$ gives $a\le(C_0/\Theta)\le\frac12$ (as $f\ge1$), and $\Theta^{-(K_x-5/2)}\le\Theta^{-1}$ (as
$K_x\ge4$, $\Theta\ge1$) gives $\eps\le a$. Apply \cref{pf:thmS:rect}.
\end{proof}

\begin{lemma}[Limit step]\label{pf:thmS:limit}
Let $D$ satisfy $\mathrm{GP}$, $\mathrm{KP}^{+}$, $\mathrm{EN}(\mathrm{jn})$, $\mathrm{CO}(V)$ and
$\mathrm{LD}(V,c)$, with $H$ complete, and let $B\in\R$. Assume that for every tower $T$, every $\eps>0$
and every $\beta$ with $\Re\thmSgram_{\eps,0,w}(\emptyset,\emptyset)/\mathrm{wt}(w)\le\beta$ for all
$w\in W(N_{\mathrm{top}})$, one has $\thmSSf(\eps,r+1,\varnothing,X_{r+1})\le B\beta$. Then
$\sum_{\rho\in R}\|\sum_pc(\rho,p)V(p)\|^2\le B\,\mathrm{jn}$ for every finite $R\subseteq\N$.
\formallabel{Formalism.FourierDecayAlt.P.limitStep}
\end{lemma}

\begin{proof}
The series $\sum_pc(\rho,p)V(p)$ converge absolutely, so it suffices to show that for every finite set $F$ of points
there is a finite $F'\supseteq F$ with $\sum_{\rho\in R}\|\sum_{p\in F'}c(\rho,p)V(p)\|^2\le B\,\mathrm{jn}$. Let $F_0$
be the set of points of level $r+1$ in $F$, $S\subseteq E\smallsetminus\{0\}$ the set of all coordinates $t_j$ of
points of $F_0$, and $T$ the tower with $\mathrm{Sg}(b,s)=S$ for all $b,s$; put $Y=X_{r+1}\supseteq F_0$ and
$F'=F\cup Y$. Since $c(\rho,p)=0$ unless $p$ has level $r+1$, $\sum_{p\in F'}c(\rho,p)V(p)=\sum_{p\in Y}c(\rho,p)V(p)$.
Now
\begin{align*}
\sum_{\rho\in R}\Bigl\|\sum_{p\in Y}c(\rho,p)V(p)\Bigr\|^2&=\Re\sum_{p,q\in Y}\Bigl(\sum_{\rho\in R}c(\rho,p)\overline{c(\rho,q)}\Bigr)
\langle V(q),V(p)\rangle\\
&\le\Re\sum_{p,q\in Y}K_{r+1}(p,q)\langle V(q),V(p)\rangle,
\end{align*}
because $K_{r+1}-\sum_\rho c(\rho,\cdot)\overline{c(\rho,\cdot)}$ and the Gram matrix of the $V(p)$ are PSD on $Y$ (Schur
product theorem). For $\delta>0$ choose $\eps>0$ so small that $\Re\thmSgram_{\eps,0,w}(\emptyset,\emptyset)/\mathrm{wt}(w)
\le\mathrm{jn}+\delta$ for the finitely many $w\in W(N_{\mathrm{top}})$ and the cutoff identity holds on $Y$; then the
right side equals $\thmSSf(\eps,r+1,\varnothing,Y)\le B(\mathrm{jn}+\delta)$. Let $\delta\to0$.
\end{proof}

\begin{lemma}[Coefficient bound]\label{pf:thmS:coef}
Let $D$ be a TB-datum, $\Theta\ge1$, $C_4\ge1$, $C_0\ge2C_4$, $f\ge1$, $\mathrm{SB}(\Theta)$, $0\le a\le1$, $0\le e\le1$ and
$2\le s\le i\le r+1$. Then
\[
\sqrt{(1+a)(1+e)\,C_4^{fs}\textstyle\prod_{1\le s'<s}\mathrm{vol}(i-s+1,s')}\le\bigl((C_0/\Theta)^f\bigr)^{s-1}.
\]
\formallabel{Formalism.FourierDecayAlt.P.OneLevelCoefficients.coef}
\end{lemma}

\begin{proof}
$(1+a)(1+e)\le4$; $\prod_{s'<s}\mathrm{vol}(i-s+1,s')\le\Theta^{-2f(s-1)}$ by $\mathrm{SB}(\Theta)$ (the indices satisfy
$b+s'\le r+2$); and $4C_4^{fs}\le(2C_4)^{2f(s-1)}\le C_0^{2f(s-1)}$, since $f(s-1)\ge1$, $s\le2(s-1)$ and $C_4\ge1$. Hence the
radicand is at most $(C_0^{f(s-1)}\Theta^{-f(s-1)})^2$.
\end{proof}

\begin{theorem}[Theorem A]\label{pf:thmS:thmA}
For all $C_1,C_4\ge1$ there are $C,C'$ such that the following holds. Let $D$ be a TB-datum, $\Theta,\mathrm{jn}\in\R$,
$H$ a Hilbert space, $V$ and $c$ as above, with $K\ge4$, $f\ge1$, $0<\Theta$, $C'\le\Theta$, $r(C/\Theta)^f\le1$,
$\mathrm{SB}(\Theta)$, $\mathrm{GP}$, $\mathrm{KP}$, $\mathrm{WP}$, $\mathrm{KS}$,
$\mathrm{E}$, $\mathrm{DT}_{\forall}$, $\mathrm{EW}$, $\mathrm{PS}$, $\mathrm{EL}$,
$\mathrm{EB}(C_4)$, $\mathrm{LS}((C_1\Theta^{-(K-5/2)})^f)$, $\mathrm{EN}(\mathrm{jn})$,
$\mathrm{CO}(V)$ and $\mathrm{LD}(V,c)$. Then for every finite $R\subseteq\N$,
\begin{equation}\label{eq:pf:thmS:thmA}
\sum_{\rho\in R}\Bigl\|\sum_pc(\rho,p)V(p)\Bigr\|^2\le\Bigl(1+C(r+1)\bigl(C/\Theta\bigr)^f\Bigr)\prod_{i=1}^{r+1}(1+e_i)\;\mathrm{jn}.
\end{equation}
\formallabel{Formalism.FourierDecayAlt.P.RecursionWithDensities.B2.theoremAAfterOfSteps}
\formallabel{Formalism.FourierDecayAlt.P.theorem_a_after}
\formallabel{Formalism.FourierDecayAlt.Final.tc_theoremAAfter}
\end{theorem}

\begin{proof}
Let $C_0=\max(3C_1,2C_4)$, let $(C'_J,C_J)$ be the constants of \cref{pf:thmS:jointA} for $C_0$, and put
$C=\max(C_0,C_J)$, $C'=\max(C'_J,3C_1,1)$. Then $\Theta\ge1$, $\Theta\ge3C_1$ and $\Theta\ge C'_J$. Put
$y=\Theta^{-(K-5/2)}\le\Theta^{-1}$ and $\mathrm{ep}=(C_1y)^f$; then $3^f\mathrm{ep}=(3C_1y)^f\le\min(1,(C_0y)^f)$ and
$2^f\mathrm{ep}\le3^f\mathrm{ep}$. Replacing all $\Theta_{i,s,t}$ by $0$ does not affect the hypotheses used and gives
$\mathrm{KP}^{+}$; by \cref{pf:thmS:limit} it suffices to show, for every tower $T$, $\eps>0$ and $\beta$ as there,
that $\thmSSf(\eps,r+1,\varnothing,X_{r+1})\le B\beta$ with $B=(1+C(r+1)(C/\Theta)^f)\prod_i(1+e_i)$.

Let $\sigma_k=\bigl(\max_{w\in W(\mathrm{bud}(k))}\max(Q(T,\eps,k,w),0)\bigr)^{1/2}$. Since $X_0=\{\emptyset\}$ and
$K_0(\emptyset,\emptyset)=1$, $Q(T,\eps,0,w)=\Re\thmSgram_{\eps,0,w}(\emptyset,\emptyset)/\mathrm{wt}(w)$, and
$\mathrm{bud}(0)=N_{\mathrm{top}}$; so $\sigma_0^2\le\beta$, and $\beta\ge0$ by $\mathrm{GP}$. By
$\mathrm{SB}(\Theta)$ and $\Theta\ge1$, $0\le\mathrm{vol}(b,s)\le1$ in the relevant range. By $\mathrm{DT}_{\forall}$
there are $\Theta'$ with $\mathrm{DT}(T,\Theta')$, so \cref{pf:thmS:onelevelafter} applies and gives
\eqref{eq:pf:thmS:onelevel} for every $w\in W(\mathrm{bud}(i))$, hence for $\sigma_i$. Its coefficients are bounded as
follows: $\sqrt{(1+3^f\mathrm{ep})(1+2^f\mathrm{ep})}\le1+3^f\mathrm{ep}\le1+(C_0y)^f$, and for $s\ge2$ the
coefficient of $\sigma_{i-s}$ is at most $((C_0/\Theta)^f)^{s-1}$ by \cref{pf:thmS:coef} (with $a=3^f\mathrm{ep}\le1$ and
$e=e_i\le1$). Thus \eqref{eq:pf:thmS:rec} holds with
$\eps=(C_0\Theta^{-(K-5/2)})^f$, $a=(C_0/\Theta)^f$, where we put $e_i:=0$ outside $1\le i\le r+1$. Since $C_0\le C$,
$r(C_0/\Theta)^f\le1$, and \cref{pf:thmS:jointA} gives
$\sigma_{r+1}^2\le(1+C_J(r+1)(C_0/\Theta)^f)\prod(1+e_i)\sigma_0^2\le B\beta$. Finally $\varnothing\in W(\mathrm{bud}(r+1))$
and $Q(T,\eps,r+1,\varnothing)=\thmSSf(\eps,r+1,\varnothing,X_{r+1})$, which is therefore $\le\sigma_{r+1}^2\le B\beta$.
\end{proof}

\begin{corollary}[Theorem A with a density]\label{pf:thmS:recdens}
For $C_1,C_4\ge1$ there are $C,C'$ such that, under the hypotheses of \cref{pf:thmS:thmA} with $C'\le\Theta$ and
$r(C/\Theta)^f\le1$, if moreover $e_i=0$ for $1\le i\le r$ and $1+e_{r+1}=\max(s_d,1)$ for some $s_d\in\R$, then
$\sum_{\rho\in R}\|\sum_pc(\rho,p)V(p)\|^2\le(1+C(r+1)(C/\Theta)^f)\max(s_d,1)\,\mathrm{jn}$.
\formallabel{Formalism.FourierDecayAlt.P.RecursionWithDensities.J1.t_RecursionWithDens}
\formallabel{Formalism.FourierDecayAlt.P.recursion_with_densities}
\end{corollary}

\begin{proof}
Under these conditions $\prod_{i=1}^{r+1}(1+e_i)=\max(s_d,1)$.
\end{proof}

\subsubsection{Mass data and the main row}\label{sec:pf:thmS:mass}

\begin{definition}[Mass data]\label{pf:thmS:massdata}
Mass data $B$ on a finite set $P$ (of places) consist of real numbers: for $p\in P$, $g_p$, $\mathrm{glk}_p$,
$\mathrm{uk}_p$, $\mathrm{nn}_p$, $\mathrm{gg}_p$, $\mu^K_p$, $\mu^{K\natural}_p$, $\mu^{K''}_p$, $\mu^S_p$, $\mu^{S'}_p$,
$\mu^{S''}_p$, $\mu^{S''\prime}_p$, $M_p$, $\ell_p$; and global numbers $\mu^K_o,\mu^S_o,\mu^{S'}_o,M_o,\ell_\infty$,
$\mu_H,\mu_{H'},\mu_{H''},\mu_{H''\prime},\mu_N,\mu_{N'}$, volumes $v_{ZZ},v_{PP},v_G,v_{Ha},v_{Ha'},v_{Na},v_{Na'},v_B,v_H,v_{H'},v_N,v_{N'}$,
and $\alpha$, $C_{oo}$, $\deg$, $\dim\pi$, $\dim\pi'$, $\rho_0$, $\rho_0'$, $d$, $d'$, $n_{0c}$, $s_d$, $e_1$, $\delta$, $K_{gs2}$,
$K_V$, $\Sigma$, $K^{\mathrm{top}}_V$, $I_N$, $I_C$, $I_{OO}$. Put $\gamma_p=\mathrm{glk}_p/\mathrm{uk}_p$,
$\Gamma_1=\prod_p\gamma_p$, $\thmStt=\prod_p(1+3/g_p)$, $\mu^K=\mu^K_o\prod_p\mu^K_p$ and similarly $\mu^{K\natural}$,
$\mu^{K''}$, $\mu^S$, $\mu^{S'}$, $\mu^{S''}$, $\mu^{S''\prime}$, $M=M_o\prod_pM_p$, $\ell=\ell_\infty\prod_p\ell_p$,
$c_0=v_Nv_H/\dim\pi$, $c_0'=v_{N'}v_{H'}/\dim\pi'$, $D_*=\deg\cdot v_G/\mu^{K''}$, $\mathrm{AA}=v_BC_{oo}(M/v_{ZZ})/\mu^K$ and
$\mathrm{GS}=\thmStt\,(d/\dim\pi)\,K_V$. $\mathrm{MI}(B)$ is the conjunction of positivity of all these numbers
(with $s_d,\delta\ge0$) and the identities
\begin{gather*}
\Sigma=\rho_0K_V,\qquad n_{0c}\rho_0=\mathrm{AA},\qquad I_N=\frac{v_{ZZ}}{\ell\,v_{PP}}\,v_{PP}I_C,\qquad I_{OO}=C_{oo}I_C,\qquad
I_{OO}\le1/\deg,\\
v_Bv_G=v_{Ha}v_{Ha'},\quad v_{Na}=v_{Na'}=1,\quad \mu^S=\mu_N\mu_H,\quad\mu^{S''}=\mu_N\mu_{H''},\quad\mu^{S'}=\mu_{N'}\mu_{H'},\\
\mu^{S''\prime}=\mu_{N'}\mu_{H''\prime},\quad
\rho_0=\frac{v_{Ha}}{v_H\mu_H},\quad v_{Na}=v_N\mu_N,\quad\rho_0'=\frac{v_{Ha'}}{v_{H'}\mu_{H'}},\quad v_{Na'}=v_{N'}\mu_{N'},\\
d=\dim\pi\frac{v_{Ha}}{v_H\mu_{H''}},\quad d'=\dim\pi'\frac{v_{Ha'}}{v_{H'}\mu_{H''\prime}},\quad\frac{\mu_H}{\mu_{H''}}=
\frac{\mu_{H'}}{\mu_{H''\prime}}=\Gamma_1,\quad
\frac{M_o}{\ell_\infty}=\frac{(\mu^K_o)^2}{\mu^S_o\mu^{S'}_o},\\
\frac{M_p}{\ell_p}=\frac{\mu^K_p\mu^{K\natural}_p}{\mu^S_p\mu^{S'}_p},\quad
\mu^S_p=\gamma_p\mu^{S''}_p,\quad\mu^{S'}_p=\gamma_p\mu^{S''\prime}_p,\quad\mu^K_p=\frac{\mathrm{gg}_p}{\mathrm{nn}_p\mathrm{uk}_p}\mu^{K''}_p,\quad
\frac{\mu^{K\natural}_p}{\mu^K_p}=\frac{\mathrm{glk}_p\mathrm{nn}_p}{\mathrm{gg}_p}.
\end{gather*}
\end{definition}

\begin{lemma}[Mass identities]\label{pf:thmS:massid}
Let $B$ satisfy $\mathrm{MI}(B)$, $e_1>0$, $K_V\ge0$, $K_{gs2}\le\thmStt\,\Gamma_1\Sigma$, and
\begin{equation}\label{eq:pf:thmS:key}
c_0c_0'K_V\le(1+\delta)\,e_1\,n_{0c}I_N.
\end{equation}
Then $\mathrm{GS}=\thmStt\,\Gamma_1\Sigma$ and $K_{gs2}\le\mathrm{GS}\le(1+\delta)e_1\,\thmStt\,dd'/D_*$. The same conclusion
holds if \eqref{eq:pf:thmS:key} is replaced by $K_V\le s_dn_{0c}K^{\mathrm{top}}_V$, $K^{\mathrm{top}}_V\ge0$, $s_d\le e_1$ and
$c_0c_0'K^{\mathrm{top}}_V\le(1+\delta)I_N$.
\formallabel{Formalism.FourierDecayAlt.P.form_gives_theorem_A}
\formallabel{Formalism.FourierDecayAlt.Final.jx_formGivesMain}
\formallabel{Formalism.FourierDecayAlt.P.mass_formula_constant_one}
\formallabel{Formalism.FourierDecayAlt.Final.counting_theorem}
\end{lemma}

\begin{proof}
The second set of hypotheses implies \eqref{eq:pf:thmS:key}, since $c_0c_0'K_V\le c_0c_0's_dn_{0c}K^{\mathrm{top}}_V\le
e_1n_{0c}(1+\delta)I_N$. (Conversely, the formal proof derives the first version from the second by replacing $s_d$ by
$e_1$ and $K^{\mathrm{top}}_V$ by $K_V/(e_1n_{0c})$.) All quantities below are positive. First,
$d/\dim\pi=v_{Ha}/(v_H\mu_{H''})=\rho_0\Gamma_1$, so $\mathrm{GS}=\thmStt\Gamma_1\rho_0K_V=\thmStt\Gamma_1\Sigma$. Next,
\begin{gather*}
\Sigma=\rho_0K_V\le\frac{(1+\delta)e_1\,\rho_0n_{0c}I_N}{c_0c_0'}=\frac{(1+\delta)e_1\,\mathrm{AA}\,I_N}{c_0c_0'},\\
\mathrm{AA}\,I_N=\frac{v_BC_{oo}M}{v_{ZZ}\mu^K}\cdot\frac{v_{ZZ}I_C}{\ell}=\frac{v_BM\,I_{OO}}{\mu^K\ell}\le\frac{v_BM}{\mu^K\ell\deg}.
\end{gather*}
The local identities give $M/\ell=\mu^K\mu^{K\natural}/(\mu^S\mu^{S'})$, $\mu^{K\natural}_p=\gamma_p\mu^{K''}_p$ (hence
$\mu^{K\natural}=\Gamma_1\mu^{K''}$), $\mu^S=\Gamma_1\mu^{S''}$ and $\mu^{S'}=\Gamma_1\mu^{S''\prime}$; so
$M/(\ell\mu^K)=\mu^{K''}/(\Gamma_1\mu^{S''}\mu^{S''\prime})$. Further $1/v_N=\mu_N$ (as $v_N\mu_N=v_{Na}=1$), so
$1/c_0=\dim\pi\,\mu_N/v_H=d\,\mu_N\mu_{H''}/v_{Ha}=d\,\mu^{S''}/v_{Ha}$ and likewise $1/c_0'=d'\mu^{S''\prime}/v_{Ha'}$. Hence
\[
\Gamma_1\Sigma\le(1+\delta)e_1\,\frac{v_B\,\mu^{K''}}{\deg\,\mu^{S''}\mu^{S''\prime}}\cdot\frac{d\,\mu^{S''}}{v_{Ha}}\cdot
\frac{d'\mu^{S''\prime}}{v_{Ha'}}=(1+\delta)e_1\,\frac{dd'\,\mu^{K''}v_B}{\deg\,v_Bv_G}=(1+\delta)e_1\frac{dd'}{D_*},
\]
and multiplying by $\thmStt>0$ gives the claim.
\end{proof}

\begin{definition}[Local model]\label{pf:thmS:localmodel}
$\mathrm{LM}(D)$, for a TB-datum $D$ over a field $E$, means: (i) for every $i$ and finite set $Y$ of points
there are finitely many \emph{places} $p$ (each a topological ring $K_p$ with a subring $R_p$, an additive character
$\psi_p$, a local datum, invertible integers, and an embedding $E\to K_p$), maps $\eta_p$ from points to local points,
class maps $\mathrm{cl}_p$, phases $d_p$ and a real function $a$ with
$K_i(x,y)=a(x)a(y)\prod_p\one[\mathrm{cl}_p(x)=\mathrm{cl}_p(y)]\,d_p(x)\overline{d_p(y)}\,K^{\mathrm{loc}}_p(\eta_p(x),\eta_p(y))$
on $Y$, $K^{\mathrm{loc}}_p$ the local kernel of the place; (ii) for $1\le s\le i\le r+1$ and $Y$ of level $i-s$, the block
system of $D$ at $(i,s)$ on $Y$ (its kernels $\thmSblk(K_i,s,t,t')$, its base kernel $K_{i-s}$ and its lattices
$\mathrm{La}$, $\mathrm{Lap}$) is a rescaled product of local systems, each a top system (if $(i,s)=(r+1,1)$), a
pull-back system or an outside system; (iii) $K_{r+1}(\thmSjoin(h,1,t),\thmSjoin(h',1,t'))=0$ for $h,h'$ of level $r$ and
$t\ne t'$ nonzero; (iv) $e_i=0$ for $1\le i\le r+1$.
\end{definition}

The following components of the local and global inputs (\cref{pf:thmS:localcores,pf:thmS:globalcores}) are used in
the proof of the main row; \ref{pf:thmS:M3} rests on \cref{pf:thmS:sizesmain}, and the others are cited.
\begin{enumerate}[label=(M\arabic*)]
\item\label{pf:thmS:M1} \emph{Local-model axioms} (\lref{Formalism.FourierDecayAlt.P.GrowthPolynomialOperator.t_AxiomsOfLocalModel}):
given the local inputs (\cref{pf:thmS:localcores}: Schur product and domination lemmas for kernel systems,
pull-back and top systems, Gram and Parseval identities for local theta kernels) and the existence $\mathrm{A}_s$ of local theta majorants for blocks of every size $s\ge1$, every
TB-datum with $\mathrm{LM}$ satisfies $\mathrm{KP}$, $\mathrm{KS}$, $\mathrm{E}$ and
$\mathrm{DT}_{\forall}$. The statements $\mathrm{A}_s$ are part of a local lemma assembled from the local inputs and the theta step (existence of local majorants for blocks of size
$s$ implies it for $s+1$; \cref{pf:thmS:thetastep}).
\item\label{pf:thmS:M2} \emph{Lattice sums} (\lref{Formalism.FourierDecayAlt.Final.jx_instLatticeSums}): there is
$C_1\ge1$ such that for every CM field $E$ unramified over its real subfield and every instance
$\mathrm{Inst}(D,e,G,C,d_e,d'_e,T,T',l)$ with $g_1,g_2,\mathrm{len}_I>0$, $C\ge0$, $\Theta>4$ and $\Theta^{-G.f}$ the largest
element of the threshold set $\mathrm{Thr}_G(C,d_e,d'_e,T,T',l)$, one has $\mathrm{SB}(\Theta)$ and
$\mathrm{LS}((C_1\Theta^{-(K-5/2)})^f)$. (Here $\mathrm{Inst}$ says that $r=G.r$, $K=5$, the places of $D$ are
the infinite places of $E$ via $e$, and the lattices, sizes and volumes of $D$ are the explicit ideals, thresholds and
norms attached to $G$.)
\item\label{pf:thmS:M3} \emph{Sizes} (\lref{Formalism.FourierDecayAlt.Final.jx_sizesToRecursion}, from
the size bounds \cref{pf:thmS:sizesmain}): for $C_A,C_0\ge1$, $C_G\ge0$, $C_\Delta>4$, $c_T,C_b>0$,
$c_TC_b>1$ and $0\le C\le C_A$, $C'$, there is $k_1$ such that the size hypotheses $\mathrm{SH}$ (with threshold
$\Theta$), $k_1\le G.k$, $16\le c_{\exp}C_f$ and $C_f\log b\le G.f\log G.k$ imply $C'\le\Theta$, $4<\Theta$,
$G.r(C/\Theta)^{G.f}\le1$, $(G.r+1)(C/\Theta)^{G.f}\le1$ and $C(G.r+1)(C/\Theta)^{G.f}\le b^{-7}$.
\item\label{pf:thmS:M4} \emph{Positivity} (\lref{Formalism.FourierDecayAlt.Final.jx_globPositive}): for
$C_\Delta>4$, $C_b>0$ there is $k_3$ such that $G.k\ge k_3$ and the hypothesis $\mathrm{Hyp}_G(b,C_\Delta,C_b)$ on the global size datum imply
$g_1,g_2,\mathrm{len}_I>0$.
\item\label{pf:thmS:M5} \emph{Poisson summation} (\lref{Formalism.FourierDecayAlt.Final.jx_poissonStepJoint}, from
\lref{Formalism.FourierDecayAlt.Final.poisson_summation_text}): the Poisson input (lattices $\Lambda_i$ in a real
inner product space with frames and wave functions compatible with $D$) implies $\mathrm{PS}$.
\item\label{pf:thmS:M6} \emph{Density} (\lref{Formalism.FourierDecayAlt.Final.tc_densChain}, shared): if
$2r+2+f+5\le k$, the $q_p\ge\ell_p$ with $\ell_p$ prime, at most $f$ indices share the same prime, and $s_d$ is at most the
supremum of the products $\prod_{p\in S}\exp(4q_p^{-(k-2r-2)})$, then $s_d\le\exp(12f2^{-(k-2r-2)})$ and $s_d\le2$.
\item\label{pf:thmS:M7} \emph{Scaling the top level} (\lref{Formalism.FourierDecayAlt.Final.tc_recHypScaleTop}):
let $D^\lambda$ be $D$ with $K_{r+1}$ multiplied by $\lambda=\max(s_d,1)$ and $e_{r+1}$ replaced by $\lambda-1$. If $s_d\le2$,
$e_{r+1}=0$, $D$ satisfies the hypotheses of \cref{pf:thmS:thmA} (with some $c$), and $c'$ is supported on points of level
$r+1$, summable against $\|V\|$, with $s_dK_{r+1}-\sum_{\rho\in R}c'(\rho,\cdot)\overline{c'(\rho,\cdot)}$ PSD on finite sets of
such points, then $D^\lambda$ satisfies the hypotheses of \cref{pf:thmS:thmA} with $c'$, has the same $r,f$, and
$e^\lambda_i=e_i$ for $i\le r$. (Scaling preserves zeros and positivity of the kernels; $0\le\lambda-1\le1$; the
dominating kernels at the top are scaled by $\lambda$; and if $K$ is PSD and $s_dK-m$ is PSD then so is $\max(s_d,1)K-m$.)
\end{enumerate}

\begin{theorem}[The main row]\label{pf:thmS:mainrowthm}
Assume the local inputs, the global inputs, the lattice link, the local-model axioms and the theta step (see \ref{pf:thmS:M1}). For every $C_4\ge1$ and $C_{\mathrm{thr}}$ there are $C_A\ge\max(C_{\mathrm{thr}},1)$ and $0\le
C\le C_A$ such that for all $C_0\ge1$, $C_G\ge0$, $C_\Delta>4$, $c_T,C_b>0$ with $c_TC_b>1$ there is $k_0$ with the
following property. Let $E$ be a CM field unramified over its maximal real subfield, $D$ a TB-datum over $E$ with a
bijection $e$ between $\{1,\dots,f\}$ and the infinite places of $E$, $G$ a global size datum, $b,C_p,C_f,\Theta\in\R$, word
data $d_e,d'_e$, $T,T'\in\N$, $d_u,d_{u'}\in\R$, $l\in\N$, such that $k_0\le G.k$, the size hypotheses
$\mathrm{SH}(G,C_A,C_0,C_G,C_\Delta,c_T,C_b,b,C_p,\Theta,d_e,d'_e,T,T',d_u,d_{u'},l)$ hold (they include
$\mathrm{Hyp}_G(b,C_\Delta,C_b)$, $\Theta>0$ and that $\Theta^{-G.f}$ is the largest element of
$\mathrm{Thr}_G(C_A,d_e,d'_e,T,T',l)$), $16\le c_{\exp}(C_\Delta,c_T,C_b)C_f$, $C_f\log b\le G.f\log G.k$,
$\mathrm{Inst}(D,e,G,C_A,d_e,d'_e,T,T',l)$ and $\mathrm{LM}(D)$. Let $\mathrm{jn}\in\R$, $H$ a Hilbert space,
$V$, $c$ as above and $s_d\in\R$, with the real inputs ($\mathrm{GP}$, $\mathrm{WP}$, $\mathrm{EW}$, the
Poisson input, $\mathrm{EL}$, $\mathrm{EB}(C_4)$, $\mathrm{EN}(\mathrm{jn})$, $\mathrm{CO}(V)$) and the
density input $\mathrm{DI}(G.k,G.r,G.f,s_d)$ (that is, $2r+2+f+5\le k$ and the hypotheses of \ref{pf:thmS:M6}).
Let $B$ be mass data with $\mathrm{MI}(B)$, the energy link $I_N=\mathrm{jn}$, $e_1=\max(s_d,1)$,
$\delta=C(r+1)(C/\Theta)^f$, and $\mathrm{RI}(D,s_d,V,c,B)$: $c$ is supported on points of level $r+1$,
$s_dK_{r+1}-\sum_{\rho\in R}c(\rho,\cdot)\overline{c(\rho,\cdot)}$ is PSD on finite sets of such points, $\sum_p|c(\rho,p)|
\|V(p)\|<\infty$, $c_0c_0'K_V\le n_{0c}M$ whenever $M$ bounds all sums $\sum_{\rho\in R}\|\sum_pc(\rho,p)V(p)\|^2$,
$K_V\ge0$ and $s_d$ is the density of $B$. Then
\[
\Gamma_1\Sigma\le(1+\delta)e_1\frac{dd'}{D_*},\qquad0\le\delta\le b^{-7},\qquad1\le e_1\le\exp\bigl(12f\,2^{-(k-2r-2)}\bigr).
\]
\formallabel{Formalism.FourierDecayAlt.P.GrowthPolynomialOperator.t_MainRow}
\end{theorem}

\begin{proof}
\emph{Constants.} Let $C_1\ge1$ be the constant of \ref{pf:thmS:M2}, and $(C_T,C'_T)$ the constants of
\cref{pf:thmS:recdens} for $(C_1,C_4)$. Put $C_A=\max(C_{\mathrm{thr}},1,|C_T|)$ and $C=|C_T|$. Given
$C_0,C_G,C_\Delta,c_T,C_b$, let $k_1$ be given by \ref{pf:thmS:M3} for $(C_A,C_0,C_G,C_\Delta,c_T,C_b,|C_T|,C'_T)$, $k_3$ by
\ref{pf:thmS:M4}, and $k_0=\max(k_1,k_3)$.

\emph{Hypotheses of Theorem A.} By \ref{pf:thmS:M3}: $C'_T\le\Theta$, $4<\Theta$, $r(|C_T|/\Theta)^f\le1$ and
$|C_T|(r+1)(|C_T|/\Theta)^f\le b^{-7}$ (with $r=G.r=D.r$ by $\mathrm{Inst}$, and $f=G.f=D.f$ via $e$). By
\ref{pf:thmS:M4}, $g_1,g_2,\mathrm{len}_I>0$; by \ref{pf:thmS:M2}, $\mathrm{SB}(\Theta)$ and
$\mathrm{LS}((C_1\Theta^{-(K-5/2)})^f)$; $K=5\ge4$, and $f\ge1$ since $E$ has an infinite place. The
local-model axioms \ref{pf:thmS:M1} give $\mathrm{KP}$, $\mathrm{KS}$, $\mathrm{E}$, $\mathrm{DT}_{\forall}$;
\ref{pf:thmS:M5} gives $\mathrm{PS}$; and $\mathrm{LD}(V,0)$ holds since $K_{r+1}$ is PSD. So
$D$ satisfies the hypotheses of \cref{pf:thmS:thmA} with the coefficients $0$.

\emph{Density and scaling.} \ref{pf:thmS:M6} gives $s_d\le\exp(12f2^{-(k-2r-2)})$ and $s_d\le2$. By the local model,
$e_i=0$ for $1\le i\le r+1$. The first three clauses of $\mathrm{RI}$ are the hypotheses of \ref{pf:thmS:M7}, which
gives the hypotheses of \cref{pf:thmS:thmA} for $D^\lambda$, $\lambda=\max(s_d,1)$, with the coefficients $c$, and
$1+e^\lambda_{r+1}=\max(s_d,1)$, $e^\lambda_i=0$ for $i\le r$. \Cref{pf:thmS:recdens} gives, for every finite $R$,
\[
\sum_{\rho\in R}\Bigl\|\sum_pc(\rho,p)V(p)\Bigr\|^2\le\bigl(1+C_T(r+1)(C_T/\Theta)^f\bigr)\max(s_d,1)\,\mathrm{jn}\le(1+\delta)e_1I_N,
\]
using $C_T(C_T/\Theta)^f\le|C_T|(|C_T|/\Theta)^f$, the energy link, and $\mathrm{jn}\ge0$ (a limit of non-negative numbers by
$\mathrm{GP}$ and $\mathrm{EN}$).

\emph{Mass identities.} The fourth clause of $\mathrm{RI}$ now gives \eqref{eq:pf:thmS:key}. Apply
\cref{pf:thmS:massid} to $B$ with $K_{gs2}$ replaced by $\thmStt\Gamma_1\Sigma$: $\thmStt\Gamma_1\Sigma\le(1+\delta)e_1\thmStt
dd'/D_*$; since $\thmStt=\prod_p(1+3/g_p)>0$, $\Gamma_1\Sigma\le(1+\delta)e_1dd'/D_*$. Finally $\delta\ge0$, $\delta\le b^{-7}$ by
\ref{pf:thmS:M3}, $e_1=\max(s_d,1)\ge1$, and $e_1\le\exp(12f2^{-(k-2r-2)})$ since both $s_d$ and $1$ are.
\end{proof}

\begin{remark}\label{pf:thmS:latticelinkunused}
The lattice-link hypothesis of the main-row statement (and hence the theorem
\lref{Formalism.FourierDecayAlt.P.GrowthPolynomialOperator.t_LatticeLink} supplying it in
\cref{pf:thmS:pbapplied}) is not used in the proof of \cref{pf:thmS:mainrowthm}. Of the local and global inputs the proof uses the components named in \ref{pf:thmS:M1}--\ref{pf:thmS:M7}, Theorem~A, the
recursion with densities and the mass identities.
\end{remark}

% ---------------------------------------------------------------------------
\subsection{The loop rows}\label{sec:pf:thmS:loops}

\begin{theorem}[Loop rows]\label{pf:thmS:loops}
Assume $\mathrm{FD}$ and the mass-ratio identity for every datum. Let $\mathcal C$ be a class with a twist
selection, all of whose members satisfy the size condition $\sfname{TuSize}$ (eventually, the size data satisfy
a lower bound on the level primes and $\log T_u\le C_{\mathrm{Tu}}\,m\,\Lambda$ for every word $u$), and let $\sigma$ be a spectral selection for
$\mathcal C$. Then $\mathcal{L}_{3}(\mathcal C,\sigma)$.
\formallabel{GlobalTheory.LoopRowsByName.loopRows_of_rowsFinDim_alt}
\formallabel{GlobalTheory.T.loopRows_of_large_gap_discreteSeries}
\formallabel{GlobalTheory.LoopRowsCloseGap.loop_rows_closed_gap}
\end{theorem}

The proof reduces the loop rows to three ``row estimates'' along families. For data $D$, size data $L$, term ties $X$ (an identification of the terms of $D$ with explicit flags
and lattices) and loop ties $Y$ for a pair of words, the row estimates $\mathrm{RE}_{\mathrm{nd}}(C)$, $\mathrm{RE}_{\mathrm{deg}}(C)$ and $\mathrm{RE}_{1}(C)$ with constant $C$ bound the non-degenerate, degenerate and corank-one contributions to the loop rows
by explicit lattice sums with constant $C$. They are provided by \cref{pf:thmS:loopbounds}, which is proved, together with
\cref{pf:thmS:betauniq}, in \cref{sec:pf:thmS:loopsproofs}.

\begin{theorem}[Row estimates along families]\label{pf:thmS:loopbounds}
Assume the Siegel mass ratio,
the gap hypotheses $\mathrm{GH}$ (branching support, uniqueness of Bessel models with large gap,
multiplicity one of the lowest $K$-type, and the formal-degree identity) and the pointwise bound for discrete
series matrix coefficients. For all constants $(C_g,c_b,C_u,c_T,C_T,C_0,A,C_1,c_V,C_\Gamma)$ there are $C_*$ and
$\vartheta$ such that for every family $F$ with $C_\Gamma(F)=C_\Gamma$, $s_{C'}(F),s_{C_l}(F)\ge\vartheta$ and
$\mathrm{MF}(F;C_0,\dots)$: (i) eventually, under $\mathrm{CH}_0$, the volume--mass formula at $(D,u)$
holds for every word $u$ (for normalised flags and every set of row representatives, $g_a(D)$ times the mass of the rows
times $\dim\pi_u$ equals $d^S_u$); (ii) for all term ties $X$ and loop ties $Y$ along $F$ (with constant $C_1$) satisfying
the tie conditions, injectivity of the terms of word pairs, the $\beta$-lattice property of word pairs, and eventually
the small-cell volume bound for terms, and for every $C\ge C_*$: eventually, under $\mathrm{CH}$ with $A$, the three
row estimates hold with constant $C$.
\formallabel{GlobalTheory.T.loopRowBounds_along_family_largeGaps}
\end{theorem}

\begin{theorem}[The corank-one lattice property]\label{pf:thmS:betauniq}
Let $D$ be a datum with hermitian space $V$, $u,u'$ words, $\Lambda$ a lattice, $F$ a flag for $T_u$ and $F',F''$ flags for
$T_{u'}$, such that the fixed-lattice hypothesis holds for $(\Lambda,F)$ and $(F,F')$, $(F,F'')$ are generic terms over
$\Lambda$. Let $\beta_1,\beta_2$ be linear automorphisms of $V$ such that $(F,F',F'',\beta_\iota)$ is a pair of corank $s_0=1$
for $\iota=1,2$. Then $\beta_2x-\beta_1x\in X(F'')$ for all $x$ orthogonal to the span $X(F')$ of the flag $F'$.
\formallabel{beta_l_unique}
\end{theorem}

Finally we use the following statement of the development, the \emph{evaluation of the row estimates}
(\lref{Formalism.L2.SG.LoopsTies.loopRowS_of_raised}): for a spectral selection with raised
thresholds and loop ties $Y$, if for every $a_F$, closedly and for every $C$ above a threshold $C_{0}(a_F,\cdot)$, the three
row estimates hold eventually under the count hypotheses $\mathrm{CH}(\dots,C_0,a_F(c_{\mathrm{out}}))$, and the
volume--mass formula holds closedly, then $\mathcal{L}_{C_0}$ holds. (The proof evaluates the lattice sums of the
row estimates by the loop-kernel bounds, giving $\mathrm{R}_{\mathrm{big}}$ and $\mathrm{R}_{\mathrm{nsc}}$ with
$\thmSepsL(D)=e^{-\frac13kf\log k}$.)

\begin{lemma}\label{pf:thmS:loopsub}
Let $\mathcal C^g\subseteq\mathcal C$ be the subclass of families on which the gap condition holds eventually, with the
restricted selection. If the gap condition is closed on $\mathcal C$ ($\thmSClosed(\mathcal C,F\mapsto(F\in\mathcal C
\Rightarrow\mathrm{Gap}_{\mathrm{ev}}(F)))$) and $\mathcal{L}_{C_0}(\mathcal C^g,\sigma)$ holds, then
$\mathcal{L}_{C_0}(\mathcal C,\sigma)$ holds.
\formallabel{Formalism.L2.SG.LoopsTies.loopRowS_of_sub}
\end{lemma}

\begin{proof}
For a given $a_F$ take $\eps_1$ the smaller of the two values and, for $\eps\le\eps_1$, the supremum of the two staged
thresholds. For $F\in\mathcal C$ with large constants the second statement gives $\mathrm{Gap}_{\mathrm{ev}}(F)$, so $F\in\mathcal C^g$
and the first statement applies.
\end{proof}

\begin{proof}[Proof of \cref{pf:thmS:loops}]
The gap hypotheses $\mathrm{GH}$ and the Whittaker--Bessel statement are the two components of the uniqueness theorem for Bessel
models with large gap and the Whittaker--Bessel statement
(\lref{GlobalTheory.T.bessel_multiplicity_one_large_gap_and_whittaker_bessel}, via
\lref{GlobalTheory.GapIn.hyp_of_uniqGap}); the pointwise bound for discrete-series coefficients is
\lref{GlobalTheory.T.discreteSeriesCoef_pointwiseBound_unconditional}; Wu decay is
\lref{GlobalTheory.wuDecayOfGap}; and the rank row $\mathcal{R}_{3}(\mathcal C)$ holds by \cref{pf:thmS:ranks}.
Given these:

(a) $\sfname{TuSize}$ makes the gap condition closed on $\mathcal C$ (\lref{Formalism.L2.SG.GapClosed.gapEv_closed_of});
by \cref{pf:thmS:loopsub} it suffices to treat $\mathcal C^g$.

(b) The representation property of CM fields and the integrability of weighted coefficients (theorems of the development) and the two decay statements for matrix coefficients (components of $\mathrm{GH}$) are made available;
$\mathrm{Fin}$ holds by \cref{pf:thmS:finxc}; the existence of $R$-measures, the local Bessel $J$-factors of all word
pairs and the covolume formula are the shared theorems
\lref{GlobalTheory.RMeasAll.rMeasEx_all}, \lref{GlobalTheory.T.localBesselJFactor_all_words_of_gap}, and \lref{GlobalTheory.CovNFormulaPf.main}.

(c) \Cref{pf:thmS:loopbounds} with the Siegel mass ratio (\lref{Formalism.HermitianLattices.Ext.L3.siegelMassRatio}) provides
constants $C_*(\cdot)$, $\vartheta(\cdot)$ as functions of the constants; the stabiliser covolume bound
\lref{GlobalTheory.T.stabiliserCovolume_relevantTerm_le_heightPow} (for every relevant term with non-zero weight,
the covolume of its stabiliser is at most a power $C r_S^2$ of an explicit height) provides a constant $c_V(\cdot)$.

(d) The volume--mass formula holds closedly on $\mathcal C$: combine the rank row with \cref{pf:thmS:loopbounds}(i), enlarging the
staged threshold so that $s_{C'}(F),s_{C_l}(F)\ge\vartheta$.

(e) Apply the evaluation of the row estimates to $\mathcal C^g$, the selection with raised thresholds, $C_1=1$, the threshold function
$C_0(a_F,l)=C_*(\dots,a_F(\cdot),1,c_V(\dots),l_{17})$ and the loop ties built from the Bessel data and the measure data of
(b). Its hypothesis on the row estimates is \cref{pf:thmS:loopbounds}(ii) applied to the raised family: the tie conditions hold for
the raised selection, injectivity of terms holds, the $\beta$-lattice property is \cref{pf:thmS:betauniq}, the small-cell volume
bound follows from (c), and the constants are large by a further enlargement of the staged threshold (including
$s_{C_l}(F)\ge1$). With (d), the evaluation gives $\mathcal{L}_{3}(\mathcal C^g,\sigma)$.
\end{proof}

% ---------------------------------------------------------------------------
\subsection{Proofs of the deferred results}\label{sec:pf:thmS:deferred}

This subsection proves results which were cited above, or in other sections of Part~II, as statements of the formal
development: three elementary inputs of the carrier bound (\cref{sec:pf:thmS:elem}); the density of $K$-finite vectors and
its consequence for Bessel functionals (\cref{sec:pf:thmS:kfin}); the Capelli identity behind the capacity of the carrier
(\cref{sec:pf:thmS:capelli}); the leaf condition and the joint tensor identity of the carrier (\cref{sec:pf:thmS:joint}); the
lattice facts and the size bounds behind the main row (\cref{sec:pf:thmS:latfacts,sec:pf:thmS:sizes}); the theta step and
the local and global inputs of the main row (\cref{sec:pf:thmS:theta,sec:pf:thmS:cores}); and the row estimates along
families and the corank-one lattice property stated in \cref{sec:pf:thmS:loops} (\cref{sec:pf:thmS:loopsproofs}). Their
deeper ingredients are cited.

\subsubsection{Elementary inputs of the carrier}\label{sec:pf:thmS:elem}

For a datum $D$ with type primes $g_0,g_1$ we put $g_S(D)=g_0g_1$.

\begin{lemma}[The trace hypothesis]\label{pf:thmS:hyptt}
Let $D$ be a datum and $u,u'$ words of $D$. Then $T_u$ and $T_{u'}$ are primes, $T_u,T_{u'}\ge3$, and
$T_u\equiv T_{u'}\pmod{g_S(D)}$.
\formallabel{Formalism.MainNotions.SchurRelationsDiscreteSeries.hypTT_all}
\end{lemma}

\begin{proof}
The level primes are prime by definition of a datum. By \lref{Formalism.MainNotions.S.T_mod}, every level prime
satisfies $T_w\equiv1\pmod{g_S(D)^2}$ (it is a norm $\pi\bar\pi$ with $\pi\equiv1$ modulo an integer divisible by
$g_S(D)^2$). As $g_0,g_1\ge2$, $G=g_S(D)\ge4$. A prime $T\equiv1\pmod{G^2}$ is not $2$, since $2<G^2$; hence $T\ge3$.
Finally $T_u\equiv1\equiv T_{u'}$ modulo $G^2$, hence modulo $G$.
\end{proof}

\begin{lemma}[An elementary inequality]\label{pf:thmS:acle}
Let $c,C_l,L\in\R$ and $r,k,f\in\N$ with $C_l\ge1000+|c|$, $f\ge2$ and $L\ge\frac12$, and put $n=2r+k+1$. Then
\[
\max\bigl(c\,(r+1)^2(rn+1),\,0\bigr)\le\tfrac12\,C_l\,n^{19}f^5L.
\]
\formallabel{Formalism.L2.SG.Pf.CarrierXAll.ac_le}
\end{lemma}

\begin{proof}
Since $1\le r+1\le n$ and $rn+1\le n^2+1\le2n^2$, we have $0\le(r+1)^2(rn+1)\le2n^4\le2n^{19}$. Since $f^5\ge32$ and
$L\ge\frac12$, the right side is at least $8C_ln^{19}$, which is non-negative and at least $2|c|n^{19}\ge c\,(r+1)^2(rn+1)$.
\end{proof}

Let $X$ be a Hida setting for a datum $D$ and $C$ a choice for $X$, with parameter $\theta$. Put $k=k(D)$, $r=r_S(D)$,
$n=2r+k+1$, let $h$ be the hermitian form of $H_{k,r}$ on $\C^n$, $W_b$ the span of the $k+1$ middle basis vectors
$w_0,\dots,w_k$, $T(z)=-h(z,z)$, and let $U_W(z)$ be the stabiliser of $z$ and of the $2r$ hyperbolic basis vectors in
$\Uni(H_{k,r})$. For a word $u$ and an embedding type $\tau$ the choice transports the local data of $X$ at the place of
$\tau$ to $\Uni(H_{k,r})$: a vector $z_u=z_C(u,\tau)\in\C^n$, a representation $\pi_u$ of its stabiliser, the discrete series
representation $\sigma_\tau$ of $\Uni(H_{k,r})$ and a measure $\mu_G$ on $\Uni(H_{k,r})$; together with the stabiliser measures,
the measure on the unipotent group $N$ and the orbit measures of the choice they form the \emph{local window}
$W_C(u,u',\tau)$ of a pair of words.

\begin{lemma}[Goodness of the weights]\label{pf:thmS:goodw}
For words $u,u'$ and an embedding type $\tau$ the local window $W_C(u,u',\tau)$ is good: $\pi_u$ and $\pi_{u'}$ have the column
types of $u$ and $u'$ at $\tau$; $z_u,z_{u'}\in W_b$ with $T(z_u)=T_u$ and $T(z_{u'})=T_{u'}$; $\mu_G$ is a Haar measure on
$\Uni(H_{k,r})$; the stabiliser measures are Haar probability measures on $U_W(z_u)$ and $U_W(z_{u'})$; the measure on $N$ is
the Haar measure normalised by the Lie algebra with the constant $\kappa_N(D,\theta)$; and the two orbit measures are the
quotients of $\mu_G$ by the products of the stabiliser measures with the measure on $N$.
\formallabel{Formalism.L2.SG.Pf.AmpWinRX.goodW_all}
\end{lemma}

\begin{proof}
The conditions on the stabiliser measures, on the measure on $N$ and on the orbit measures are part of the definition of a
choice. By \lref{Formalism.L2.SG.LocWSOf.goodW_of}, the window is therefore good as soon as three transfers hold for $u$ and
$u'$: the column-type transfer ($\pi_u$ has the column type of $u$ at $\tau$), the $z$-transfer ($z_u\in W_b$ and $T(z_u)=T_u$)
and the Haar transfer ($\mu_G$ is a Haar measure). These hold for every choice, word and type by
\lref{Formalism.L2.SG.Pf.colsTypeTr}, \lref{Formalism.L2.SG.Pf.zTr}, and \lref{Formalism.L2.SG.Pf.haarTr} respectively.
\end{proof}

\subsubsection{$K$-finite vectors}\label{sec:pf:thmS:kfin}

Let $K_H=\Uni(H_{k,r})\cap\Uni(n)$, a compact group, and let $\sigma$ be a continuous unitary representation of $\Uni(H_{k,r})$ on a
Hilbert space $\mathcal H$. For a word $w=(Z_1,\dots,Z_\ell)$ of elements of the Lie algebra of $\Uni(H_{k,r})$ write
$d\sigma(w)=d\sigma(Z_1)\cdots d\sigma(Z_\ell)$; a vector $v$ is \emph{smooth} if all $d\sigma(w)v$ are defined, and
\emph{$K$-finite} if it lies in a finite-dimensional subspace stable under $\sigma(K_H)$.

\begin{theorem}[Density of $K$-finite vectors]\label{pf:thmS:kfindense}
Let $\sigma$ be a discrete series representation of $\Uni(H_{k,r})$ with a Harish-Chandra parameter $\lambda$ (for some
signature). For every smooth vector $v$, every finite list $w_1,\dots,w_L$ of words in the Lie algebra and every $\delta>0$
there is a $K$-finite vector $a$ with $\sum_{j=1}^L\|d\sigma(w_j)(v-a)\|<\delta$.
\formallabel{Formalism.L2.SG.Pf.CarrierXAll.kFinDense_holds}
\formal{GlobalTheory.KFinDenseDS.kFin_dense}
\formal{GlobalTheory.KFinAvg.kDense}
\end{theorem}

\begin{proof}
A discrete series representation is a continuous unitary representation, and for a smooth vector $v$ and a word $w$ the map
$g\mapsto d\sigma(w)\sigma(g)v$ is continuous. Let $\mu$ be the Haar probability measure of $K_H$ and
\[
p(g)=\tfrac12\bigl(1+\tfrac1n\Re\tr g\bigr)\qquad(g\in\Uni(n)).
\]
For unitary $g$ the eigenvalues have modulus one, so $|\tr g|\le n$ and $0\le p(g)\le1$; moreover $p(g)=1$ only if all
eigenvalues equal $1$, that is $g=1$, unitary matrices being diagonalisable. By compactness, for every neighbourhood $U$ of
$1$ there is $\delta_U<1$ with $p\le\delta_U$ on $\Uni(n)\setminus U$. The power $p^m$ is a polynomial of degree at most $m$ in
the entries of $g$ and $\bar g$; these polynomials form a finite-dimensional space $\mathcal P_m$, stable under the
translations $b\mapsto b(x\,\cdot)$ because the entries of $xg$ are linear in those of $g$.

Put $p_m=p^m/\int_{K_H}p^m\,d\mu$; the denominator is positive because $p$ is continuous with $p(1)=1$ and $\mu$ charges every
non-empty open subset of $K_H$. The $p_m$ form an approximate identity at $1$: for continuous $\phi\colon K_H\to\mathcal H$,
$\int p_m\phi\,d\mu\to\phi(1)$. Indeed, given $\eta>0$ choose a neighbourhood $U$ of $1$ with $\|\phi-\phi(1)\|\le\eta$ on
$U\cap K_H$, a number $\delta'\in(\delta_U,1)$ and a neighbourhood $V\subseteq U$ of $1$ on which $p\ge\delta'$; then
\[
\Bigl\|\int p_m(\phi-\phi(1))\,d\mu\Bigr\|\le\eta+2\sup\|\phi\|\cdot\frac{\delta_U^m}{\mu(V\cap K_H)\,\delta'^m},
\]
and the last term tends to $0$.

Let $a_m=\int_{K_H}p_m(g)\,\sigma(g)v\,d\mu(g)$. (i) $a_m$ is $K$-finite: for $x\in K_H$, invariance of $\mu$ gives
$\sigma(x)a_m=\int p_m(x^{-1}g)\sigma(g)v\,d\mu(g)$ with $g\mapsto p_m(x^{-1}g)$ in $\mathcal P_m$, so $\sigma(x)a_m$ lies in the
space $\{\int b(g)\sigma(g)v\,d\mu(g):b\in\mathcal P_m\}$, which is finite-dimensional and, by the same computation,
$\sigma(K_H)$-stable. (ii) For each $j$, differentiating under the integral sign (legitimate since $v$ is smooth and $K_H$ is
compact) and using $\int p_m\,d\mu=1$,
\[
d\sigma(w_j)(v-a_m)=\int_{K_H}p_m(g)\bigl(d\sigma(w_j)v-d\sigma(w_j)\sigma(g)v\bigr)\,d\mu(g),
\]
which tends to $0$ by the approximate identity, applied to the continuous function $g\mapsto d\sigma(w_j)\sigma(g)v$. Choose $m$
with $\|d\sigma(w_j)(v-a_m)\|<\delta/(L+1)$ for all $j$; then $a=a_m$ satisfies
$\sum_j\|d\sigma(w_j)(v-a)\|\le L\delta/(L+1)<\delta$.
\end{proof}

Let now $B$ be Wu--Bessel data for the choice $C$. For a word $u$ and a type $\tau$ they provide a \emph{Bessel functional}
$b=b_B(u,\tau)$ on the smooth vectors of $\sigma_\tau$, which is additive and satisfies a Sobolev bound: there are $C_w\in\R$ and
finitely many words $w_j$ with $|b(y)|\le C_w\sum_j\|d\sigma_\tau(w_j)y\|$ for smooth $y$. The \emph{local Bessel function} is
$\varphi_B(u,\tau)(m)=b(\sigma_\tau(m)v_0)$, where $v_0$ is the base vector of the lowest $K$-type.

\begin{proposition}[Bessel functionals on $K$-finite vectors]\label{pf:thmS:kfinnz}
Assume the conclusion of \cref{pf:thmS:kfindense}. Let $D$ be a datum satisfying the gap hypothesis, $X$ a Hida setting for $D$,
$C$ a choice and $B$ Wu--Bessel data for $C$, $u$ a word and $\tau$ an embedding type. Then there is a $K$-finite vector $a$ for
$\sigma_\tau$ with $b_B(u,\tau)(a)\ne0$.
\formallabel{Formalism.L2.SG.Pf.CarrierXAll.kFinNZ_of}
\end{proposition}

\begin{proof}
Write $\sigma=\sigma_\tau$, $b=b_B(u,\tau)$ and $\varphi=\varphi_B(u,\tau)$. Under the gap hypothesis every choice satisfies the
local factor conditions (\lref{Formalism.L2.SG.Pf.LocF.locFactor_holds}), and then $\varphi$ has non-zero finite norm in
$L^2(\nu_S)$, $\nu_S$ the orbit measure of the choice (\lref{Formalism.L2.SG.Pf.GoodS.l2SG_phiOf_ne_zero_and_lt_infty}). The
vector $z_u$ lies in $W_b$ with $T(z_u)=T_u>0$ (\lref{Formalism.L2.SG.Pf.zTr}) and $\mu_G$ is a Haar measure
(\lref{Formalism.L2.SG.Pf.haarTr}), so $\nu_S$ gives measure zero to the complement of $\Uni(H_{k,r})$
(\lref{Formalism.L2.SG.Pf.GoodS.nu_conc}, which states this for the quotient of a Haar measure of $\Uni(H_{k,r})$ by a Haar
measure of the stabiliser of such a vector). If $\varphi$ vanished on $\Uni(H_{k,r})$ it would vanish $\nu_S$-almost everywhere;
hence there is $m\in\Uni(H_{k,r})$ with $b(\sigma(m)v_0)=\varphi(m)\ne0$.

The representation $\sigma$ is a discrete series representation with the Harish-Chandra parameter $\lambda_D(\tau)$
(\lref{Formalism.L2.SG.Pf.dsTr} and \lref{Formalism.MainNotions.S.Inh.amp_isHCParam_lamD}); its base vector $v_0$ is a lowest
$K$-type vector (\lref{Formalism.L2.SG.Pf.isLowUnitH_base_vector}) and hence smooth
(\lref{Formalism.L2.SG.Pf.WuDec.lowSmoothSK}); so $y=\sigma(m)v_0$ is smooth
(\lref{Formalism.MainNotions.LowestKTypeMultiplicity.smooth_orbit}). Let $C_w$, $w_j$ be the constants of the Sobolev bound
for $b$ and $\delta=|b(y)|/(|C_w|+1)>0$. By density there is a $K$-finite $a$ with $\sum_j\|d\sigma(w_j)(y-a)\|<\delta$.
$K$-finite vectors are smooth (\lref{Formalism.L2.SG.Pf.k_finite_smooth}), so $y-a$ is smooth and $b(y)=b(y-a)+b(a)$. If
$b(a)=0$, then
\[
|b(y)|=|b(y-a)|\le C_w\sum_j\|d\sigma(w_j)(y-a)\|\le|C_w|\,\delta<|b(y)|,
\]
a contradiction. Hence $b(a)\ne0$.
\end{proof}

\subsubsection{The Capelli identity}\label{sec:pf:thmS:capelli}

For $z\in W_b$, a \emph{frame} of $z$ is a family $B=(B_1,\dots,B_k)$ in $W_b$ with $h(B_a,B_b)=\delta_{ab}$ and $h(B_a,z)=0$.
Let $\pi$ be a continuous representation of $U_W(z)$ on a finite-dimensional Hilbert space $W$ and $B$ a frame. Let $E^B_{ab}$
be the endomorphism $x\mapsto B_a\,h(x,B_b)$ of $\C^n$; its hermitian and skew-hermitian parts lie in the Lie algebra of
$U_W(z)$, and we put $E_{ab}=d\pi(E^B_{ab})$, where $d\pi$ is the complex-linear extension of the differential of $\pi$. The
\emph{Capelli operator} is the column determinant
\[
\mathrm{Cap}_B(\pi)(w)=\sum_{\varsigma\in S_k}\operatorname{sgn}(\varsigma)\,(E+w)_{\varsigma(1)1}(E+w-1)_{\varsigma(2)2}\cdots
(E+w-k+1)_{\varsigma(k)k}\qquad(w\in\C),
\]
where $(E+x)_{ab}=E_{ab}+x\,\delta_{ab}$ and products are compositions in the written order. We say that $\pi$ has
\emph{Harish-Chandra type} $u\in\R^k$ if there is $g\in\Uni(H_{k,r})$ fixing the hyperbolic basis vectors with
$g(\sqrt{T(z)}\,w_k)=z$ such that $x\mapsto\pi(g\,\mathrm{blk}(x)\,g^{-1})$, $\mathrm{blk}$ the block embedding of $\Uni(k)$ used in
the development, is a non-zero irreducible unitary representation of $\Uni(k)$ with highest weight $\Lambda(u)$, the lowest
$K$-type attached to $u$; a direct computation with its definition gives $\Lambda(u)_j=u_j-\frac{k-1}2+(j-1)$ ($1\le j\le k$).

\begin{theorem}[Capelli identity]\label{pf:thmS:capelli}
Let $z\in W_b$ with $T(z)>0$, and let $\pi$ be a continuous representation of $U_W(z)$ on a finite-dimensional Hilbert space $W$,
of Harish-Chandra type $u\in\R^k$. Then for every frame $B$ of $z$ and every $w\in\C$,
\[
\mathrm{Cap}_B(\pi)(w)=\prod_{j=1}^k\Bigl(w-\frac k2+\frac12+u_j\Bigr)\cdot\mathrm{id}_W.
\]
\formallabel{Formalism.L2.SG.Pf.CarrierT.capacity_from_hermitian_char}
\end{theorem}

\begin{proof}
Let $g$ be as in the definition of the Harish-Chandra type; the vectors $B^g=(gw_0,\dots,gw_{k-1})$ form a frame of $z$
(\lref{Formalism.FourierDecayAlt.P.QrCap.fr_frame}).

\emph{Step 1.} For every frame $B$ the matrices $E^B_{ab}$ lie in the complexified Lie algebra of $U_W(z)$
(\lref{Formalism.FourierDecayAlt.P.QrD.frame_lie}), and on such matrices the complexified differential of $\pi$ is the
combination of the differentials of the hermitian and skew-hermitian parts (\lref{Formalism.FourierDecayAlt.P.QrD.ecRep_eq});
hence $\mathrm{Cap}_B(\pi)(w)$ is the Capelli determinant $\operatorname{qdet}(E_B)(w)$ of the $k\times k$ matrix
$E_B=(E_{ab})$ with entries in $\End(W)$.

\emph{Step 2.} The Capelli determinant does not depend on the frame: $\operatorname{qdet}(E_B)(w)=\operatorname{qdet}(E_{B'})(w)$
for any two frames of $z$ (\lref{Formalism.FourierDecayAlt.P.QrCap.capChange}; the change of frame $P_{ca}=h(B_a,B'_c)$ is
invertible with inverse $Q_{bd}=h(B'_d,B_b)$, and the Capelli determinant is invariant under the induced change of $E_B$).

\emph{Step 3.} For the frame $B^g$, $\operatorname{qdet}(E_{B^g})(w)$ acts on $W$ as $\prod_{j=1}^k(w+\Lambda(u)_j-(j-1))$
(\lref{Formalism.FourierDecayAlt.P.QrCap.core}). This is the classical evaluation of the Capelli determinant on an irreducible
module of highest weight $\Lambda$ \cite{HoweUmeda1991}: the operators $E_{ab}$ of the frame $B^g$ satisfy the commutation relations
of $\mathfrak{gl}_k$, the diagonal torus acts through them by $\pi(g\,\mathrm{blk}(\mathrm{diag}(e^{i\theta_c}))\,g^{-1})=
\exp(\sum_ci\theta_cE_{cc})$, the Capelli determinant is central and so acts on the irreducible module $W$ by a scalar, and the
scalar is computed on a highest weight vector, which is annihilated by the raising operators.

\emph{Step 4.} Since $\Lambda(u)_j-(j-1)=u_j-\frac{k-1}2$, the scalar of Step~3 is $\prod_j(w-\frac k2+\frac12+u_j)$, and Steps
1--3 give the claim.
\end{proof}

\subsubsection{The leaf condition and the joint tensor identity}\label{sec:pf:thmS:joint}

We use the carrier construction of Part~II for $X$, $C$, Wu--Bessel data $B$, words $u,u'$ and a bijection $e$ between the
infinite places of $E$ and the embedding types which is compatible with $X$, under the normalisation $v(D)=\mathrm{rd}_E/2$ of
the host. The places are indexed by $j$, with types $\tau_j$. For vectors $a_j,b_j$ of $\sigma_{\tau_j}$,
$\xi_j\in V_{u,v_j}$, $\xi'_j\in V_{u',v_j}$ and scalars $c_j$, the \emph{coefficient function} $f_{c,j}$ is the
operator-valued matrix coefficient of $(a_j,b_j)$, multiplied by $c_j$ and by the rank-one operator attached to
$(\xi_j,\xi'_j)$, and conjugated by a fixed torus element of the datum. From these coefficient functions, grids, a cut, indices
$i'_0$, top lifts, type projections and a constant $A_0$, the construction produces a \emph{carrier} $\mathcal P$: a Hida setting
over $E$ with level towers, the test functions $f_{c,j}$, parameters $\nu'_j$, $m_{0,j}$ and isometries
$e_P(v)\colon W_p(v)\to V_{u,v}$, $e_Q(v)\colon W_q(v)\to V_{u',v}$. The carrier satisfies the \emph{pinning covariance} if for
every infinite place $v$, every $\gamma$ in the local unitary group and all measures $\mu_{R,v},\nu$ satisfying the pin hypotheses
($G_v(\gamma)>0$, the matrix of $\gamma$ is monomial and its transfer $\mathrm{tr}_v(\gamma)$ to the model group is relevant for
the base tower, $\mu_{R,v}$ is a Haar measure on the stabiliser of $\gamma$ and $\nu$ the corresponding quotient measure), the
local orbital integral equals the transported model orbital integral:
\[
J_v\bigl(\nu;W_{u,v},W_{u',v};\gamma\bigr)=e_P(v)\circ\mathrm{orb}_v\bigl(\mathrm{tr}_v(\gamma),\mathrm{tr}_{v*}\mu_{R,v},f_{c,v}\bigr)\circ
e_Q(v)^{-1}.
\]
A function $F$ on the model group lies in the \emph{coefficient class} $\mathrm{CC}(A,m)$ if it is of class $C^m$ and
$\mathrm N(g)^A\|\partial_L(w)\partial_R(w')F(g)\|$ is bounded for all words $w,w'$ of total length at most $m$ in unit Lie algebra
elements, $\mathrm N$ being the norm function of the development.

\begin{proposition}[The leaf condition]\label{pf:thmS:xleaf}
Assume the uniqueness of Bessel models with large gap. Let $D$ be a datum with size data satisfying the count hypotheses
$\mathrm{CH}_0$ with constants $(C',C_l,\dots)$, $C_l\ge1000$, and let $X,C,B,u,u',e$ be as above. Assume the trace hypothesis for
$(u,u')$, the goodness of the weights and of the Harish-Chandra parameters at every place, $T_u,T_{u'}>0$, and let
$(a_j,b_j,\xi_j,\xi'_j)$ be pinned at scale $1$ ($a_j,b_j$ are $K$-finite and
$\langle b_B(u,\tau_j)(a_j),\xi_j\rangle\langle\xi'_j,b_B(u',\tau_j)(b_j)\rangle=1$), with $\mathrm{scl}(B,u,u',\tau_j)\ne0$ and
$c(B,u,\tau_j)\ne0$ for all $j$. Let $h_j\in\R$ be such that $\lambda_D(\tau_j)$ is a Harish-Chandra parameter with gap at least
$h_j$ and the gap condition of Wu decay holds at $h_j$, and let $A_c\ge0$ with $A_c\ge C_{\mathrm{orb}}(r+1)^2(rn+1)$ and
$A_c\le h_j/2$ for all $j$, where $C_{\mathrm{orb}}$ is the constant of
\lref{Formalism.FourierDecayAlt.P.orbital_integral_converges} (convergence of the orbital integrals of the Bessel kernels).
Let $c_j$ be the normalising scalars of the carrier, defined from $c(B,u,\tau_j)$, $\mathrm{scl}(B,u,u',\tau_j)$ and the
amplitude vector provided by the uniqueness of Bessel models. Then the leaf condition holds for $(a,b,\xi,\xi',c)$: the $a_j,b_j$
are $K$-finite, and for all top lifts, type projections and $A_0$ the carrier built from the coefficient functions of
$(a,b,\xi,\xi',c)$ satisfies the pinning covariance.
\formallabel{GlobalTheory.CarrierGap.xLeafAt_of_pinsG}
\end{proposition}

\begin{proof}
$K$-finiteness of $a_j,b_j$ is part of the pinning. Fix the top lifts, type projections and $A_0$. By
\lref{Formalism.L2.SG.Pf.CoefCCT.fcT_inCC}, every coefficient function $f_{c,j}$ lies in $\mathrm{CC}(A_c,m)$ for every $m$:
this lemma states that for $K$-finite vectors of a discrete series representation whose Harish-Chandra parameter has gap at least
$h_j$, with the gap condition of Wu decay at $h_j$ and $0\le A_c\le h_j/2$, the conjugated and rescaled operator-valued matrix
coefficients lie in these classes. Then \lref{GlobalTheory.CarrierGap.pinCov_ccOneG} gives the pinning covariance: it states
that under the hypotheses of the proposition, if every coefficient function with the normalising scalars $c_j$ lies in
$\mathrm{CC}(A_c,1)$, the carrier built from them satisfies the pinning covariance for all top lifts, type projections, $A_0$ and
models. (Its proof reduces the covariance to an identity, at each place, between the local orbital integral and the model orbital
integral of $f_{c,j}$, which holds for the normalising scalar $c_j$.)
\end{proof}

\begin{theorem}[Joint tensor identity]\label{pf:thmS:jointtensor}
Let $X,C,u,u',e$ be as above, and let $\mathcal P$ be the carrier built, with its standard models, from arbitrary constants,
grids, coefficient functions $f_{c,j}$, a cut, indices $i'_0$, top lifts, type projections and $A_0$. Let $i_0=i'_0$,
$\mathrm{sc}>0$, and let $\nu_q$ be pinned quotient measures: $\nu_q(\gamma)=\mu_\infty$ for regular $\gamma$, and for the
representative $\gamma_w$ of a point $w$ with a rational representative, $\nu_q(\gamma_w)$ is the quotient of a Haar measure on the
stabiliser $\mathrm{St}_\infty(\gamma_w)$ whose image in coordinates is the stabiliser measure of $w$. Assume that $\mathcal P$ is
built and satisfies the pinning covariance. Then for every tower point $p$ whose point $w$ of level $r+1$ has a rational
representative and $G_v(\gamma_w)>0$ at all infinite places $v$, all frames $(h,h')$ adapted to $\gamma_w$ with scale
$\mathrm{sc}$ and indices $i_0,i'_0$, and all $x=(x_v)_v$,
\[
\pi_u(h^{-1})\,\mathrm{OO}\bigl(\nu_q(\gamma_w);W_u,W_{u'};\gamma_w\bigr)\,\pi_{u'}(h')\Bigl(\bigotimes_vx_v\Bigr)
=\bigotimes_ve_P(v)\Bigl(\frac{\mathrm{vec}_v(\mathrm{std}(p_v))}{b_{\mathrm{st}}(\kappa_G,p_v)}\,e_Q(v)^{-1}x_v\Bigr).
\]
Here $\mathrm{OO}$ is the archimedean orbital integral of the global Whittaker--Bessel data $W_u,W_{u'}$, $p_v$ the localisation
of $p$ at $v$ and $\mathrm{std}(p_v)$ its standardisation, $\mathrm{vec}_v$ the model vector operator of the base tower at $v$ with
the parameters $\nu'_v,m_{0,v}$ and the test function $f_{c,v}$, and $b_{\mathrm{st}}(\kappa_G,\cdot)$ the standard volume factor
with $\kappa_G=\kappa(g,\mathrm{rd}_E)$.
\formallabel{Formalism.L2.SG.Pf.CarrierXAll.jointTensor_compact_test}
\end{theorem}

\begin{proof}
Call a point $w$ \emph{admissible} if it has a rational representative and $G_v(\gamma_w)>0$ for all $v$.

\emph{Inputs.} The local stabilisers of the standard vectors have the geometric description of
\lref{Formalism.L2.SG.Pf.Carrier.JointQuot.stVGeo_of_sec}, applied with the section of
\lref{Formalism.L2.S.L6.UnitaryStandardVectors.stVSec_holds}: at every admissible $w$ and every place $v$ the local stabiliser
$\mathrm{St}_v(\gamma_w)$ is a subgroup, parametrised continuously by its free coordinates, with translations and inversion
measure-preserving in coordinates. From this, \lref{Formalism.L2.SG.Pf.Carrier.JointQuot.locUni_of} provides, for every Haar
measure $\mu_R$ on the global stabiliser normalised by the stabiliser measure, inversion-invariant normalised Haar measures
$\mu_{R,v}$ on the local stabilisers and measurable functions $b_v<\infty$ with $\int b_v(sx)\,d\mu_{R,v}(s)=1$ for all $x$; and
\lref{Formalism.L2.SG.Pf.Carrier.JointQuot.globProd_of} shows $\mu_R=\prod_v\mu_{R,v}$. The characters of the Whittaker--Bessel
data satisfy $\psi_u(s)\overline{\psi_{u'}(\gamma_w^{-1}s\gamma_w)}=1$ on the stabilisers of admissible points
(\lref{Formalism.L2.SG.Pf.Carrier.JointQuot.stRel_holds}). By construction, the carrier built from $X$ is compatible with
$X$, its base frames are those of $X$ (as $i_0=i'_0$), and its embeddings of $E$ are those of $X$ at the corresponding places.

\emph{Step 1: factorisation.} By \lref{Formalism.L2.SG.Pf.Carrier.JointQuot.ooQuot_pin} there are local measures
$\nu_l(w,v)$, quotients of the normalised local Haar measures of the stabilisers, such that for every admissible $w$
\[
\mathrm{OO}\bigl(\nu_q(\gamma_w);W_u,W_{u'};\gamma_w\bigr)\Bigl(\bigotimes_vy_v\Bigr)=\bigotimes_vJ_v\bigl(\nu_l(w,v);W_{u,v},W_{u',v};
\gamma_{w,v}\bigr)y_v :
\]
the global quotient measure is the product of the local ones (Fubini, with $\prod_vb_v$ as cut-off function), and the integrand
is a tensor product of local integrands.

\emph{Step 2: the places.} Fix $v$. Three statements about carriers built from $X$ --
\lref{Formalism.L2.SG.Pf.Carrier.relTr_carrierOf} (the transfer of $h_v^{-1}\gamma_{w,v}h'_v$ to the model group is relevant
for the base tower), \lref{Formalism.L2.SG.Pf.Carrier.measPinAt_of} (the transfer of the conjugated local Haar measure of the
stabiliser is $b_{\mathrm{st}}(\kappa_G,p_v)$ times the model measure of the tower at the transferred element) and
\lref{Formalism.L2.SG.Pf.Carrier.orbVecAt_carrierOf} (given this, the model orbital integral of any test function at the
transferred element is $b_{\mathrm{st}}(\kappa_G,p_v)^{-1}$ times the model vector operator at $\mathrm{std}(p_v)$) -- combine by
\lref{Formalism.L2.SG.Pf.Carrier.JointPinJB.jointBridge_of_open} into a bridge: there are $\gamma''$ and measures
$\mu_{R,v},\nu$ satisfying the pin hypotheses such that
$\pi_{u,v}(h_v^{-1})\,J_v(\nu_l(w,v);\gamma_{w,v})\,\pi_{u',v}(h'_v)=J_v(\nu;\gamma'')$, and the model orbital integral of
$f_{c,v}$ at $\mathrm{tr}_v(\gamma'')$ equals $b_{\mathrm{st}}(\kappa_G,p_v)^{-1}\mathrm{vec}_v(\mathrm{std}(p_v))$. The pinning
covariance identifies $J_v(\nu;\gamma'')$ with $e_P(v)\circ\mathrm{orb}_v(\cdots)\circ e_Q(v)^{-1}$. Hence
\[
\pi_{u,v}(h_v^{-1})\,J_v\bigl(\nu_l(w,v);\gamma_{w,v}\bigr)\,\pi_{u',v}(h'_v)=e_P(v)\circ
\frac{\mathrm{vec}_v(\mathrm{std}(p_v))}{b_{\mathrm{st}}(\kappa_G,p_v)}\circ e_Q(v)^{-1}.
\]

\emph{Step 3.} The adapted frames act on the global tensor product factor by factor,
$\pi_u(h^{\pm1})\bigotimes_vy_v=\bigotimes_v\pi_{u,v}(h_v^{\pm1})y_v$. Combining this with Steps 1 and 2 gives the identity.
\end{proof}

\subsubsection{Lattice facts}\label{sec:pf:thmS:latfacts}

Let $E$ be a CM field with maximal real subfield $E^+$. We call $E$ \emph{unramified} if the different $\mathfrak D$ of $E$ is
generated by the different of $E^+$. Put $\deg=[E^+:\Q]$ and $\mathrm{rd}=|d_E|^{1/[E:\Q]}$. For $g\in\N_{>0}$ and $k,\nu\in\Z$ let
$\mathfrak a(g,k,\nu)=\mathfrak D^\nu g^k$. For fractional ideals $A,B$ and $s\in\N$ let
$\mathrm{tw}_s(A,B)=\{x\in A:x-(-1)^s\bar x\in B\}$, and put
\begin{gather*}
P=\mathrm{tw}_s\bigl(\mathfrak a(g,1,\nu),\mathfrak a(g,2,\nu)\bigr),\qquad
P^\sharp=\mathrm{tw}_s\bigl(\mathfrak a(g,-1,-\nu),\mathfrak a(g,0,-\nu)\bigr),\\
\mathfrak a^\vee=\mathfrak a(g,-1,-\nu-1),\qquad\mathfrak a^\sharp=\mathfrak a(g,-1,-\nu).
\end{gather*}
For a $\Z$-module $L\subseteq E$ let $L^{*g}=\{x\in E:\Tr_{E/\Q}(xy/g)\in\Z\text{ for all }y\in L\}$, and for $c\in\R$ let
$\Phi_c(L)=c\,\iota(L)\subseteq E_\R$, where $\iota$ is the mixed embedding.

\begin{theorem}[Lattice facts]\label{pf:thmS:latfacts}
Let $E$ be an unramified CM field, $g>0$, $\nu\in\Z$ and $s\in\N$. Then:
\begin{enumerate}
\item $\mathfrak a(g,2,\nu)\subseteq P\subseteq\mathfrak a(g,1,\nu)$ and $[\mathfrak a(g,1,\nu):P]=[P:\mathfrak a(g,2,\nu)]=g^{\deg}$;
\item $P^{*g}\subseteq\mathfrak a^\vee$ and $[\mathfrak a^\vee:P^{*g}]=g^{\deg}$;
\item $P^\sharp\subseteq\mathfrak a^\sharp$ and $[\mathfrak a^\sharp:P^\sharp]=g^{\deg}$;
\item $\sqrt{N(\mathfrak a^\vee)}=(g\,\mathrm{rd}^{\nu+1})^{-\deg}$;
\item $\Phi_{\mathrm{rd}^{-\nu}}(P)$ is a discrete subgroup of $E_\R$ which spans $E_\R$, of covolume $(\mathrm{rd}/2)^{\deg}g^{3\deg}$.
\end{enumerate}
\formallabel{Formalism.FourierDecayAlt.P.lattice_facts_general}
\end{theorem}

\begin{proof}
We use the following facts.
\begin{enumerate}[label=(\alph*)]
\item $[E:\Q]=2\deg$, and $E$ has $\deg>0$ infinite places (\lref{Formalism.FourierDecayAlt.P.degreeFacts_statement}).
\item Each $\mathfrak a(g,k,\nu)$ is stable under complex conjugation, since $\mathfrak D$ is stable under the automorphisms of $E$
and $g\in\Q$.
\item The relative different $\mathfrak D_{E/E^+}$ is trivial: by transitivity $\mathfrak D=\mathfrak D_{E/E^+}\cdot\mathfrak
D_{E^+}\OO_E$, and $\mathfrak D=\mathfrak D_{E^+}\OO_E\ne0$. In particular it is prime to every $g$.
\item If $B\subseteq A$ and $B$ is stable under conjugation, then $B\subseteq\mathrm{tw}_s(A,B)\subseteq A$: for $x\in B$,
$\bar x\in B$ and so $x-(-1)^s\bar x\in B$.
\item If $A\ne0$ is stable under conjugation and $\mathfrak D_{E/E^+}$ is prime to $g$, then
$[A:\mathrm{tw}_s(A,gA)]=[\mathrm{tw}_s(A,gA):gA]=g^{\deg}$ (\lref{Formalism.FourierDecayAlt.P.twistedLatticeIndex}).
\item $P^{*g}=\mathrm{tw}_{s+1}\bigl(\mathfrak a^\vee,\mathfrak a(g,0,-\nu-1)\bigr)$: the inclusion $\subseteq$ is
\lref{Formalism.FourierDecayAlt.P.dual_lattice_subset}, and $\supseteq$ is
\lref{Formalism.FourierDecayAlt.P.traceDualGe}, which uses that $E$ is unramified.
\item $\sqrt{N(\mathfrak a(g,k,\nu))}=(g^k\mathrm{rd}^\nu)^{\deg}$: indeed $N(\mathfrak D)=|d_E|=\mathrm{rd}^{2\deg}$ by (a),
$N(g^k\OO_E)=g^{2k\deg}$, and the norm is multiplicative.
\item For a fractional ideal $B\ne0$, $c\ne0$ and a $\Z$-module $L\supseteq B$ with $[L:B]$ finite, $\Phi_c(L)$ is discrete, spans
$E_\R$ and has covolume $|c|^{2\deg}N(B)2^{-\deg}\sqrt{|d_E|}/[L:B]$ (\lref{Formalism.FourierDecayAlt.P.covolumeIndex}).
\end{enumerate}
Now (1) is (d) and (e) for $A=\mathfrak a(g,1,\nu)$, $gA=\mathfrak a(g,2,\nu)$. (2): by (f) and (d), $P^{*g}\subseteq\mathfrak a^\vee$,
and by (f) and (e) for $A=\mathfrak a^\vee$, $gA=\mathfrak a(g,0,-\nu-1)$, the index is $g^{\deg}$. (3) is (d) and (e) for
$A=\mathfrak a^\sharp$, $gA=\mathfrak a(g,0,-\nu)$. (4) is (g) for $\mathfrak a(g,-1,-\nu-1)$. (5): apply (h) with
$B=\mathfrak a(g,2,\nu)$, $L=P$ and $c=\mathrm{rd}^{-\nu}\ne0$; by (1), $[P:B]=g^{\deg}$, and by (g) and (a),
$N(B)=g^{4\deg}\mathrm{rd}^{2\nu\deg}$ and $\sqrt{|d_E|}=\mathrm{rd}^{\deg}$. The covolume is
\[
\mathrm{rd}^{-2\nu\deg}\cdot g^{4\deg}\mathrm{rd}^{2\nu\deg}\cdot2^{-\deg}\cdot\mathrm{rd}^{\deg}\cdot g^{-\deg}=(\mathrm{rd}/2)^{\deg}g^{3\deg}.
\qedhere
\]
\end{proof}

\subsubsection{Sizes}\label{sec:pf:thmS:sizes}

For a global size datum $G$ we write $k=G.k$, $r=G.r$, $m=G.m=2r+2$, $f=G.f$, $N=10m$, $\Delta=G.\Delta$, $I=\Delta/2-\frac52$ (the
length $\mathrm{len}_I$), $g=g_1g_2$ for its type primes and $\kappa$ for its angle; the hypothesis $\mathrm{Hyp}_G(b,C_\Delta,C_b)$
contains in particular $m\le f\le k$, $(kf\log b)^{C_\Delta}\le\Delta$, $k^{C_b}\le b_{w,1}\le b^2$ for the node values at the
places $w$, $b>1$, $g_1,g_2\ge(kb)^8$, and $v\ge1$, so that the angle $\kappa=\pi/v$ satisfies $0<\kappa\le\pi$.

\begin{theorem}[Size bounds]\label{pf:thmS:sizesmain}
Let $C_A\ge1$, $C_0\ge1$, $C_G\ge0$, $C_\Delta>4$, $c_T,C_b>0$ with $c_TC_b>1$, and
$c_{\exp}=\min\{(C_\Delta-4)/2,(c_TC_b-1)/4\}$. There is $k_0$ with the following property. Let $G$ be a global size datum and
let $b,C',C_f,\Theta\in\R$, word data $d_e,d'_e$ of length $k$, $T,T'\in\N$, $d_u,d_{u'}\in\R$ and $l\in\N$ satisfy: $k\ge k_0$;
$\mathrm{Hyp}_G(b,C_\Delta,C_b)$; $1\le l\le k-2$ and $d_e(l)\ne d'_e(l)$; $b^{c_Tm}\le T$ and $b^{c_Tm}\le T'$; $8C_G\le C'$
and $C'm^5f\log b\le k$; the mass ratios, the lower bound and the column count of the size datum $G.\mathrm{size}(d_e,T,d_u)$,
the column count of $G.\mathrm{size}(d'_e,T',d_{u'})$, the choice of $T$ for both with the constants $C_0$ and
$\exp(C_Gm^5f\log b)$, and the lower bound $\frac12$ for the formal degrees at every place; $\Theta>0$, and $\Theta^{-f}$ is the
largest element of the threshold set $\mathrm{Thr}_G(C_A,d_e,d'_e,T,T',l)$. Then $k^{c_{\exp}}\le\Theta$; and if moreover
$16\le c_{\exp}C_f$ and $C_f\log b\le f\log k$, then $r(C_A/\Theta)^f\le1$ and $C_A(r+1)(C_A/\Theta)^f\le b^{-7}$.
\formallabel{Formalism.FourierDecayAlt.Final.sz_sizesMain}
\end{theorem}

\begin{proof}
Put $C_t=2C_0e^{1/4}\ge1$; $c_{\exp}>0$ since $C_\Delta>4$ and $c_TC_b>1$. For $\Delta,N,I$ and $TT>0$ let
\[
\mu(\Delta,N,I,TT)=\min\Bigl\{\frac{\sqrt{\Delta/(kN^3)}}{160C_tC_A},\ \frac{I}{160\pi C_tC_A^2\sqrt k\,N},\
\Bigl(\frac{\sqrt{TT}}{160\pi C_tC_A\sqrt k}\Bigr)^{1/2}\Bigr\}.
\]
The power bound \lref{Formalism.FourierDecayAlt.P.theta_power_bound} provides $k_1$ such that for $k\ge k_1$, $2\le m\le f\le
k$, $N=10m$, $I=\Delta/2-\frac52$, $k^{C_b/2}\le b$, $(kf\log b)^{C_\Delta}\le\Delta$ and $(b^{c_Tm})^2\le TT$ one has
$\pi C_AN/I\le1$ and $k^{c_{\exp}}\le\mu(\Delta,N,I,TT)$. The final smallness
\lref{Formalism.FourierDecayAlt.P.final_smallness} provides $k_2$ such that for $k\ge k_2$, $2r+2\le f$, $b\ge1$,
$16\le c_{\exp}C_f$, $C_f\log b\le f\log k$ and $k^{c_{\exp}}\le\Theta$ one has $r(C_A/\Theta)^f\le1$ and
$C_A(r+1)(C_A/\Theta)^f\le b^{-7}$. Put $k_0=\max(7,k_1,k_2)$ and let the data be given.

\emph{Elementary consequences.} We have $k\ge7$, $m\ge2$ and $f\ge m\ge2$. Since $k^{C_b}\le b^2$, $k^{C_b/2}\le b$; as $b>1$,
$T,T'\ge b^{c_Tm}>1$, so $T,T'\ge2$ and $(b^{c_Tm})^2\le TT'$. The power bound applies with $TT=TT'$: $\pi C_AN/I\le1$ and
$k^{c_{\exp}}\le\mu(\Delta,N,I,TT')$. Since $k^{c_{\exp}}\ge1$ and $C_tC_A\ge1$, the first entry of $\mu$ gives
$\Delta\ge160^2kN^3$; hence $N^3\le\Delta/2$, $\Delta>0$ and $I>0$. Also $g_1,g_2\ge(kb)^8\ge4$, so $g\ge16\ge2\pi$.

\emph{Positivity.} By the interlacing of the parameters $\lambda$ and $u(d_e)$ at every place
(\lref{Formalism.FourierDecayAlt.P.interlacing}) and the strict monotonicity of $\lambda$
(\lref{Formalism.FourierDecayAlt.P.lamStrictAnti}), the Vandermonde products of $u(d_e)$, $u(d'_e)$ and $\lambda$ are positive;
with the lower bound for the formal degrees, the formal-degree quantities $V(w)$ and $d_\sigma(w)$ are positive at every place $w$.

\emph{Critical heights.} Corollary~A (\lref{Formalism.FourierDecayAlt.P.corollary_a_of_proposition_a}: under the mass
ratios, the lower bound, the column count, the choice of $T$ with constants $C_0$, $\exp(C_Gm^5fL)$, $8C_G\le C'$,
$C'm^5fL\le k$ and $8f2^{-k}\le\log2$, the product $\prod_wT_{\mathrm{crit}}(w)/T$ of the critical heights lies between
$(2C_0)^{-1/k}$ and $(2C_0)^{1/k}e^{1/4}\le C_t$), applied with $L=\log b$ to both size data (note $8f2^{-k}\le\log2$ for
$f\le k$, $k\ge7$), gives $\prod_wT_{\mathrm{crit}}(w)/T\le C_t$ and $\prod_wT'_{\mathrm{crit}}(w)/T'\le C_t$. Put
$t_w=\bigl(T_{\mathrm{crit}}(w)T'_{\mathrm{crit}}(w)/(TT')\bigr)^{1/2}\ge0$; then $(\prod_wt_w)^2\le C_t^2$, so $\prod_wt_w\le C_t$.

\emph{The $\beta$-function.} The boundary bound \lref{Formalism.FourierDecayAlt.P.beta_bound_on_boundary} with its constant
$A_0\in(0,160]$ (\lref{Formalism.FourierDecayAlt.P.a_zero_le_one_sixty}) gives, for $Y$ in the cut interval at $l$ with
$\frac52\le v_l+Y\le D\le\Delta/2$: $0\le g^m\beta_w(Y)\le160\,t_w\sqrt k\sqrt{D/\Delta}$. Taking suprema, the long-cut supremum
$B_1(w)$ (with $D=N^3$) and the supremum $B(w)$ over the whole interval (with $D=\Delta/2$) satisfy $0\le B_1(w)\le160t_w\sqrt
k\sqrt{N^3/\Delta}$ and $0\le B(w)\le160t_w\sqrt k$.

\emph{Thresholds.} The threshold products \lref{Formalism.FourierDecayAlt.P.threshold_products} (for $C=C_A$, $A_0=160$,
$C_t$, $\kappa\in(0,\pi]$, $N,I,g,\Delta,TT'$, using $\pi C_AN/I\le1$, $\prod_wt_w\le C_t$ and the bounds for $B_1,B$) bound the
four kinds of elements of the threshold set: the products of the $x$-thresholds at $i=1$ by
$C_t(160C_A)^f(kN^3/\Delta)^{f/2}$, those at $i\ge2$ by $C_t(160C_A\sqrt k)^f(\pi C_AN/I)^{2f}$, the products of the
$\lambda$-thresholds with $s\ge2$ by $C_t(160C_A\sqrt k)^f(\pi C_AN/I)^f$, and the products of the $\lambda'$-thresholds by
$\sqrt{C_t}\,(160\pi C_A\sqrt k/\sqrt{TT'})^{f/2}$. The largest element $\Theta^{-f}$ is of one of these kinds, so it is bounded
by the corresponding quantity; the lower bound \lref{Formalism.FourierDecayAlt.P.theta_lower_bound} (if $\Theta^{-f}$ is
at most one of these four quantities, then $\mu(\Delta,N,I,TT')\le\Theta$) gives $\mu(\Delta,N,I,TT')\le\Theta$, hence
$k^{c_{\exp}}\le\Theta$. Finally, if $16\le c_{\exp}C_f$ and $C_f\log b\le f\log k$, the final smallness applies, since
$k\ge k_2$, $2r+2=m\le f$ and $b\ge1$.
\end{proof}

\subsubsection{The theta step}\label{sec:pf:thmS:theta}

A \emph{local datum} $(K,R,\psi)$ consists of a Hausdorff commutative topological ring $K$ which is Artinian and reduced, an open
compact subring $R$ containing a unit $\varpi$ of $K$ with $\varpi^mx\in R$ for every $x\in K$ and some $m$, such that every
$x\in K$ has $x+r\in K^\times$ for some $r\in R$, and an additive character $\psi$ with $\overline{\psi(x)}=\psi(-x)$ and
$\psi(aR)=1\iff a\in R$. For $x\in K$ let $|x|_R=[xR:xR\cap R]/[R:R\cap xR]$. A \emph{point} $h$ consists of a list of block
sizes, of total size $|h|$, and units of $K$; $\mathrm{pt}(h)$ is the associated matrix, $\mathrm{cs}(h)=|\det\mathrm{pt}(h)|_R^{-1}$,
$h\oplus(s,t)$ is obtained by appending a block of size $s$ with the unit $t$, and $h\in\GL(R)$ means that $\mathrm{pt}(h)$ and its
inverse have entries in $R$. The local kernel $N_0(h,h')=K^{\mathrm{loc}}(h,h')$ is the finite sum over pairs of cosets of the
development, and for $s\ge1$, units $t,t'$,
\[
N_s(h,h',t,t')=|tt'|_R^{s(s-1)/2}\,N_0\bigl(h\oplus(s,t),h'\oplus(s,t')\bigr).
\]
For $s\ge1$, a finite set $F$ of points and $\Theta\colon K\to\{\text{kernels on points}\}$, the \emph{theta properties}
$\mathrm{TP}_s(F,\Theta)$ are:
(T1) $\Theta(x+y)=\Theta(x)$ for $y\in R$;
(T2) $\Theta(x)$ is PSD on $F$;
(T3) there is an open compact $R$-submodule $L$ of $K$ with $\Theta(x)=0$ on $F\times F$ for $x\notin L$;
(T4) $\Theta(t)-N_s(\cdot,\cdot,t,t)$ is PSD on $F$ for every unit $t$;
(T5) $M_s-\sum_{x\in K/R}\Theta(x)$ is PSD on $F$, where $M_s(h,h')=\mathrm{cs}(h)^s\mathrm{cs}(h')^sN_0(h,h')$;
(T6) $\Theta(x)(h,h')=0$ if $|\det\mathrm{pt}(h)|_R\ne|\det\mathrm{pt}(h')|_R$;
(T7) if all points of $F$ lie in $\GL(R)$, then $\one_R(x)N_0-\Theta(x)$ is PSD on $F$ for every $x$.
\emph{Local theta majorants exist for blocks of size $s$} if for every $m$ and every finite set $F$ of points of total size $m$
there is $\Theta$ with $\mathrm{TP}_s(F,\Theta)$. The block-kernel statement (A2) asserts that for every local datum and $s\ge1$
the kernels $((h,t),(h',t'))\mapsto N_s(h,h',t,t')$ are PSD on finite sets of pairs (point of a fixed total size, unit) and
satisfy a support condition; the single-block Parseval identity asserts that for points $h,h'$ of equal total size and unit
representatives $c(x)$ of the classes $x\in K/R$, $x\mapsto N_1(h,h',c(x),c(x))$ has finite support and sum
$\mathrm{cs}(h)\mathrm{cs}(h')N_0(h,h')$.

\begin{theorem}[The theta step]\label{pf:thmS:thetastep}
Assume (A2) and the single-block Parseval identity. Let $(K,R,\psi)$ be a local datum in which every non-zero integer is a unit,
and $s\ge1$. If local theta majorants exist for blocks of size $s$, they exist for blocks of size $s+1$.
\formallabel{Formalism.FourierDecayAlt.P.lattice_average_induction}
\end{theorem}

\begin{proof}
Let $F$ be a finite set of points of total size $m$, and write $M=M_{s+1}$. We use three results of the development.
\begin{enumerate}[label=(\alph*)]
\item (Cut module.) There is an open compact $R$-submodule $L_c\supseteq R$ of $K$ such that $N_{s+1}(h,h',t,t)=0$ for
$h,h'\in F$ and units $t\notin L_c$, with $L_c=R$ if all points of $F$ lie in $\GL(R)$
(\lref{Formalism.FourierDecayAlt.P.finite_family_truncation_module}, from the support statement
\lref{Formalism.FourierDecayAlt.P.added_block_kernel_support}: $N_s(h,h',t,t')\ne0$ implies $\det\mathrm{pt}(h)\,t\in R$
and $\det\mathrm{pt}(h')\,t'\in R$).
\item (Exhaustion; \lref{Formalism.FourierDecayAlt.P.rootStableExhaustion}.) There are open compact sets
$\Omega_1\subseteq\Omega_2\subseteq\cdots$ of units, each stable under multiplication by the $(s+1)$-th roots of unity, whose union
contains every unit.
\item (Majorants on $\Omega$; \lref{Formalism.FourierDecayAlt.P.theta_omega_bound}.) Under (A2), the Parseval identity and
the existence of majorants for size $s$, for $L_c$ as in (a) and each such $\Omega$ there is $\Theta_\Omega$ with (T1), (T2),
(T5), (T6) for $s+1$, vanishing outside $L_c$, and with (T4) for the units in $\Omega$. This is the core of the argument; its
proof uses the existence for size $s$, $p$-adic stationary phase, Gauss sums of quadratic forms and averages over lattices.
\end{enumerate}
Let $\Theta_j=\Theta_{\Omega_j}$.

\emph{Compactness.} The quotient $L_c/R$ is finite, since $L_c$ is compact and covered by the disjoint open cosets of $R$. Let $S$
be a set of representatives. As $\Theta_j$ is $R$-periodic and vanishes outside $L_c$, the sum in (T5) is the finite sum over
$S$. For every $x$ the kernel $M-\Theta_j(x)$ is PSD: if $x\equiv y\in S$ it is the sum of $M-\sum_{y'\in S}\Theta_j(y')$ and of
the PSD kernels $\Theta_j(y')$, $y'\ne y$; if $x\notin L_c$, then $\Theta_j(x)=0$ and $M$ itself is such a sum. If $k$ and $M-k$
are PSD then $0\le k(h,h)\le M(h,h)$ and $|k(h,h')|^2\le k(h,h)k(h',h')$ (a $2\times2$ principal minor), so
$|k(h,h')|^2\le M(h,h)M(h',h')$. Hence the finitely many sequences $j\mapsto\Theta_j(y)(h,h')$, $y\in S$, $h,h'\in F$, are
bounded, and by the Bolzano--Weierstrass theorem there is a strictly increasing $\phi$ along which all of them converge. Let
$\Theta(x)$ be the limit at the representative of $x$ if $x\in L_c$, and $\Theta(x)=0$ otherwise; then
$\Theta_{\phi(j)}(x)\to\Theta(x)$ on $F\times F$ for every $x\in K$.

\emph{Limit.} (T1) holds by construction, and (T3) with $L=L_c$. PSD kernels on $F$ form a closed set; this gives (T2), and also
(T4): a unit $t$ lies in $\Omega_j$ for all large $j$, so $\Theta_{\phi(j)}(t)-N_{s+1}(\cdot,\cdot,t,t)$ is PSD for large $j$. The
PSD kernels $M-\sum_{y\in S}\Theta_{\phi(j)}(y)$ converge to $M-\sum_{y\in S}\Theta(y)$, the kernel of (T5); and (T6) holds since
the relevant entries vanish along the sequence. Finally, if all points of $F$ lie in $\GL(R)$, then $L_c=R$, and
$\mathrm{cs}(h)=1$ for $h\in F$ (the determinant is a unit of $R$, and multiplication by a unit of $R$ preserves $R$), so
$M=N_0$. For $x\in R$ the kernel $N_0-\Theta(x)=M-\Theta(x)$ is PSD, as a limit of the PSD kernels $M-\Theta_{\phi(j)}(x)$; for
$x\notin R=L_c$ both $\one_R(x)N_0$ and $\Theta(x)$ vanish. This is (T7).
\end{proof}

\subsubsection{The local and global inputs}\label{sec:pf:thmS:cores}

A kernel $k$ on a finite set $F$ is PSD if $\sum_{x,y\in F}a_x\overline{a_y}k(x,y)$ is real and non-negative for all $a\in\C^F$.
A \emph{local system} $S=(N,\Theta,B,L,L')$ on a set $\alpha$ over a field $E$ consists of kernels $N(t,t')$ and $\Theta(t)$ on
$\alpha$ ($t,t'\in E$), a kernel $B$ and additive subgroups $L,L'$ of $E$. It is \emph{valid} on a finite set $F$ if:
$N(t,t')(h,h')\ne0$ with $t,t'\ne0$ and $h,h'\in F$ implies $t^{-1}-t'^{-1}\in L'$; $N(t,t)$ is PSD for $t\ne0$; $\Theta(t)$
depends only on $t$ modulo $L$; $\Theta(t)$ is PSD; $\Theta(t)-N(t,t)$ is PSD for $t\ne0$; and $B-\sum_{t\in R}\Theta(t)$ is PSD
for every finite set $R$ of elements pairwise incongruent modulo $L$. Products of local systems are taken entrywise, with
$L=\bigcap_pL_p$ and $L'=\bigcap_pL'_p$; a rescaling by $a\colon\alpha\to\R$ multiplies $N$, $\Theta$, $B$ by $a(h)a(h')$; a phase
twist by $d\colon E\to\C$ multiplies $N(t,t')$ by $d(t)\overline{d(t')}$. The block system of a TB-datum at $(i,s)$ has kernels
$K_i(\thmSjoin(h,s,t),\thmSjoin(h',s,t'))$, majorants $\Theta_{i,s}$, base kernel $(1+e_i)K_{i-s}$ and the lattices
$\mathrm{La}_{i-s+1,s}$, $\mathrm{Lap}_{i-s+1,s}$. For a local datum $(K,R,\psi)$, a field embedding $\sigma\colon E\to K$, a unit
$u$, a map $\eta$ to points and a family $\Theta$, the \emph{pull-back system} of size $s$ has kernels
$N_s(\eta h,\eta h',u\sigma(t),u\sigma(t'))$, majorants $\Theta(u\sigma(t))(\eta h,\eta h')$, base kernel $M_s(\eta h,\eta h')$
and lattices $\sigma^{-1}(u^{-1}R)$, $\sigma^{-1}(uR)$; the \emph{top system} (with a shift $\tau_0\in K$) has kernels
$\delta_{tt'}\Theta(u\sigma(t)-\tau_0)$, the same majorants, base kernel $M_1$ and lattices $\sigma^{-1}(u^{-1}R)$, $0$; an
\emph{outside system} has given kernels $N$, majorants $\one_L(t)$ and base kernel $1$.

\begin{theorem}[The local inputs]\label{pf:thmS:localcores}
The following statements hold.
\begin{enumerate}[label=(LI\arabic*)]
\item\label{pf:thmS:LI1} The three notions of positive semi-definiteness of kernels on finite sets used in the development (a
positive semi-definite matrix; a hermitian kernel with non-negative real part of the quadratic form; a kernel with real
non-negative quadratic form) coincide.
\item\label{pf:thmS:LI2} (Schur product.) Finite products of PSD kernels are PSD.
\item\label{pf:thmS:LI3} (Product domination.) If $k_p$ and $K_p-k_p$ are PSD for all $p$ in a finite set, then
$\prod_pK_p-\prod_pk_p$ is PSD.
\item\label{pf:thmS:LI4} (Diagonal scaling and phase.) If $k$ is PSD, so are $a(x)a(y)k(x,y)$ for real $a$ and
$d(x)\overline{d(y)}k(x,y)$ for complex $d$.
\item\label{pf:thmS:LI5} (Block diagonal.) If $k$ is PSD on every fibre of $c\colon F\to\beta$, then $\one[c(x)=c(y)]\,k(x,y)$ is
PSD on $F$.
\item\label{pf:thmS:LI6} (Operations on systems.) Products of systems valid on $F$ are valid on $F$; rescalings of valid systems are
valid; phase twists of valid systems by $d$ with $|d|=1$ are valid; an outside system is valid if its kernels satisfy the support
condition and $N(t,t)$, $\one_L(t)-N(t,t)$ are PSD for $t\ne0$; and if $k_p$ and $\one_{L_p}(t)-k_p$ are PSD for all $p$, then
$\one_{\bigcap_pL_p}(t)-\prod_pk_p$ is PSD.
\item\label{pf:thmS:LI7} (Pull-back and top systems.) For a family $\Theta$ with $\mathrm{TP}_s(\eta(F),\Theta)$, $s\ge1$, and points
$\eta(h)$ of equal total size, the pull-back system is valid on $F$; for $\Theta$ with $\mathrm{TP}_1(\eta(F),\Theta)$ the top
system is valid on $F$; and if moreover all $\eta(h)$ lie in $\GL(R)$ (and $\tau_0\in R$ for the top system), the kernels
$\one_L(t)-N(t,t)$ of these systems are PSD.
\item\label{pf:thmS:LI8} (Global from local.) If the block system of a TB-datum at $(i,s)$ agrees on $Y$ with a rescaled product of
local systems valid on $Y$ (kernels, base kernel and $L$, with $\bigcap_pL'_p$ contained in its $L'$), then it is valid on $Y$ for
suitable majorants.
\item\label{pf:thmS:LI9} (Uniform majorants.) If for every finite subset $F$ of a set $A$ some majorant makes a system valid on
$F$, then one majorant, vanishing outside $A\times A$ and PSD everywhere, makes it valid on every finite subset of $A$.
\item\label{pf:thmS:LI10} (Axioms.) If all block systems of a TB-datum are valid on finite sets of the right level, $0\le e_i\le1$,
the majorants vanish off that level, and $K_{r+1}(\thmSjoin(h,1,t),\thmSjoin(h',1,t'))=0$ for $h,h'$ of level $r$ and $t\ne t'$
non-zero, then the datum satisfies the kernel-support and theta-domination axioms; if only the existence of majorants for each
finite set is known, there are majorants for which these axioms hold and all $\Theta_{i,s}(t)$ are PSD.
\item\label{pf:thmS:LI11} (Kernels of local type.) If $K(x,y)=a(x)a(y)\prod_pk_p(x,y)$ on $Y$ with all $k_p$ PSD, then $K$ is PSD;
and the local kernels $N_0(\eta x,\eta y)$ of a local datum are PSD on every finite set.
\item\label{pf:thmS:LI12} (Domination.) The textual domination of a TB-datum implies its minimal domination; and if all block
systems at all towers are rescaled products of valid local systems, then the textual domination holds at every tower and the
datum satisfies the domination axiom.
\item\label{pf:thmS:LI13} (Local theta kernels.) For every local datum: (A2); the Gram property of $N_0$ (it is PSD on finite sets of
points); $N_0(h,h')=1$ for points of equal size in $\GL(R)$ when $\psi(R)=1$; the single-block Parseval identity; the existence of
a single-block family $\Theta$ with $\Theta(t)=N_1(\cdot,\cdot,t,t)$ on units, sum over $K/R$ equal to $M_1$, and
$\mathrm{TP}_1(F,\Theta)$ for all $F$ of fixed size; the support statement of \cref{pf:thmS:thetastep}(a); the block-diagonality
of $N_0$ and $N_s$ in $|\det\mathrm{pt}|_R$; and the assembly of the local lattice-average lemma from these statements and the
theta step.
\end{enumerate}
\formallabel{Formalism.FourierDecayAlt.P.HilbertSchmidtNorm.localCores}
\end{theorem}

\begin{proof}
\ref{pf:thmS:LI1}: a positive semi-definite matrix is hermitian with real non-negative quadratic form; conversely, if
$Q(a)=\sum a_x\overline{a_y}k(x,y)$ is real for all $a$, then by polarisation $k(y,x)=\overline{k(x,y)}$, and for a hermitian
kernel $Q(a)$ is real, so the three conditions agree. \ref{pf:thmS:LI2} is the Schur product theorem \cite[Thm.~7.5.3]{HornJohnson2013}
(by induction; the empty product is the constant kernel $1$, which is PSD). \ref{pf:thmS:LI3}: by induction, writing
\[
\prod_{p\in P}K_p-\prod_{p\in P}k_p=(K_q-k_q)\prod_{p\ne q}K_p+k_q\Bigl(\prod_{p\ne q}K_p-\prod_{p\ne q}k_p\Bigr),
\]
each term is a product of PSD kernels ($K_p=k_p+(K_p-k_p)$ is PSD). \ref{pf:thmS:LI4}: substitute $a_xa(x)$, resp. $a_xd(x)$, for
$a_x$ in the quadratic form. \ref{pf:thmS:LI5}: the quadratic form is the sum over $b\in\beta$ of the quadratic forms on the fibres.
The first part of \ref{pf:thmS:LI11} follows from \ref{pf:thmS:LI2} and \ref{pf:thmS:LI4}, and the second from the Gram property
in \ref{pf:thmS:LI13}: the quadratic form of $N_0(\eta x,\eta y)$ at $a$ is that of $N_0$ on $\eta(Y)$ at
$b_h=\sum_{\eta x=h}a_x$. The first part of \ref{pf:thmS:LI12} is a restriction of the quantifiers: the minimal domination
consists of the conditions of the textual one for the relevant $t$.

The remaining statements are theorems of the development: \ref{pf:thmS:LI6} is
\lref{Formalism.FourierDecayAlt.Final.jx_productLemma}, \lref{Formalism.FourierDecayAlt.Final.jx_scaleOk}, \lref{Formalism.FourierDecayAlt.Final.jx_phaseOk}, \lref{Formalism.FourierDecayAlt.Final.jx_outsideSysOk}, and \lref{Formalism.FourierDecayAlt.Final.jx_outsideDomination};
\ref{pf:thmS:LI7} is
\lref{Formalism.FourierDecayAlt.Final.jx_localPullback}, \lref{Formalism.FourierDecayAlt.Final.jx_localUnitHeads}, \lref{Formalism.FourierDecayAlt.Final.jx_topPullback}, and \lref{Formalism.FourierDecayAlt.Final.jx_topUnitHeads};
\ref{pf:thmS:LI8} is \lref{Formalism.FourierDecayAlt.Final.jx_globalOfLocal}; \ref{pf:thmS:LI9} is
\lref{Formalism.FourierDecayAlt.Final.jx_uniformMajorantOn}; \ref{pf:thmS:LI10} is
\lref{Formalism.FourierDecayAlt.Final.jx_axiomsFromSystems} and \lref{Formalism.FourierDecayAlt.Final.jx_axiomsFromFamilies}; the
second part of \ref{pf:thmS:LI12} is \lref{Formalism.FourierDecayAlt.Final.tc_domAfterOfLocal}; and the statements of
\ref{pf:thmS:LI13} are, in this order,
\lref{Formalism.FourierDecayAlt.P.second_gram_and_vanishing}, \lref{Formalism.FourierDecayAlt.P.kernel_gram_property}, \lref{Formalism.FourierDecayAlt.P.zero_kernel_all_ones}, \lref{Formalism.FourierDecayAlt.P.first_kernel_parseval}, \lref{Formalism.FourierDecayAlt.P.first_theta_kernel_equality}, \lref{Formalism.FourierDecayAlt.P.first_theta_equals_first_kernel}, \lref{Formalism.FourierDecayAlt.P.added_block_kernel_support}, \lref{Formalism.FourierDecayAlt.P.added_block_kernel_block_diagonal}, and \lref{Formalism.FourierDecayAlt.P.lattice_average_lemma_from_parts}
(the fifth and sixth giving the single-block family, the first with its periodicity, value on units and sum, the second with the
theta properties). Each of them states exactly the corresponding item.
\end{proof}

\begin{theorem}[The global inputs]\label{pf:thmS:globalcores}
The following statements hold.
\begin{enumerate}[label=(GI\arabic*)]
\item\label{pf:thmS:GI1} Theorem~A in the form of \cref{pf:thmS:thmA}, and the recursion with densities \cref{pf:thmS:recdens}.
\item\label{pf:thmS:GI2} (Scaling at the top.) If $\lambda\ge0$, $e_{r+1}=0$ and $\Theta$ is a minimal domination of a TB-datum $D$
at a tower, then $\Theta$ with the top majorants multiplied by $\lambda$ is a minimal domination of $D$ with $K_{r+1}$
multiplied by $\lambda$ and $e_{r+1}$ replaced by $\lambda-1$.
\item\label{pf:thmS:GI3} (Sum bound with maximum.) If $k$ and $s_dk-m$ are PSD, so is $\max(s_d,1)k-m$.
\item\label{pf:thmS:GI4} (Scaled recursion hypotheses.) Statement \ref{pf:thmS:M7}.
\item\label{pf:thmS:GI5} (Densities.) If $n\ge2$, $\ell_p\le q_p$ with $\ell_p$ prime and at most $f$ indices of $S$ share a prime,
then $\sum_{p\in S}q_p^{-n}\le3f2^{-n}$; if $18f\le2^n$, then $\exp(12f2^{-n})\le2$; if $2r+2+f+5\le k$, then
$18f\le2^{k-2r-2}$; and statement \ref{pf:thmS:M6}.
\item\label{pf:thmS:GI6} (Poisson summation.) Statement \ref{pf:thmS:M5} and its two inputs: for a lattice $\Lambda$ in a real
inner product space $V$ with coordinates $\pi_v$, and a continuous $\varphi\ge0$ on $\Lambda+x$ with
$|\varphi(c)|\le C\prod_v(1+|\pi_v(c)|)^{-3}$ and summable Fourier transform on the dual lattice, the Poisson summation formula holds
at $x$ and $\sum_{l\in S}\varphi(l+x)$ is bounded by the Poisson bound of $(\Lambda,\varphi)$; and its form for lattices given by
embeddings and dual maps.
\item\label{pf:thmS:GI7} (Box sums.) There is $C$ such that for every totally complex number field $E$, non-zero fractional ideal
$\mathfrak b$, $x_w>0$, $\Theta>2$ and $K_x\ge3$ with $\Theta^{d}=\sqrt{N\mathfrak b}/\prod_wx_w$ ($d$ the number of infinite
places $w$) and every finite set $S$ of non-zero elements of $\mathfrak b$:
$\sum_{l\in S}\prod_w\min(1,(x_w/|l|_w)^{K_x})\le(C\Theta^{-(K_x-5/2)})^{d}$.
\item\label{pf:thmS:GI8} (Lattices and ideals.) \Cref{pf:thmS:latfacts}; for $\nu_{r,i,s}=2r+3-2i-s$:
$\mathfrak a^\vee$ at $\nu_{r,i,s}$ equals $\mathfrak a^\sharp$ at $\nu_{r,i,s-1}$ and $\mathfrak a(g,-1,-(2r+4-2i-s))$, and
$\nu_{r,i,s}=2(r+1)-i-(i-1+s)$; the norm formula $\sqrt{N(\mathfrak a(g,k,\nu))}=(g^k\mathrm{rd}^\nu)^{\deg}$; and the degree facts
$[E:\Q]=2\deg$, $E$ has $\deg>0$ infinite places.
\item\label{pf:thmS:GI9} (Thresholds.) The frame-product identities expressing the thresholds $\lambda$ and $\lambda'$ through their
normalised versions, and the identity of the $x$-thresholds with the $\lambda$-thresholds at $s=1$.
\item\label{pf:thmS:GI10} (Lattice sums.) The joint lattice-sum bound (from the box sums), the instance sizes (from
\ref{pf:thmS:GI8} and \ref{pf:thmS:GI9}), and statement \ref{pf:thmS:M2}.
\item\label{pf:thmS:GI11} (Sizes.) \Cref{pf:thmS:sizesmain} and statements \ref{pf:thmS:M3} and \ref{pf:thmS:M4}.
\item\label{pf:thmS:GI12} (Masses.) The mass identities \cref{pf:thmS:massid}, in both forms.
\end{enumerate}
\formallabel{Formalism.FourierDecayAlt.P.HilbertSchmidtNorm.globalCores}
\end{theorem}

\begin{proof}
\ref{pf:thmS:GI1}, \ref{pf:thmS:GI11} (first part) and \ref{pf:thmS:GI12} are proved in this section, and so is the first part of
\ref{pf:thmS:GI8}. \ref{pf:thmS:GI2}: each condition of the minimal domination at level $r+1$ (the majorants minus the blocks of
$K_{r+1}$ are PSD, periodicity, and the base kernel $(1+e_{r+1})K_{r+1-s}$ minus sums of majorants is PSD) is multiplied by
$\lambda\ge0$, since $1+(\lambda-1)=\lambda$ and $e_{r+1}=0$; the conditions at the other levels do not change.
\ref{pf:thmS:GI3}: $\max(s_d,1)k-m=(s_dk-m)+(\max(s_d,1)-s_d)k$, a sum of PSD kernels. In \ref{pf:thmS:GI5}: if $18f\le2^n$, then
$12f2^{-n}\le\frac23$ and $e^{2/3}<2$; if $2r+2+f+5\le k$, then $2^{k-2r-2}\ge2^{f+5}\ge32(f+1)\ge18f$; the first statement is
\lref{Formalism.FourierDecayAlt.Final.tc_densSum} and the last \lref{Formalism.FourierDecayAlt.Final.tc_densChain}. In
\ref{pf:thmS:GI8}: $\mathfrak a^\vee$ at $\nu_{r,i,s}$ is $\mathfrak a(g,-1,-(\nu_{r,i,s}+1))$, $\mathfrak a^\sharp$ at
$\nu_{r,i,s-1}$ is $\mathfrak a(g,-1,-\nu_{r,i,s-1})$, and $\nu_{r,i,s-1}=\nu_{r,i,s}+1=2r+4-2i-s=2(r+1)-i-(i-1+s)+1$; the norm
formula is step (g) of the proof of \cref{pf:thmS:latfacts}, and the degree facts are
\lref{Formalism.FourierDecayAlt.P.degreeFacts_statement}.

The remaining statements are theorems of the development: \ref{pf:thmS:GI4} is
\lref{Formalism.FourierDecayAlt.Final.tc_recHypScaleTop}; \ref{pf:thmS:GI6} is
\lref{Formalism.FourierDecayAlt.Final.jx_poissonStepJoint}, \lref{Formalism.FourierDecayAlt.Final.poisson_summation_text}, and \lref{Formalism.FourierDecayAlt.Final.jx_poissonJoint};
\ref{pf:thmS:GI7} is \lref{Formalism.SA.P.smoothedSumBound}; \ref{pf:thmS:GI9} is
\lref{Formalism.FourierDecayAlt.P.frame_product_identities} and \lref{Formalism.FourierDecayAlt.P.lattice_x_eq_lambda_one};
\ref{pf:thmS:GI10} is
\lref{Formalism.FourierDecayAlt.P.lattice_sums_joint}, \lref{Formalism.FourierDecayAlt.P.instance_sizes}, and \lref{Formalism.FourierDecayAlt.Final.jx_instLatticeSums};
and the second part of \ref{pf:thmS:GI11} is
\lref{Formalism.FourierDecayAlt.Final.jx_sizesToRecursion} and \lref{Formalism.FourierDecayAlt.Final.jx_globPositive}. Each of them
states exactly the corresponding item.
\end{proof}

\subsubsection{Row estimates along families and the corank-one property}\label{sec:pf:thmS:loopsproofs}

\begin{proof}[Proof of \cref{pf:thmS:loopbounds}]
Fix the constants $(C_g,c_b,C_u,c_T,C_T,C_0,A,C_1,c_V,C_\Gamma)$.

\emph{Constants.} Under the Siegel mass ratio, \lref{GlobalTheory.T.rowCounts_uniform_constants_families} provides constants
$c_3,c_H,c_S,c_4$ and a threshold $\vartheta_1$ such that for every family $F$ with $s_{C'}(F),s_{C_l}(F)\ge\vartheta_1$ the four
row-count conditions of the admissibility package hold eventually, for all pairs of words: the local height-determinant bound
with constant $3$ for every $C\ge c_3$, the extension-class bound with $c_H$, the matching of archimedean signatures with $c_S$,
and the row bound with $c_4$. Under the Siegel mass ratio, $\mathrm{GH}$ and the pointwise bound,
\lref{GlobalTheory.T.amplitude_bounds_uniform_large_gap_families} provides $C_a\ge0$, $C_2$, $C_h\ge0$ and a threshold
$\vartheta_2$ such that for every family with $s_{C'},s_{C_l}\ge\vartheta_2$ and the mass formulas $\mathrm{MF}$, for every $C_{pr}$
and $\delta>0$, the amplitude bounds with these constants hold eventually. Let $U=(c_3,c_H,c_V,c_S,c_4,C_a,C_2,C_h)$, let $C_*$
be the threshold $c_0(C_T,A,C_0,C_\Gamma,U)$ of the row estimates below, and $\vartheta=\max(40,\vartheta_1,\vartheta_2)$.

\emph{(i).} Let $F$ be a family with $C_\Gamma(F)=C_\Gamma$, $s_{C'}(F),s_{C_l}(F)\ge\vartheta$ and $\mathrm{MF}$. The raw-mass and
genus-mass identity holds for every datum and word (\lref{GlobalTheory.T.rows_raw_mass_and_genus_mass}). Hence, as
$s_{C_l}(F)\ge40$, \lref{GlobalTheory.CritLevelGap.along_family} applies (with the formal-degree identity, a component of
$\mathrm{GH}$): it gives the existence of the family rows, the three critical-level conditions (F6)--(F8) of the admissibility
package, the bound for the quotients of Weyl dimensions and Vandermonde products, and statement (i).

\emph{(ii).} Let $X,Y$ be as in (ii). The family is admissible with the constants $U$: the components of admissibility are the
four row-count conditions of the first step, the small-cell volume bound (a hypothesis of (ii)) and, for every $C_{pr}$, the
constants $C_h,C_a$ with the amplitude bounds of the first step for all $\delta>0$. By
\lref{GlobalTheory.T.admissible_family_satisfies_row_estimates} (for an admissible family with the conditions of (i), the tie
conditions, injectivity of terms and the $\beta$-lattice property, and every $C\ge c_0(C_T,A,C_0,C_\Gamma(F),U)$, eventually and
under $\mathrm{CH}$ with $A$ the three row estimates hold with constant $C$), statement (ii) follows, since $C_\Gamma(F)=C_\Gamma$.
\end{proof}

\begin{proof}[Proof of \cref{pf:thmS:betauniq}]
Let $V$ be the hermitian space of $D$, of dimension $k+2r+1$ ($r=r_S$), $U$ the $E$-span of the lattice generated by $X(F)$,
$X(F')$ and the lifts of the flags, $H=U^\perp$, and $U''$ the analogous space for $(F,F'')$. For an isometry $\beta$ let
$\mathrm{Th}(\beta)$ be the space of $x\in H$ for which there is $y\in U''^\perp$ with $x-y\in X(F)$ and $\beta x-y\in X(F'')$; the
corank is $s_0(\beta)=k-\dim\mathrm{Th}(\beta)$.

(a) Since $(F,F')$ and $(F,F'')$ are terms, their Gram values are non-zero, and
\lref{PI.UnitaryGroupMassFormula.BetaLUniq.gram_matrix_nondegenerate} gives $\dim U=\dim U''=2r+2$ and
$U\cap U^\perp=U''\cap U''^\perp=0$ (the $2r+2$ flag vectors spanning $U$ are linearly independent, and no non-zero vector of $U$
is orthogonal to $U$).

(b) $\beta_1$ and $\beta_2$ agree on $H$ modulo $X(F'')$. Indeed $\dim H=\dim V-\dim U=k-1$ by (a), and
$\mathrm{Th}(\beta_\iota)\subseteq H$ has dimension $k-s_0(\beta_\iota)=k-1$, so $\mathrm{Th}(\beta_\iota)=H$. For $x\in H$ choose
$y_\iota\in U''^\perp$ with $x-y_\iota\in X(F)$ and $\beta_\iota x-y_\iota\in X(F'')$. Then $y_1-y_2=(x-y_2)-(x-y_1)\in
X(F)\subseteq U''$ and $y_1-y_2\in U''^\perp$, so $y_1=y_2$ by (a), and $\beta_2x-\beta_1x=(\beta_2x-y_2)-(\beta_1x-y_1)\in X(F'')$.

(c) Every isotropic vector orthogonal to $X^+(F')=X(F')+Ez(F')$ lies in $X(F')$
(\lref{PI.UnitaryGroupMassFormula.BetaLUniq.weight_space_anisotropic}).

(d) If $v\perp X^+(F')$, $v\notin X(F')$, $v\perp H$ and $\beta_1^{-1}\beta_2v-\mu v\in X(F')$ with $\mu\in E$, then $\mu=1$
(\lref{PI.UnitaryGroupMassFormula.BetaLUniq.eigenvalue_last_line_eq_one}, which proves this for generic pairs of corank one).

(e) The pair conditions include that $\beta_1,\beta_2$ are isometries mapping the flag of $F'$ onto that of $F''$ and $z(F')$ to
$z(F'')$. The line reduction \lref{PI.UnitaryGroupMassFormula.BetaLUniq.line_reduction_holds} states that for such
isometries, (a) for $(F,F')$, agreement modulo $X(F'')$ on $H$, (c) and (d) imply $\beta_2x-\beta_1x\in X(F'')$ for all
$x\perp X(F')$. By (a)--(d) its hypotheses hold.
\end{proof}

\begin{remark}\label{pf:thmS:rem2}
(1) The density of $K$-finite vectors (\cref{pf:thmS:kfindense}) enters the formal proof as an instance argument of the lemmas
which use it; the instance is discharged by the averaging argument given above. Only the continuity and unitarity of $\sigma$
are used; the Harish-Chandra parameter plays no role.

(2) In \cref{pf:thmS:acle} the hypothesis $C_l\ge1000+|c|$ is used only through $C_l\ge|c|$ (which also gives $C_l\ge0$).

(3) The congruence of \cref{pf:thmS:hyptt} is deduced from the stronger congruence $T_u\equiv1\pmod{g_S(D)^2}$ for every word.

(4) The local and global inputs (\cref{pf:thmS:localcores,pf:thmS:globalcores}) are conjunctions of $32$ and $27$ statements of the
development; several of them are implications whose hypotheses are other members of the conjunction (for instance the joint
Poisson statement assumes the Poisson summation statement), and the conjunctions are proved by proving each member.
\end{remark}

% ---------------------------------------------------------------------------
\subsection{Remarks on the formal statements}\label{sec:pf:thmS:remarks}

\begin{remark}\label{pf:thmS:rem1}
(1) In \cref{pf:thmS:lg1} the slot $R$ is the slot of all cores, and the large-base condition holds trivially for it; the
three rows hold at every slot $Q$ and for every linkage (\cref{pf:thmS:everyslot}), because their contents
(\cref{pf:thmS:atchoice}) hold for all classes and all raw selections. The row $\mathsf{hTF}$ does not depend on the
selection at all.

(2) Several lemmas of \cref{sec:pf:thmS:fam,sec:pf:thmS:loops} take auxiliary propositions as instance arguments: the formal-degree identity (the formal degree at $\tau$ equals the Harish-Chandra formal degree of the corresponding parameter), the representation property of CM fields (for a CM field unramified over its real subfield and totally positive
real $a,b,n$ there are $s,t$ with $as\bar s+bt\bar t=n$), the integrability of a weighted matrix coefficient of discrete series with large gap, and two decay statements for such coefficients. All are discharged by theorems of
the development (\lref{Formalism.L2.SG.OwnFacts.FD.factDsigHC}, \lref{GlobalTheory.hermRepUnram}, and \lref{GlobalTheory.Sloc.CoefWeighted.coefWeightedLoc_uncond}
and the components of $\mathrm{GH}$); the statements proved are therefore unconditional apart from their explicit
hypotheses.

(3) In the pointwise theorem (\cref{pf:thmS:pointwise}) the Hilbert--Schmidt square $\mathrm{hs}(u,u')$ of the size data is
shown to equal $\|K_0(u,u')\|^2$ (\cref{pf:thmS:packagetie}); the two conclusions are the same bound. The family statement
(\cref{pf:thmS:main}) keeps only the first.

(4) The lattice-link hypothesis of the main row is not used (\cref{pf:thmS:latticelinkunused}); it is nevertheless
supplied by a theorem in the proof of \cref{pf:thmS:pbapplied}.

(5) The constants in the rows are fixed: $C_0=3$, count constants $(C_{\mathrm{YL}}+1,\frac14,3,\frac12,C_{\mathrm{Tu}},3)$, and
in the transport $C_\Delta=5$, $c_T=\frac12$, $C_b=3$. The non-generic row is used with $\varepsilon_{\mathrm{NG}}=
2^{-(k-2r_S-5)}\exp(C(2r_S+2)^5f\log b)$, and the loop rows with $\thmSepsL(D)=\exp(-\frac13kf\log k)$; the final constant
$1+b^{-6}$ comes from \cref{pf:thmS:num}(7), where $b\ge16$.

(6) The term systems of \cref{pf:thmS:terms} are abstract: \cref{pf:thmS:pba} uses only the listed axioms, and the
identification of the concrete term system $Y.t_O(u,u')$ with an orbital expansion enters only through the hypotheses
(H5), (H7)--(H10), (H18), (H19) of the pointwise theorem.
\end{remark}
